\section{B. Good reduction and inertia characters}
\subsection{Explicit counterexamples to the mod-$p$ good-reduction criterion}

\begin{theorem}[B-RD1: a counterexample for every odd prime]\label{thm:B-RD1}
\status{Proved.}
For every prime $p\geq3$, there is a complete discretely valued field $K$
of characteristic zero, residue field $\overline{\mathbf F}_p$, and
absolute ramification index
\[
 e(K)=\frac{p^2-1}{2},
\]
containing $\mu_p$, and an elliptic curve $E/K$ whose full $p$-torsion
is $K$-rational and whose reduction is additive, potentially good and
supersingular. Its semistability defect is exactly $2$. Its inertia
character is $-2$ on the nontrivial element of the finite inertia image.
\end{theorem}
\begin{proof}
We give the field and the curve by an explicit algebraic construction.
Put $k=\overline{\mathbf F}_p$, $R=W(k)$ and $F=\operatorname{Frac}R$.
Let $m=(p-1)/2$ and choose any root $\lambda\in k$ of
\[
 H_p(T)=\sum_{i=0}^{m}\binom mi^2T^i.
\]
Such a root exists since $m\geq1$. It is neither $0$ nor $1$:
$H_p(0)=1$ and
$H_p(1)=\binom{2m}{m}\equiv(-1)^m\pmod p$.
The coefficient calculation for Frobenius on the one-dimensional space
$H^1(\mathcal O)$ of a Legendre elliptic curve gives $(-1)^mH_p(\lambda)$
as its Hasse--Witt scalar. Thus
\[
 E_0/R:\quad y^2=x(x-1)(x-[\lambda])
\]
is smooth and has supersingular special fibre. Here $[\lambda]$ is the
Teichmuller representative. The discriminant is
$16[\lambda]^2(1-[\lambda])^2$, an $R$-unit.

Let $L=F(E_0[p])$. We now verify its ramification without invoking a
Breuil--Kisin dictionary. Every nonzero geometric $p$-torsion point
specializes to the identity, since the special fibre is supersingular.
For a local parameter $T$ of the formal group at that identity, its
height is two. Consequently $[p](T)/T$ has constant coefficient $p$,
its first coefficient which is a unit modulo $p$ is that of
$T^{p^2-1}$, and all lower nonconstant coefficients are divisible by $p$.
Weierstrass preparation gives a polynomial of degree $p^2-1$ which is
Eisenstein over $R$. Thus every nonzero $P\in E_0[p]$ generates a
completely ramified extension $F(P)/F$ of degree $p^2-1$. This degree is
prime to $p$, so the extension is tame. Since $F$ has algebraically
closed residue field, its unique tame extension of that degree is
Galois cyclic and contains every such point. Hence
\[
 \operatorname{Gal}(L/F)\simeq C_{p^2-1}.
\]
Its action on $E_0[p]\setminus\{0\}$ is free and transitive. The unique
involution in this cyclic group acts as $-1$: an involution in
$\mathrm{GL}_2(\mathbf F_p)$ with any eigenvalue $+1$ would fix a
nonzero vector.

Let $H$ be this order-two subgroup and put $K=L^H$. The Weil pairing
shows $\mu_p\subset L$, and $H$ acts on it by the determinant of $-1$,
which is $1$. Thus $\mu_p\subset K$. We have $e(K)=(p^2-1)/2$.
Choose any nonzero $P\in E_0[p]$, and put $\delta=y(P)^2\in K^\times$.
A nonzero odd-order point has $y(P)\ne0$, so $L=K(\sqrt\delta)$.
Write $E_0$ as $y^2=x^3+a_2x^2+a_4x+a_6$ and define
\[
 E:\quad y^2=x^3+\delta a_2x^2+\delta^2a_4x+\delta^3a_6.
\]
Over $L$, the map $(x,y)\mapsto(\delta x,\delta\sqrt\delta\,y)$
identifies $E_0$ with $E$. For every $Q\in E_0[p]$, both $x(Q)$ and
$\sqrt\delta\,y(Q)$ are fixed by $H$. Hence every point of $E[p]$
is $K$-rational. Equivalently, its representation is the tensor product
of the scalar quadratic representation on $E_0[p]$ and its inverse.

The curve $E_0$ has good reduction over $K$, so for every prime
$\ell\ne p$ inertia acts trivially on $T_\ell(E_0)$. On $T_\ell(E)$ it
acts by the nontrivial ramified quadratic character of $L/K$. The
Neron--Ogg--Shafarevich criterion\cite[Theorem~1]{SerreTate1968Good} therefore excludes good reduction.
It has good supersingular reduction over $L$, so its bad reduction is
additive and its minimal good-reduction inertia quotient has order $2$.
The trace on its covariant Dieudonne module is $-2$ at the involution,
since that involution acts as $[-1]$ on the special elliptic curve.
\end{proof}

\begin{theorem}[B-RD2: full $3$-torsion already over $\mathbf Q_3(\zeta_3)$]\label{thm:B-RD2}
\status{Proved, with an exact computation.}
Let $K=\mathbf Q_3(\zeta)$, where $\zeta^2+\zeta+1=0$. Then $K$ has
absolute ramification index $2$ and residue field $\mathbf F_3$. The curve
\[
 E/K:\quad y^2+y=x^3
\]
has $E[3]\subset E(K)$, additive potentially supersingular reduction,
and semistability defect exactly $6$. Its inertia character takes the
value $-2$ on the element of order $2$ of its finite inertia image.
\end{theorem}
\begin{proof}
The discriminant and $j$-invariant are $\Delta=-27$ and $j=0$.
The third division polynomial is
\[
 \psi_3(x)=3x(x^3+1).
\]
Besides the identity, the following are the eight distinct $3$-torsion
points:
\[
 (0,0),\ (0,-1),\qquad
 (x,y)\quad\text{with }x\in\{-1,-\zeta,-\zeta^2\},\quad
 y\in\{\zeta,\zeta^2\}.
\]
They exhaust $E[3]$, which has degree nine in characteristic zero.
The exact program \texttt{verify\_p3\_full\_torsion.py} checks all these
points, their orders, all $81$ sums, and that $(0,0),(-1,\zeta)$ form a
basis. Its recorded output is part of the computation certificate.

The displayed equation has $v_K(\Delta)=6$. Discriminants of two
Weierstrass equations of the same curve differ by a twelfth power.
Thus no equation over $K$ can have discriminant valuation zero, and
$E$ does not have good reduction.

We construct good reduction after ramification degree six. Make first
the unramified quadratic extension $K'=K(i)$, with $i^2=-1$, and choose
$\sqrt3\in K'$. Adjoin $r$ with $r^3=-1/4$, and then $u$ with
$u^2=r\sqrt3$. The polynomial for $s=r+1$ is
\[
 s^3-3s^2+3s-3/4.
\]
Its Newton polygon has the single slope $-2/3$ for the valuation
normalized on $K'$, so $K'(r)/K'$ is completely ramified of degree three.
The element $r\sqrt3$ has odd valuation in that field, so adjoining $u$
gives a completely ramified quadratic extension. Put $L=K'(r,u)$.
After completing the square and setting
\[
 x=r+u^2X,\qquad y+\tfrac12=u^3Y,
\]
the equation becomes
\[
 Y^2=X^3+\frac{3r}{u^2}X^2+X.
\]
The middle coefficient has positive valuation and the special fibre is
$Y^2=X^3+X$, which is smooth and supersingular in characteristic three.
Thus $E$ has good reduction over an extension of ramification degree six.

We verify minimality, including possible unramified extensions. Let
$F/K$ have good reduction. The discriminant valuation implies
$2\mid e(F/K)$. In a good integral Weierstrass equation over $F$, complete
the square; this is integral since $2$ is a unit. The resulting equation
has the form $Y^2=X^3+b_2X^2+b_4X+b_6$. Comparison with
$y^2=x^3+1/4$ uses $x=r_0+u_0^2X$, $y=u_0^3Y$, and forces
$v_3(u_0)=1/4$. Integrality of $b_6$ gives
\[
 v_3(r_0^3+1/4)\geq3/2.
\]
It follows that $r_0$ is a unit with residue $-1$. Writing $s_0=r_0+1$,
if $v_3(s_0)<1/3$ the term $s_0^3$ uniquely has lowest valuation in
$s_0^3-3s_0^2+3s_0-3/4$; if $v_3(s_0)>1/3$ the constant term uniquely
has lowest valuation. Both alternatives give valuation at most one.
Hence $v_3(s_0)=1/3$, so $3\mid e(F/\mathbf Q_3)$, and therefore
$3\mid e(F/K)$. We conclude $6\mid e(F/K)$, proving the defect is six.

Over the unramified extension $K'$, the extension just constructed is
Galois: $r\mapsto\zeta r$, $u\mapsto\zeta^2u$ has order three and
commutes with $u\mapsto-u$. Its inertia group is cyclic of order six.
The latter involution induces $Y\mapsto-Y$ on the good special fibre,
that is $[-1]$, so its trace is $-2$.
\end{proof}

\begin{proposition}[B-RD3: the ordinary case of R3(a)]\label{prop:B-RD3}
\status{Proved.}
Let $K$ be complete discretely valued of characteristic zero, with
perfect residue field of characteristic $p\geq3$, and assume
$\mu_p\subset K$. If an abelian variety $A/K$ has potentially good
ordinary reduction and $A[p]$ is unramified, then $A$ has good reduction.
There is no restriction on $\dim A$.
\end{proposition}
\begin{proof}
We may pass to the completed maximal unramified extension: good
reduction is equivalent before and after this passage by the
Neron--Ogg--Shafarevich criterion. The residue field is now algebraically
closed. Choose a finite Galois extension $L/K$ giving good reduction,
and let $\mathcal A/\mathcal O_L$ be the resulting abelian scheme.
For its ordinary special fibre, the connected subgroup
$\mathcal A[p]^0$ is of multiplicative type and its etale quotient has
rank $p^g$. Indeed the Cartier dual of the connected subgroup has
etale special fibre, hence is etale by Nakayama's lemma applied to its
relative differentials. Over the strictly henselian ring
$\mathcal O_L$, it is constant; thus
$\mathcal A[p]^0\simeq\mu_p^g$.

Every inertia element acts semilinearly on this subgroup by its action
on the base, composed with a matrix in $\mathrm{GL}_g(\mathbf F_p)$.
The base action fixes $\mu_p$, and its generic action on $A[p]$ is
trivial. The matrix is therefore the identity. The action on the etale
quotient is trivial for the same reason. Over the algebraically closed
residue field the connected--etale sequence splits, and there are no
nonzero homomorphisms between the connected and etale factors in either
direction. Thus inertia acts identically on $\mathcal A_k[p]$.

If $\alpha$ is one of the resulting finite-order automorphisms of
$\mathcal A_k$, then $\alpha-1$ factors through $[p]$, so
$\alpha=1+p\beta$ for some $\beta\in\operatorname{End}(\mathcal A_k)$.
The kernel of reduction modulo $p$ in the units of this endomorphism
ring has no nontrivial finite-order elements for $p\geq3$. To see this,
pass to its finite free $\mathbf Z_p$-completion. If
$\alpha=1+p^a\gamma$ with $a\geq1$ and $\gamma\notin p\operatorname{End}$,
the binomial expansion shows that a prime-to-$p$ power has the same
leading valuation, and the $p$-th power has leading valuation $a+1$;
neither can be one. This proves the assertion for any finite order.
Hence inertia is trivial on the entire special abelian variety.
Proper smooth base change now makes $T_\ell(A)$ unramified for
$\ell\ne p$, and the same good-reduction criterion concludes.
\end{proof}

\begin{corollary}[B-RD6: the manuscript's cas test is negative]\label{cor:B-RD6}
\status{Refuted, for every odd prime.}
For every $p\geq3$ there is a field $K$ as in Theorem~\ref{thm:B-RD1}
and two elliptic curves $E,E'/K$ with potentially good reduction and
isomorphic $p$-torsion, for which the inertia traces differ on an element
of order prime to $p$.
\end{corollary}
\begin{proof}
Use $E$ and $K$ from Theorem~\ref{thm:B-RD1}. Choose an ordinary
elliptic curve $C/\overline{\mathbf F}_p$; for example choose a Legendre
parameter outside the finite root set of $H_p$ and outside $\{0,1\}$.
The ordinary Serre--Tate coordinates of
\cite[Theorem~2.1]{Katz1981SerreTate} give its canonical deformation
over $W(\overline{\mathbf F}_p)$, whose $p$-divisible group is
\[
 \mu_{p^\infty}\oplus\mathbf Q_p/\mathbf Z_p.
\]
This formal elliptic scheme algebraizes: its identity section provides
the compatible relatively ample bundles $\mathcal O(3[0])$, so the
formal existence theorem applies. Let $E'$ be its generic fibre after
extension to $K$. Then $E'$ has good reduction, and $\mu_p\subset K$
gives $E'[p]\simeq(\mathbf Z/p)^2\simeq E[p]$ over $K$.
The good curve has inertia trace $2$; the bad quadratic twist has trace
$-2$ at the inertia involution. That involution has order prime to $p$.
\end{proof}

\subsection{The first torsion level: a sharp tame ramification range}

\begin{theorem}[B-RD7: the first-level tame range]\label{thm:B-RD7}
\status{Proved; independently audited by the parent agent.}
Let $p\geq5$, and permit arbitrary perfect residue fields of characteristic
$p$. There exists a complete discretely valued $K$ of characteristic zero
and absolute ramification index $e$, and an elliptic curve $E/K$ with
$E[p]\subset E(K)$, additive potentially good reduction, and semistability
defect $m$, if and only if
\[
 m\in\{2,3,4,6\},\qquad m\mid(p+1),\qquad
 (p-1)\mid e,\qquad e\geq\frac{p^2-1}{m}.
\]
In the existence direction one may take the residue field to be
$\overline{\mathbf F}_p$. In particular the set of possible $e$ is
\[
 \left\{e\in(p-1)\mathbf Z_{>0}:
 e\geq\frac{p^2-1}{m_{\max}(p)}\right\},
\]
where $m_{\max}(p)=2,6,4,6$ for $p\equiv1,5,7,11\pmod {12}$,
respectively. For $p=3$, the same assertion with a specified tame defect
$m=2$ or $4$ holds: defect $2$ occurs exactly at even $e\geq4$, and
defect $4$ occurs at every positive even $e$. No assertion about wild
3-primary inertia is included in the latter statement.
\end{theorem}

\begin{proof}
We first prove necessity. Completion of the maximal unramified extension
does not change $e$ or the inertia image. The ordinary case is excluded
by Proposition~\ref{prop:B-RD3}. For $p\geq5$, potential-good inertia is
tame by \cite[Corollary~2 of Theorem~2]{SerreTate1968Good}. Thus the
minimal good-reduction extension $L/K$ is cyclic of order $m$, with
supersingular reduction $C$. Over an algebraically closed field of
characteristic greater than three, writing $C$ in short Weierstrass form
shows that its origin-preserving automorphisms have the form
$(x,y)\mapsto(u^2x,u^3y)$, with $u^4a=a$ and $u^6b=b$.
The possibilities are $m=2,3,4,6$. The curves with $j=0$ are supersingular
exactly when $p\equiv-1\pmod3$, and those with $j=1728$ exactly when
$p\equiv-1\pmod4$. Both assertions follow directly by extracting the
coefficient of $x^{p-1}$ from $(x^3+b)^{(p-1)/2}$ and
$(x^3+ax)^{(p-1)/2}$, respectively. Consequently $m\mid p+1$.
The same discussion for tame inertia in characteristic three gives
$m=2$ or $4$.

Normalize the valuation $v_L$ by $v_L(L^\times)=\mathbf Z$.
An average over the tame cyclic action provides a formal parameter $t$
at the identity of the good model which is an exact eigenvector.
If its eigenvalue is a primitive $m$th root $\zeta^q$, then every
$K$-rational nonzero $p$-torsion point $P$ satisfies
\[
 v_L(t(P))\equiv q\pmod m,\qquad \gcd(q,m)=1.
\]
All these valuations are positive, since $C$ is supersingular. Choose a
point $P$ whose parameter has maximum valuation $A$, and a point $Q$
whose parameter has minimum valuation $B$. If $A>B$ they are independent;
if $A=B$, choose any basis. The $p-1$ nonzero multiples of $P$ have
valuation $A$. The other $p(p-1)$ points have valuation $B$: if $A>B$,
this follows from the linear term of the formal group law; if $A=B$,
it follows from the assumed common maximum and minimum.
The constant term of $[p](t)/t$ is $p$, and its distinguished degree is
$p^2-1$. Weierstrass preparation therefore gives
\[
 m e=v_L(p)=(p-1)(A+pB).                         \tag{B.7.1}
\]
Here $A\geq B\geq1$ and $A\equiv B\equiv q\pmod m$.
Since $p\equiv-1\pmod m$, the displayed identity implies both
$(p-1)\mid e$ and $me\geq p^2-1$.

We give a deformation-theoretic converse with the equivariance made
explicit. Put $k=\overline{\mathbf F}_p$. Choose a supersingular elliptic
curve $C/k$ with an automorphism $\alpha$ of order $m$. For $m=3,6$
choose $j=0$, for $m=4$ choose $j=1728$ (the unique supersingular
$j$ in characteristic three is also zero), and for $m=2$ choose any
supersingular curve. The derivative of $\alpha$ is a primitive $m$th
root; a finite automorphism of order prime to $p$ with derivative one
is the identity, by averaging a local coordinate.

The framed deformation functor of $C$, with full Drinfeld $p$-level
structure, is represented by a complete regular local ring $R$ of
dimension two. If $P,Q$ are the universal labelled basis, their formal
coordinates form a regular system of parameters. These are the
precise statements of \cite[Theorem~2.1.2 and
Proposition~2.2.3(i)]{CarterStrauch2021LubinTate}; Serre--Tate identifies
the height-two formal-group deformation with the elliptic deformation
\cite[Theorem~1.2.1]{Katz1981SerreTate}. The universal formal elliptic
curve algebraizes: the identity gives the compatible ample line
bundles $\mathcal O(3[0])$, so the projective formal existence theorem
applies over the complete noetherian ring $R$.

Changing the framing by $\alpha$ gives an action of $C_m$ on this
functor, hence a semilinear action on the universal elliptic curve.
It preserves the labels $P,Q$. For any initial formal coordinate $t_0$,
replace it by
\[
 t=\frac1m\sum_{i=0}^{m-1}\zeta^{-i}\alpha^{i*}t_0,
\]
where $\zeta$ is the derivative eigenvalue. The linear coefficient is
a unit, so $t$ remains a coordinate and $\alpha^*t=\zeta t$.
Thus $x=t(P)$ and $y=t(Q)$ are regular parameters satisfying exactly
$\alpha(x)=\zeta x$, $\alpha(y)=\zeta y$. The action fixes $W(k)$.
Cohen's presentation theorem gives an equivariant surjection
\[
 W(k)[[X,Y]]\longrightarrow R,\qquad X\mapsto x,\quad Y\mapsto y,
\]
where $\alpha$ multiplies $X,Y$ by $\zeta$.

We detail its single relation, to avoid assuming that an arbitrary
monomial specialization has the intended ramification. The full-level
factorization gives, in $R$,
\[
 p=u\prod_{(a,b)\in\mathbf F_p^2\setminus\{0\}}
          t([a]P+[b]Q),\qquad u\in R^\times.
\]
Lift $u$ to a unit $U$ and each factor to a power series $T_{a,b}$
in $W(k)[[X,Y]]$ by using the multiplication and addition power
series of the formal group, lifting their coefficients. These lifts
have zero constant term and linear term $aX+bY$ (using integral
representatives for $a,b$); crucially, $T_{a,0}$ is divisible by $X$
with unit quotient when $a\ne0$. Then
\[
 F_0=p-U\prod_{(a,b)\ne0}T_{a,b}
\]
lies in the kernel and generates it: it has coefficient one on $p$
modulo the square of the maximal ideal of $W(k)[[X,Y]]$.
Replace $F_0$ by its $C_m$-average $F$. This still generates the kernel,
has the same linear term $p$, and is invariant.

Write $t=e/(p-1)$ and choose
\[
 B=1,\qquad A=1+m\left(t-\frac{p+1}{m}\right).
\]
These are positive integers, $A\geq B$, and $A\equiv B\equiv1\pmod m$.
Reduce the relation modulo $p$, and give $X,Y$ weights $A,B$. If $A>B$, the weighted initial term of the
product in $F_0$ is a nonzero scalar times
$X^{p-1}Y^{p(p-1)}$, of weight
$d=(p-1)(A+pB)$. If $A=B$, the weighted initial term is the nonzero
homogeneous polynomial $\prod_{(a,b)\ne0}(aX+bY)$ of degree $p^2-1$.
In either case that initial form is invariant, since its total degree
is $p^2-1$, divisible by $m$. Averaging therefore preserves it.
Choose a Teichmuller unit $c$ arbitrarily when $A>B$, and with
$\bar c\notin\mathbf F_p$ when $A=B$. Then
\[
 R_L=W(k)[[\pi]]/(F(\pi^A,c\pi^B))
\]
is a complete discrete valuation ring with uniformizer $\pi$ and
$v_L(p)=d$: modulo $p$ its defining equation has first nonzero term
of degree $d$, while its coefficient on $p$ is a unit. This also
proves directly that the specified specialization exists.

The action $\pi\mapsto\zeta\pi$ preserves this equation, by invariance
of $F$ and $A\equiv B\equiv1\pmod m$. It has order $m$, acts trivially
on the residue field, and makes the map $R\to R_L$ equivariant.
Let $L=\operatorname{Frac}R_L$ and $K=L^{C_m}$. Then $L/K$ is tame
of ramification degree $m$, and
\[
 e(K)=d/m=e.
\]
The generic elliptic curve descends to $K$. Since the labels of the
level structure were preserved, its full $p$-torsion is rational over
$K$. The inertia action on the good special fibre is the faithful
order-$m$ action $\alpha$. The Neron--Ogg--Shafarevich criterion gives
exactly semistability defect $m$, so reduction is additive. Finally the
Weil pairing implies $\mu_p\subset K$. This proves every sufficiency
assertion made in the theorem, including $p=3$ and $m=2,4$.
\end{proof}

\subsection{A missing wild semistability defect in characteristic three}

\begin{theorem}[B-RD8: defect twelve with full rational $3$-torsion]\label{thm:B-RD8}
\status{Proved; exact computation recorded.}
The list $2,3,4,6$ in B.2 is incomplete in residue characteristic three.
Put
\[
 K=\mathbf Q_3(i,\zeta),\qquad i^2=-1,\quad
 \zeta^2+\zeta+1=0,\qquad \pi=\zeta-1,\quad t=1+i\pi.
\]
Then $K$ has absolute ramification index $2$ and residue field
$\mathbf F_9$. The elliptic curve
\[
 E:\quad y^2=x^3-27t(t^3+8)x+54(t^6-20t^3-8)
\]
has full $3$-torsion over $K$ and additive potentially supersingular
reduction, with semistability defect exactly $12$.
\end{theorem}
\begin{proof}
Write $a=-27t(t^3+8)$ and $b=54(t^6-20t^3-8)$.
Set $\delta=1+2\zeta=3+2\pi$, so $\delta^2=-3$.
The eight nonzero $3$-torsion points are
\[
 \left(-9t^2,\ \pm12\delta(t^3-1)\right)
\]
and, for $\epsilon\in\{1,\zeta,\zeta^2\}$,
\[
 \left(3t^2+12\epsilon t+12\epsilon^2,
       \ \pm36(\epsilon t^2+\epsilon^2t+1)\right).
\]
Indeed their four $x$-coordinates give the factorization
\[
 3x^4+6ax^2+12bx-a^2
 =3(x+9t^2)\prod_{\epsilon^3=1}
       (x-3t^2-12\epsilon t-12\epsilon^2).
\]
Substitution verifies their $y$-coordinates; the points are distinct
and exhaust the eight nonzero geometric $3$-torsion points.
The exact certificate \texttt{verify\_p3\_defect12.py} performs these
calculations in
$\mathbf Q[i,\pi]/(i^2+1,\pi^2+3\pi+3)$, verifies the curve equation
and the tripling identity for every point, and verifies independence
of points with different $x$-coordinates. No numerical approximation
is used.

Normalize $v_K(\pi)=1$. Direct calculation gives
\[
 \Delta=2^{12}3^9(t^3-1)^3,\quad
 v_K(t^3-1)=v_K(t^3+8)=3,\quad v_K(\Delta)=27.
\]
In particular $E$ is not good over $K$, since the valuation of a
Weierstrass discriminant can change only by a multiple of $12$.
We next both construct good reduction and prove the required minimality.

In the $2$-division polynomial $x^3+ax+b$, substitute
$x=3\pi^2 Z$. Dividing by $(3\pi^2)^3$ gives
\[
 Z^3+A Z+B,
 \qquad A=4i\pi+3\pi+6,\quad B=-12i+20\pi+30.
\]
Both $A$ and $B$ have valuation one. The polynomial is Eisenstein,
so a root $z$ generates a totally ramified cubic extension $K_1/K$,
and $v_K(z)=1/3$. Set $x_0=3\pi^2 z$. Then
\[
 v_K(x_0)=13/3,\qquad v_K(3x_0^2+a)=9.
\]
The second equality can also be read from the derivative
$3z^2+A$, whose valuation is one. Adjoin $u$ satisfying
\[
 u^4=3x_0^2+a.
\]
In the normalized valuation of $K_1$, the right side has value $27$,
prime to four. Thus this is a totally ramified degree-four extension;
for $L=K_1(u)$ we have $e(L/K)=12$ and $v_K(u)=9/4$.
The substitution $x=x_0+u^2X$, $y=u^3Y$ gives
\[
 Y^2=X^3+\frac{3x_0}{u^2}X^2+X.
\]
The middle coefficient has valuation $11/6>0$. The special fibre is
$Y^2=X^3+X$, smooth and supersingular in characteristic three.
Therefore good reduction is obtained with ramification degree twelve.

Conversely, let $F/K$ be any finite extension over which $E$ has good
reduction. Its discriminant valuation implies
$12\mid27e(F/K)$, hence $4\mid e(F/K)$. Since $2$ is invertible,
the $2$-torsion of a good elliptic scheme is finite etale; all its
points are defined over the maximal unramified extension of $F$.
The Eisenstein $2$-division root has $K$-valuation $1/3$ after the
indicated scaling, so $3\mid e(F/K)$. Consequently
$12\mid e(F/K)$. Together with the constructed extension, this proves
that the semistability defect is exactly twelve.
\end{proof}

\subsection{The complete first-level range in residue characteristic three}

\begin{lemma}[B-RD9: a variable-ramification Hesse family]\label{lem:B-RD9}
\status{Proved.}
Let $K$ have algebraically closed residue field of characteristic three,
absolute ramification index $e$, and a uniformizer $\varpi$. Suppose
$\mu_3\subset K$. Choose an integer $r$ with
$0<2r<e$ and $3\nmid e+r$, and put $t=1+\varpi^r$. The Hesse
Weierstrass curve of Theorem~\ref{thm:B-RD8} has full rational
$3$-torsion and semistability defect
\[
 m=\frac{12}{\gcd(4,e+r)}.
\]
\end{lemma}
\begin{proof}
The torsion formulas in Theorem~\ref{thm:B-RD8} are identities in $t$,
so apply unchanged. Write $w=\varpi^r$ and first scale the $2$-division
polynomial by $x=3X$. After translating $X=w^2+Z$, the resulting cubic is
\[
 Z^3+3w^2Z^2-(12w^3+18w^2+36w+27)Z
       +12w^4-36w^3-117w^2-108w-54.
\]
Its coefficients, from constant through quadratic, have valuations
$e+4r,e+3r,e+2r$. The inequalities $2r<e$ make the displayed leading
terms unique. Its Newton polygon has the single slope
$-(e+4r)/3$. Because $3\nmid e+r$, adjoining a root $z$ gives a
completely ramified cubic extension $K_1/K$. Put $x_0=3(w^2+z)$.
Then
\[
 v_K(x_0)=e+2r,\qquad v_K(3x_0^2+a)=3(e+r).
\]
The derivative computation giving the second equality has a unique
lowest term, of valuation $e+3r$, before multiplying by $3^2$.

Choose $u$ with $u^4=3x_0^2+a$. Since the residue field is algebraically
closed and $4$ is invertible, adjoining this root has ramification
degree $4/\gcd(4,e+r)$ over $K_1$. The equation obtained from
$x=x_0+u^2X$, $y=u^3Y$ is
\[
 Y^2=X^3+(3x_0/u^2)X^2+X,
\]
and the middle coefficient has valuation $(e+r)/2>0$. Thus it is good.
Conversely $v_K(\Delta)=9(e+r)$, which forces the tame factor
$4/\gcd(4,e+r)$ in any good-reduction ramification degree. The
$2$-division cubic forces the additional factor three, exactly as in
Theorem~\ref{thm:B-RD8}. This proves minimality.
\end{proof}

\begin{theorem}[B-RD10: all first-level defects at $p=3$]\label{thm:B-RD10}
\status{Proved; independently audited by the parent agent.}
Allowing arbitrary perfect residue fields, a potentially good additive
elliptic curve with full rational $3$-torsion exists at every positive
even absolute ramification index. The precise possible defects are
\[
\begin{array}{c|c}
\text{defect}&\text{absolute ramification indices}\\\hline
2& e\in2\mathbf Z,\ e\geq4\cr
3& e\in2\mathbf Z,\ e\geq4\cr
4& e\in2\mathbf Z,\ e\geq2\cr
6& e\in2\mathbf Z,\ e\geq2\cr
12& e\in2\mathbf Z,\ e\geq2.
\end{array}
\]
Over the particular field $\mathbf Q_3(\zeta_3)$, Theorem~\ref{thm:B-RD2}
already gives an example of defect six.
\end{theorem}
\begin{proof}
The special fibre is supersingular by Proposition~\ref{prop:B-RD3}.
Over an algebraically closed field of characteristic three it is
isomorphic to $y^2=x^3-x$. Its automorphisms are
\[
 (x,y)\longmapsto(u^2x+r,u^3y),\qquad
 u^4=1,\quad r^3-r=0.
\]
Thus its automorphism group has order twelve, is $C_3\rtimes C_4$,
and the only possible nontrivial subgroup orders are $2,3,4,6,12$.
The inertia image embeds in this group.

Over a good-reduction extension of ramification degree $m$, all eight
nonzero $3$-torsion parameters have positive integral valuations.
Their sum is $v_L(3)=me$, by the constant term of $[3](T)/T$.
Consequently $me\geq8$, which rules out $m=2,3$ at $e=2$.
The Weil pairing gives $\mu_3\subset K$, hence $e$ is even.
These are all the necessary restrictions in the table.

The tame cases $m=2,4$ are Theorem~\ref{thm:B-RD7}.
We explain the remaining constructions explicitly. For any desired
even $e$, a convenient base with algebraically closed residue field is
\[
 K_e=\operatorname{Frac}\left(
 W(\overline{\mathbf F}_3)[\varpi]/
 (\varpi^e+3\varpi^{e/2}+3)\right).
\]
The polynomial is Eisenstein, and $1+\varpi^{e/2}$ is a primitive
cube root of unity. In all the cases below the variable-ramification
family is taken over $K_e$.

For defect twelve and $e\geq4$, use $r=1$ unless $3\mid e+1$, in
which case use $r=3$. The second case first occurs at $e=8$, so always
$0<2r<e$. The integer $e+r$ is odd and prime to three. The missing
case $e=2$ is Theorem~\ref{thm:B-RD8}.

For defect six with $e\equiv0\pmod4$, use $r=2$, replacing it by
$r=6$ when $3\mid e+2$. These choices work for every such $e\geq8$.
For $e\equiv2\pmod4$, use $r=4$, replacing it by $r=8$ when
$3\mid e+4$. These choices work except at $e=2,6,14$.
In each applicable case, $e+r\equiv2\pmod4$ and $3\nmid e+r$,
so Lemma~\ref{lem:B-RD9} applies. The case $e=2$ is
Theorem~\ref{thm:B-RD2}. The case $e=4$ is obtained from the
order-twelve example by the tame quadratic extension: its inertia
image then has order six. The case $e=14$ follows from the order-six
example at $e=2$ by a tame degree-seven extension.
For $e=6$, start with the order-six example over a base of algebraically
closed residue field, and adjoin a cube root of a uniformizer.
Its good-reduction cubic subfield is obtained by adjoining a cube root
of $-1/4$, whose Kummer class has valuation zero; the uniformizer class
has valuation one. These cyclic cubic extensions are distinct and
linearly disjoint. Thus the order-six inertia image remains unchanged.
All finite extensions here are totally ramified.

For defect three with $e\equiv0\pmod4$, use $r=4$, replacing it by
$r=8$ when $3\mid e+4$. These choices work for every such $e\geq12$.
For $e\equiv2\pmod4$ and $e\geq6$, use $r=2$, replacing it by
$r=6$ when $3\mid e+2$; only $e=10$ is exceptional.
The examples at $e=4,8$ come from the order-six example at $e=2$
by tame extensions of degree two and four: they remove its order-two
inertia and retain the order-three subgroup.
For $e=10$, work over $K_{10}$ and put $t=1+\varpi^2$.
In the $2$-division cubic scaled by $x=3X$, translate
$X=\varpi^4+\varpi^6+Z$. Its constant, linear and quadratic
coefficients have valuations $22,16,14$, respectively. Its Newton
polygon therefore has the single slope $-22/3$, forcing a ramified
cubic extension. Here $v_K(\Delta)=108$, and the good-model
construction has $v_K(u^4)=36$; its fourth root requires no additional
ramification over the cubic extension. Hence the defect is exactly
three. The exact certificate \texttt{verify\_p3\_ranges.py} verifies these
coefficient valuations by reducing polynomials modulo
$\varpi^{10}+3\varpi^5+3$, and separately audits all the index recipes
for every even $e\leq10000$.

For the assertions about tame base change, one can work throughout
over completed maximal unramified bases. The order-six example has
cyclic inertia $C_6$. The order-twelve example has the full group
$C_3\rtimes C_4$, whose tame quotient is $C_4$; a tame quadratic
extension leaves the order-six subgroup. These observations also
verify that the stated defects are not merely upper bounds.
\end{proof}

\begin{corollary}[B-RD11: the cyclotomic field]\label{cor:B-RD11}
\status{Proved.}
For $p\geq5$, no elliptic curve over $\mathbf Q_p(\zeta_p)$ can have
full rational $p$-torsion and additive potentially good reduction.
For $p=3$, such a curve exists.
\end{corollary}
\begin{proof}
For $p\geq7$, the necessary inequality in
Theorem~\ref{thm:B-RD7} gives
$e\geq(p^2-1)/6>p-1$, contradicting $e(\mathbf Q_p(\zeta_p))=p-1$.
For $p=5$, the only possible defect is six. We justify that the good
extension may be chosen totally ramified, so its residue field really
is $\mathbf F_5$. Choose a minimal integral short Weierstrass equation
$y^2=x^3+ax+b$, and put $\delta=v_K(\Delta)$. Since the reduction is
potentially good, $v_K(j)\geq0$, giving $3v_K(a)\geq\delta$.
The discriminant identity then gives $2v_K(b)\geq\delta$ as well.
Put $m'=12/\gcd(12,\delta)$. Over the totally ramified extension
$K(\varpi^{1/m'})$, scaling by $u=\varpi^{\delta/12}$ gives integral
coefficients and unit discriminant. Conversely the discriminant
forces $m'$ to divide every good-reduction ramification degree.
Thus $m'=6$, and this construction gives the required degree-six
extension with unchanged residue field. Equality now holds in the
valuation bound: over this good-reduction extension every nonzero
$5$-torsion formal parameter has valuation one. Dividing those
parameters by a uniformizer and reducing gives an injective
$\mathbf F_5$-linear map
\[
 E[5](K)\longrightarrow k_L=k_K=\mathbf F_5.
\]
Additivity follows from the linear term of the formal group law;
injectivity follows because every nonzero parameter has valuation
one. Its source would have dimension two and its target dimension
one, a contradiction. The claim for $p=3$ is
Theorem~\ref{thm:B-RD2}.
\end{proof}

\subsection{All torsion levels for tame potentially good elliptic curves}

\begin{theorem}[B-RD12: the sharp tame higher-level criterion]\label{thm:B-RD12}
\status{Proved; independently audited by the parent agent.}
Fix $p\geq5$, $n\geq1$, and $m\in\{2,3,4,6\}$ with $m\mid p+1$.
Assume $e$ is a positive multiple of $p-1$, and put
$s=v_p(e/(p-1))$. There is a complete discretely valued $K$ of absolute
ramification index $e$ and an elliptic curve with full rational
$p^n$-torsion and semistability defect $m$ if and only if
\[
 s\geq n-1,\qquad
 e\geq\frac{p^2-1}{m}\,
        p^{\max(2n-2,\,4n-4-s)}.                  \tag{B.12.1}
\]
Arbitrary perfect residue fields are allowed; the constructions use
$\overline{\mathbf F}_p$. The same assertion holds for $p=3$ with a
specified tame defect $m=2$ or $4$. In that assertion no restriction on
wild defects is claimed here; their sharp higher-level criterion is proved separately in Theorem~\ref{thm:B-RD26-wild}.
\end{theorem}
\begin{proof}
We first obtain an exact integer criterion. Pass to completed maximal
unramified bases and a tame cyclic degree-$m$ good-reduction extension
$L/K$. Use the eigenparameter from Theorem~\ref{thm:B-RD7}.
Let $A$ and $B$ be the maximum and minimum valuations of a primitive
$p^n$-torsion point. Points $P,Q$ achieving them form a basis if $A>B$:
otherwise $Q=aP+pR$ with $a$ a unit; multiplication by $p$ strictly
increases every positive formal valuation, contradicting minimality
of $B$. If $A=B$, choose any basis. We have
\[
 A\geq B\geq1,\qquad A\equiv B\equiv q\pmod m,
 \qquad \gcd(q,m)=1.                              \tag{B.12.2}
\]

The full-level deformation ring has the factorization of
\cite[Proposition~2.2.3(vii)]{CarterStrauch2021LubinTate}.
Iterated $j$ times, it expresses the coordinate of $[p^j]Q$ as a unit
times a product of $p^j$ primitive level-$n$ coordinates, each raised
to $p^j$. The relevant lines have representatives $(a,1)$ with
$a\equiv0\pmod {p^{n-j}}$. All these coordinates have valuation $B$,
because the contribution $[a]P$ has valuation strictly greater than
$B$. Consequently
\[
 v_L(t([p^j]Q))=p^{2j}B\quad(0\leq j\leq n-1).   \tag{B.12.3}
\]
This factorization applies to the actual good elliptic curve by its
map from the universal framed deformation ring with full level.

Primitive vectors modulo scalar units are the lines of
$\mathbf P^1(\mathbf Z/p^n)$. There are $p^n$ representatives $(a,1)$,
all of valuation $B$: when $A=B$, a cancellation would contradict the
chosen maximum $A$. There is the line $(1,0)$, of valuation $A$.
For $1\leq j<n$, there are $(p-1)p^{n-j-1}$ representatives $(1,b)$
with $v_p(b)=j$, each of valuation $\min(A,p^{2j}B)$. Again equality
cannot cause cancellation beyond $A$, since $A$ is the maximum among
primitive points. The series $[p^n](T)/[p^{n-1}](T)$ has constant term
$p$, distinguished degree $(p^2-1)p^{2n-2}$, and precisely the
primitive $p^n$-torsion roots. Taking its constant term gives
\[
 me=(p-1)p^{n-1}\left(
 A+p^nB+\sum_{j=1}^{n-1}(p-1)p^{n-j-1}
                \min(A,p^{2j}B)\right).          \tag{B.12.4}
\]

We prove that every pair of integers satisfying (B.12.2) and (B.12.4)
can be realized; for the stated range it will suffice to use $B=q=1$.
Use the full Drinfeld $p^n$-level deformation ring $R_n$, whose regular
parameters $x,y$ are the universal basis coordinates
\cite[Theorem~2.1.2 and Proposition~2.2.3(i)]{CarterStrauch2021LubinTate}.
The eigenparameter averaging and equivariant Cohen presentation of
Theorem~\ref{thm:B-RD7} apply verbatim to this level.
We justify the weighted specialization, including its possible
cancellations, as follows.

Work modulo $p$ in the power-series presentation and give $X,Y$
weights $A,B$. Lift the full-level factorization (B.12.3), rather than
an arbitrary expression for $[p^j]y$. Each factor with representative
$(a,1)$ and $a\equiv0\pmod {p^{n-j}}$ has weighted initial form $Y$,
up to a nonzero constant. Therefore a lift of $[p^j]y$ has weighted
initial form $c_jY^{p^{2j}}$, with $c_j\in k^\times$. It follows that
for a primitive line $(1,b)$, $v_p(b)=j$, one can choose a lift whose
initial form is the lower-weight term of
\[
 X+c_bY^{p^{2j}},\qquad c_b\ne0,
\]
or that entire binomial when the two weights agree. The other
primitive lines are handled by their linear terms. These choices are
legitimate lifts because the factorization identities hold in $R_n$;
changes by the defining relation do not alter their images.

The constant-term identity for $[p^n](T)/[p^{n-1}](T)$ gives a generator
of the Cohen kernel of the form $p$ minus a unit times the product of
all primitive coordinates. Average this generator under $C_m$.
Its weighted initial form modulo $p$ is preserved: every initial
factor has scalar character $\zeta$, since
$p^{2j}\equiv1\pmod m$, and the total number
$(p^2-1)p^{2n-2}$ of factors is divisible by $m$.
Now specialize $X=\pi^A$, $Y=c\pi^B$, choosing the residue of $c$
to avoid the finitely many roots of the displayed initial forms.
The resulting relation has linear term $p$ and first nonzero term
modulo $p$ of degree $d$, where $d$ is the right side of (B.12.4).
It defines a complete DVR with $v_L(p)=d$.
Every primitive-coordinate weight is strictly smaller than $d$, so
terms in its chosen lift divisible by $p$ cannot change that weight.
Thus all the asserted valuations hold after specialization, with no
circular assumption on $v_L(p)$.

For $A\equiv B\equiv1\pmod m$, the action $\pi\mapsto\zeta\pi$
preserves the relation and the specialization. Tame descent gives a
curve over the invariant field, of absolute ramification $d/m$,
with full rational $p^n$-torsion and exact defect $m$, as in
Theorem~\ref{thm:B-RD7}. For a general common residue $q$ in
(B.12.2), replace the action on $\pi$ by multiplication by
$\zeta^{q^{-1}}$, where the inverse is taken modulo $m$.

It remains to simplify (B.12.4), which also proves that $B=1$ suffices.
Choose $r\in\{0,\ldots,n-1\}$ so that
$p^{2r}B\leq A\leq p^{2r+2}B$, with the upper bound omitted for
$r=n-1$. Summing the geometric series in (B.12.4) gives
\[
 me=(p-1)p^{2n-r-2}(A+p^{2r+1}B).                \tag{B.12.5}
\]
It implies the two conditions
\[
 p^{2n-r-2}(p-1)\mid e,\qquad
 e\geq \frac{p^2-1}{m}p^{2n+r-2}.                \tag{B.12.6}
\]
Conversely suppose these conditions hold for some $r$. Increase $r$
until it is $n-1$, or until the second inequality fails for $r+1$;
the divisibility remains true. Put
\[
 B=1,\qquad
 A=\frac{me}{(p-1)p^{2n-r-2}}-p^{2r+1}.
\]
Then $A$ is an integer, $A\equiv1\pmod m$, and
$p^{2r}\leq A<p^{2r+2}$ unless $r=n-1$, when only the lower bound
is needed. Hence it realizes (B.12.4) by the construction above.
Finally the least permissible $r$ is
$\max(0,2n-2-s)$; such an $r\leq n-1$ exists exactly when
$s\geq n-1$. Substitution gives (B.12.1).
\end{proof}

\begin{corollary}[B-RD13: the least good-forcing level in the tame case]
\label{cor:B-RD13}
\status{Proved; independently audited by the parent agent.}
Let $p\geq5$, $e\in(p-1)\mathbf Z_{>0}$, and let $m_{\max}(p)$ be
as in Theorem~\ref{thm:B-RD7}. If $m_{\max}(p)e<p^2-1$, put $N=0$.
Otherwise put
\[
 s=v_p(e/(p-1)),\qquad
 L=\left\lfloor\log_p\frac{m_{\max}(p)e}{p^2-1}\right\rfloor,
 \qquad
 N=1+\min\left(s,\left\lfloor\frac L2\right\rfloor,
                     \left\lfloor\frac{L+s}{4}\right\rfloor\right).
\]
Among potentially good elliptic curves over all such fields of
absolute ramification $e$, the least positive level $n$ at which
full rational $p^n$-torsion forces good reduction is $N+1$.
For $p=3$, the same formula with $m_{\max}=4$ gives the least level
within the class of potentially good curves with tame inertia image.
\end{corollary}
\begin{proof}
For a fixed level, the largest allowed defect gives the weakest
inequality in Theorem~\ref{thm:B-RD12}. Its inequalities are equivalent
to $n-1\leq s$, $2n-2\leq L$, and $4n-4-s\leq L$.
Solving them gives exactly the displayed largest bad level $N$.
The construction realizes it, and its lower levels, while every
potentially good curve outside the bad cases is good.
\end{proof}

\begin{remark}[The scope of ``forces good reduction'']\label{auto:B:b9f011fa8019}
\status{Proved under stated hypotheses}

The potentially good hypothesis is essential. If $K$ contains
$\mu_{p^n}$, a split Tate curve with parameter $q=a^{p^n}$,
$0<|a|<1$, has full rational $p^n$-torsion and split multiplicative
reduction. Thus the preceding least-level assertion is not a
criterion for all elliptic curves without that hypothesis.
\end{remark}

\subsection{Cyclotomic bounds in every dimension and a character lower bound}

\begin{lemma}[B-RD14: full torsion forces the cyclotomic extension]
\label{lem:B-RD14}
\status{Proved.}
Let $A/K$ be an abelian variety of positive dimension over a
characteristic-zero field. If $A[p^n]$ is rational over $K$, then
$\mu_{p^n}\subset K$. A principal polarization, or a polarization of
degree prime to $p$, is not needed for this implication.
\end{lemma}
\begin{proof}
Choose any polarization defined over $K$. Its Weil pairing gives a
nonzero $G_K$-equivariant integral alternating form
\[
 b:T_pA\times T_pA\longrightarrow\mathbf Z_p(1).
\]
Divide all its coefficients by their largest common power of $p$.
The resulting form is still integral and equivariant, and at least
one coefficient is a unit. Since $G_K$ acts trivially on $T_pA/p^n$,
equivariance yields
$(\chi_p(\sigma)-1)b(x,y)\equiv0\pmod {p^n}$ for every $x,y$.
Choose a pair with unit value; then
$\chi_p(\sigma)\equiv1\pmod {p^n}$ for every $\sigma\in G_K$.
This is precisely the claimed cyclotomic containment.
\end{proof}

\begin{proposition}[B-RD15: the unconditional reading of the least-level question]
\label{prop:B-RD15}
\status{Proved; the final implication is explicitly vacuous.}
Fix $p\geq3$, a positive multiple $e$ of $p-1$, and $g\geq1$, and put
$s=v_p(e/(p-1))$. Among \emph{all} abelian varieties of dimension $g$
over all complete discretely valued fields with perfect residue field
and absolute ramification $e$, the least positive $n$ for which full
rational $p^n$-torsion implies good reduction is $s+2$.
At this level full rational $p^n$-torsion cannot occur at all.
Consequently a substantive good-reduction threshold must specify the
potentially good class, as in Corollary~\ref{cor:B-RD13}.
\end{proposition}
\begin{proof}
Lemma~\ref{lem:B-RD14} and the ramification degree
$[\mathbf Q_p(\mu_{p^n}):\mathbf Q_p]=p^{n-1}(p-1)$ imply
$p^{n-1}(p-1)\mid e$. Thus full rational $p^n$-torsion is impossible
for $n\geq s+2$.

Conversely, for each $1\leq n\leq s+1$, choose a complete field $K$
with residue field $\overline{\mathbf F}_p$, absolute ramification $e$,
and $\mu_{p^n}\subset K$. Such a field is obtained by taking the
cyclotomic extension of $\operatorname{Frac}W(\overline{\mathbf F}_p)$
and a further totally ramified extension of the required degree.
By \cite[Theorem~1, p.6 of the author version]{Tate1995Uniformization},
a split Tate curve with parameter $q=a^{p^n}$, where $a$ is a
uniformizer, has all its $p^n$-torsion represented by
$\zeta_{p^n}^i a^j$ in $K^\times/q^{\mathbf Z}$, and has split
multiplicative reduction. For dimension $g$, take its product with
$g-1$ canonical ordinary elliptic lifts. Their $p$-divisible groups
are $\mu_{p^\infty}\oplus\mathbf Q_p/\mathbf Z_p$, so they too have
full rational $p^n$-torsion. The product still has bad reduction.
This proves the asserted least level and explains its vacuity.
\end{proof}

\begin{proposition}[B-RD16: necessary lower bounds for character determination]
\label{prop:B-RD16}
\status{Proved.}
Let $h\geq2$. For $p\geq5$ and absolute ramification $e$, let $N$
be the largest bad elliptic level from Corollary~\ref{cor:B-RD13}.
If $N\geq1$, any universal least index $\nu$ in B.4 must satisfy
$\nu\geq N$. More precisely, for each $1\leq n\leq N$ there are two
height-$h$ BT groups satisfying the manuscript's monodromy condition,
with isomorphic $p^n$-torsion, whose characters on inertia elements of
order prime to $p$ differ. For $p=3$, the same lower bound holds using
the tame value $m_{\max}=4$ in that corollary.
\end{proposition}
\begin{proof}
Take the elliptic curve $E/K$ supplied by Theorem~\ref{thm:B-RD12}
at level $n$, and a canonical ordinary elliptic lift $E'/K$.
Lemma~\ref{lem:B-RD14} gives $\mu_{p^n}\subset K$, so both
$E[p^n]$ and $E'[p^n]$ are constant of rank $p^{2n}$.
Their inertia characters differ: the ordinary good curve has trace
two, whereas the nontrivial tame automorphism of the supersingular
good reduction has trace different from two. Append the same constant
etale BT group of height $h-2$ to both. Their characters each increase
by $h-2$, preserving the difference, and their truncated groups remain
isomorphic. These groups are potentially good, hence satisfy the
stated monodromy condition without using any extension argument.
Finally the manuscript's $M(\nu)$ is $M[p^{\nu+1}]$; failure at
$n=N$ gives $\nu\geq N$.
\end{proof}

\subsection{The monodromy condition and the Fontaine comparison}

\begin{proposition}[B-RD4: Lead B' needs an additional hypothesis]\label{prop:B-RD4}
\status{Refuted as stated.}
For every odd prime $p$, there is a height-two Barsotti--Tate group over
$K=\mathbf Q_p(\zeta_p)$ satisfying the manuscript's monodromy condition
whose rational Tate module is not potentially Hodge--Tate, and therefore
is not potentially semistable. Its manuscript inertia character is the
constant function $2$. Thus the monodromy condition in 9:11 cannot be
identified with potential semistability, and Fontaine's usual
Weil--Deligne representation need not be defined for its Tate module.
\end{proposition}
\begin{proof}
Write $\chi:G_K\to\mathbf Z_p^\times$ for the cyclotomic character.
Since $\mu_p\subset K$, it takes its values in $1+p\mathbf Z_p$. Put
\[
 \rho(g)=\begin{pmatrix}1&p^{-1}\log\chi(g)\\0&1\end{pmatrix}.
\]
This is a continuous representation on $T=\mathbf Z_p^2$, because the
logarithm is additive on $1+p\mathbf Z_p$ and takes values in
$p\mathbf Z_p$. Over a characteristic-zero field, the equivalence
between finite etale group schemes and finite continuous Galois modules
applied to $T/p^nT$ gives compatible finite etale group schemes $M[p^n]$.
The transition maps are induced by multiplication by $p$ on the lattices,
and the resulting inductive system is a height-two Barsotti--Tate group.
The invariant first coordinate and the quotient both give the constant
group $\mathbf Q_p/\mathbf Z_p$. They extend over $\mathcal O_K$, so the
specified two-step filtration satisfies the manuscript condition and
its character is $1+1=2$.

For $V=T\otimes\mathbf Q_p$, the Sen operator is
\[
 \Theta_V=\frac1p\begin{pmatrix}0&1\\0&0\end{pmatrix}.
\]
Indeed the logarithm of a unipotent matrix $1+aN$, with $N^2=0$, is
$aN$, and division by $\log\chi(g)$ gives the displayed operator.
It remains nonzero and nilpotent after restriction to any finite
extension. Sen's criterion says that a representation is Hodge--Tate
exactly when this operator is diagonalizable with integral eigenvalues;
the identical unscaled example and the criterion are explicitly given
in \cite[\S II.1.2, p.~9 of the author's version]{Berger2004Introduction}.
The operator is not diagonalizable. Potentially semistable
representations are potentially Hodge--Tate, so no finite extension
makes this representation semistable.
\end{proof}

\begin{proposition}[The corrected Fontaine dictionary]
\label{prop:B-RD4-dictionary}
\status{Proved with the additional potential semistability hypothesis.}
Let $M$ satisfy the manuscript monodromy condition, and assume in
addition that $V=T_p(M)\otimes_{\mathbf Z_p}\mathbf Q_p$ is potentially
semistable. For arbitrary perfect residue field, its manuscript
character is the inertia trace of Fontaine's $p$-module of Deligne
$D_{\mathrm{pst}}(V)$. When the residue field is finite, this is the
inertia character of the usual Weil--Deligne representation. With
\[
 D_{\mathrm{cov}}(G)=D_{\mathrm{contra}}(G^D),
\]
the precise comparison for a good graded group $G$ over a finite
extension $L/K$ is
\[
 D_{\mathrm{cov}}(G_{k_L})[1/p]
       \simeq D_{\cris,L}(T_pG\otimes\mathbf Q_p)(-1).
                                                        \tag{B-RD4.1}
\]
The Tate twist in this formula does not change the inertia trace.
\end{proposition}
\begin{proof}
Choose a finite Galois extension $L/K$ over which $V$ is semistable
and every graded group has good reduction. All subrepresentations
and quotients in the filtration are semistable over $L$.
Crystalline comparison for a BT group, with the indicated
covariant-versus-contravariant convention, gives (B-RD4.1)
functorially and compatibly with descent; see
\cite[Lemma~2.2.4 and Corollary~2.2.6, author pp.26--27; Appendix~A.2
and~A.6]{Kisin2006Crystalline}. One can check the convention on
$\mathbf Q_p/\mathbf Z_p$ and $\mu_{p^\infty}$: their covariant
Frobenius slopes are respectively one and zero, whereas the
crystalline Tate-module slopes are zero and minus one.

The semistable period functor is exact on semistable
representations. Thus its filtration has graded terms the
crystalline period modules of the good graded groups, and traces
of every finite descent-inertia element add. Exactness and the
passage from the completed maximal unramified extension back to
$K$ are treated in \cite[\S\S5.1.5--5.1.7,
pp.156--158]{Fontaine1994Semi}; alternatively the exact
$D_{\mathrm{pst}}$ functor is stated in
\cite[\S2.3.4, p.337]{Fontaine1994WD}.
The crystalline Tate twist is unramified in this period-module
sense, so the right side of (B-RD4.1) has the same inertia trace as
its untwisted version. Summing proves the assertion.

Over an arbitrary perfect residue field the natural object is the
$p$-module of Deligne, with semilinear Galois action, Frobenius and
monodromy; its definition is
\cite[\S\S1.1.1--1.1.2, pp.322--323]{Fontaine1994WD}.
If the residue field is finite, the corresponding linear
Weil--Deligne action is
$\rho_0(w)=\varphi^{-a(w)}\rho(w)$, as in
\cite[\S1.3.5, p.329; \S2.3.7, pp.339--340]{Fontaine1994WD}.
For inertia $a(w)=0$, so the trace just computed is exactly the
Weil--Deligne inertia trace. No finite-residue-field hypothesis is
needed for the preceding Deligne-module statement.
\end{proof}

\begin{proposition}[B-RD5: the height-one case of B.4]\label{prop:B-RD5}
\status{Proved.}
Suppose $\mu_p\subset K$. For height $h=1$ and every prime $p$, the manuscript
character on prime-to-$p$ inertia elements is determined by $M(0)$.
The least index in its indexing is therefore $\nu=0$.
\end{proposition}
\begin{proof}
For $p=2$, the finite special inertia action is contained in
$(\mathbf Q_2^\times)_{\mathrm{tors}}=\{1,-1\}$. Every odd-order
inertia element therefore acts trivially. Its requested character is
identically one, so $\nu=0$ suffices. We now assume $p\geq3$.

A height-one group satisfying the monodromy condition itself acquires
good reduction. Its good model has dimension zero or one. In dimension
zero it is etale, and in dimension one its Cartier dual is etale.
Over the completed maximal unramified extension of a good-reduction
field its model is respectively $\mathbf Q_p/\mathbf Z_p$ or
$\mu_{p^\infty}$. Its finite descent inertia character therefore takes
values in the torsion subgroup of $\mathbf Z_p^\times$, namely
$\mu_{p-1}$. In dimension zero, reduction modulo $p$ gives precisely
this character. In dimension one, the Tate module acquires additionally
the cyclotomic factor, whose reduction is trivial under $\mu_p\subset K$.
Thus in both cases the finite character is the unique Teichmuller lift
of the one-dimensional mod-$p$ inertia character. The same conclusion
holds after summing characters of height-one potentially good graded
pieces: the trace on a prime-to-$p$ element is the Brauer character of
the mod-$p$ representation, and is independent of the chosen filtration.
\end{proof}


\subsection{A finite bound for the manuscript character in every height}

\begin{lemma}[Comparison after a common good-reduction extension]
\label{lem:B-RD17-common}
\status{Proved under stated hypotheses}
Let $M,N$ be height-$h$ BT groups over $K$, each equipped with a
filtration by BT subgroups whose graded pieces have good reduction
over the same finite Galois extension $L/K$. Let $c$ be an integer
such that $p^c$ times every generic-fibre homomorphism between finite
flat commutative $p$-group schemes over $\mathcal O_L$ extends over
$\mathcal O_L$. If
\[
 M[p^n]\simeq N[p^n],\qquad n>2ch^2,
                                                        \tag{B-RD17.1}
\]
then their manuscript characters agree on all prime-to-$p$ inertia
elements. No good-reduction assertion about the extensions between
the graded pieces is assumed.
\end{lemma}
\begin{proof}
Use the given isomorphism to identify the two generic finite groups
with $X$. The two BT filtrations induce filtrations
\[
 0=X_0\subset\cdots\subset X_r=X,
 \qquad 0=Y_0\subset\cdots\subset Y_s=X.
\]
Write $h_i$ and $h'_j$ for the heights of their respective graded BT
groups. Thus $\sum_i h_i=\sum_j h'_j=h$.
The Zassenhaus refinement has common generic pieces
\[
 Q_{ij}=
 \frac{X_i\cap Y_j}
 {(X_{i-1}\cap Y_j)+(X_i\cap Y_{j-1})}.
                                                        \tag{B-RD17.2}
\]
In the good model of the $i$th graded group at level $p^n$, take
schematic closures of the subgroup filtration induced by the
$Y_j$. Consecutive quotients give finite flat models $G_{ij}$ of
$Q_{ij}$. Refine the good model of the $j$th graded group from $N$
in the same way, obtaining another model $H_{ij}$. These quotients
need not be truncated BT groups; only their finite flatness is used.
All closures and quotients are stable under the descent action.
The descent actions on the good BT models exist uniquely by Tate's
extension theorem for homomorphisms over a normal base.

Let $k_L$ be the residue field, and extend it algebraically if
necessary. The special Dieudonne module of $G_{ij}$ is a subquotient
of the free $W(k_L)/p^n$-module belonging to the $i$th graded BT
model. Consequently it can be generated by at most $h_i$ elements
over the DVR $W(k_L)$. Similarly the module for $H_{ij}$ has at most
$h'_j$ generators. Here exactness of Dieudonne theory over a perfect
field is applied only to special fibres of exact finite flat
sequences.

The generic identification of $G_{ij}$ and $H_{ij}$, and its inverse,
extend after multiplication by $p^c$. The two composites are
$p^{2c}$, because this equality holds generically and the source
coordinate algebras are flat. Uniqueness of these extensions makes
them compatible with inertia. Fix a prime-to-$p$ element $\gamma$
of inertia, of order $m$. After extending the perfect residue field
so that all $m$th roots of unity are present, the integral
idempotents
\[
 \frac1m\sum_{a=0}^{m-1}\zeta^{-a}\gamma^a
\]
split every special Dieudonne module into $\zeta$-eigensummands.
On a fixed eigensummand, the kernel of either comparison map is
annihilated by $p^{2c}$, and its cokernel is also annihilated by
$p^{2c}$. A module with at most $d$ generators over a DVR has
$p^{2c}$-torsion of length at most $2cd$, and any quotient killed
by $p^{2c}$ has length at most $2cd$. It follows that the two
$\zeta$-eigenlengths differ by at most
\[
 2c\max(h_i,h'_j)\leq2ch_i h'_j.
                                                        \tag{B-RD17.3}
\]
(The last inequality uses that all nonzero graded heights are
positive integers.)

Exactness along each refinement adds these lengths. A good BT group
of height $d$ has a free integral Dieudonne module of rank $d$, and
its $p^n$-torsion contributes $n$ times the multiplicity of each
$\gamma$-eigenvalue in that module. Thus, if $a_\zeta,b_\zeta$ are
the multiplicities in the sums of graded isocrystals for $M,N$,
\[
 \bigl|n(a_\zeta-b_\zeta)\bigr|
 \leq 2c\sum_{i,j}h_i h'_j=2ch^2.
\]
Under (B-RD17.1), the integer $a_\zeta-b_\zeta$ is zero. This holds
for every eigenvalue and every prime-to-$p$ inertia element, proving
the assertion. In particular, applying the same argument to two
filtrations of a single $M$, with arbitrarily large $n$, proves
independence of the prime-to-$p$ character from the chosen filtration.
\end{proof}

\begin{lemma}[A uniform ramification degree for the graded pieces]
\label{lem:B-RD17-degree}
\status{Proved under stated hypotheses}
Put $d=4h^2$ and
\[
 D_{p,h}=
 \begin{cases}
  |\GL_d(\mathbf F_p)|,&p>2,\\
  |\GL_d(\mathbf Z/4\mathbf Z)|
     =2^{d^2}|\GL_d(\mathbf F_2)|,&p=2.
 \end{cases}
\]
For two height-$h$ groups satisfying the monodromy condition, their
graded pieces have simultaneous good reduction over a finite Galois
extension $L/K$ with ramification index at most $D_{p,h}$.
\end{lemma}
\begin{proof}
Choose any finite Galois good-reduction extension $L_0/K$ for all
the graded pieces. Let $I_0$ be its inertia group. Their direct sum
has height $2h$; denote its special Dieudonne isocrystal by $D$.
The image $H$ of $I_0$ lies in
$\End_F(D)^\times$, where $F$ denotes semilinear Frobenius.
This endomorphism algebra has $\mathbf Q_p$-dimension at most
$(2h)^2=d$. Here is an elementary justification not requiring a
classification of isocrystals. In any semilinear Frobenius module,
linearly independent fixed vectors over the fixed field
$W(k_{L_0})[1/p]^{F}=\mathbf Q_p$ remain linearly independent over
$W(k_{L_0})[1/p]$: normalize a shortest relation to have one
coefficient $1$, apply Frobenius, and subtract. A shorter relation
would result unless every coefficient were fixed. Apply this to
conjugation by Frobenius on the endomorphism module, of rank $d$.

Left multiplication embeds $H$ in a general linear group over
$\mathbf Q_p$ of degree at most $d$. A finite group preserves a
$\mathbf Z_p$-lattice, obtained by summing the translates of any
lattice. For $p>2$, the kernel of reduction modulo $p$ in
$\GL_d(\mathbf Z_p)$ is torsion free. For $p=2$, the kernel of
reduction modulo $4$ is torsion free. The $p$-adic logarithm, or the
binomial valuation argument, proves both assertions. Thus
$|H|\leq D_{p,h}$.

Let $J$ be the kernel of $I_0\to H$. It is normal in
$\operatorname{Gal}(L_0/K)$, and $L=L_0^J$ is Galois over $K$ with
$e(L/K)=|H|$. Each rational graded Tate module is crystalline over
$L$. Indeed crystalline comparison identifies its special
isocrystal, with its descent action, with the crystalline period
module, up to the dual or Tate twist dictated by the convention.
The twist has trivial inertia action. The group $J$ acts trivially,
and $L_0,L$ have the same residue field. Hence
\[
 D_{\cris,L}(V)=D_{\cris,L_0}(V)^J
\]
has full dimension. The period-invariant and unramified-completion
statements used here are also explicitly treated in
\cite[\S\S5.1.5--5.1.7, pp.156--158]{Fontaine1994Semi}.
The required comparison, including functoriality,
is given in \cite[Lemma~2.2.4 and the proof of Corollary~2.2.6]{Kisin2006Crystalline}.
Corollary~2.2.6 of that paper, valid for arbitrary perfect residue
field and every prime, now gives a good BT group with this rational
Tate module.

For completeness, the actual graded lattice, rather than just its
rational representation, also has a good model. If $T_0$ is the Tate
lattice of one good model, rescale the desired lattice $T$ so that
$p^aT_0\subset T\subset T_0$. Take the schematic closure of
$T/p^aT_0$ in the finite flat group of $p^a$-torsion, and quotient the
good BT group by this finite flat subgroup. Its Tate lattice is
$p^{-a}T$, which is isomorphic to $T$. Finite flat quotients of BT
groups give BT groups, level by level. This argument uses only the
rational existence theorem at $p=2$, and does not assume an integral
$p=2$ equivalence in Kisin's Theorem~2.2.7.
\end{proof}

\begin{theorem}[B-RD17: a finite universal index exists]
\label{thm:B-RD17}
\status{Proved; independent proof and primary-source audit passed.}
For every prime $p$, every absolute ramification index $e$, and every
height $h\geq1$, a finite universal character-determination index in
B.4 exists. With $D_{p,h}$ as above, set
\[
 q=\left\lceil\log_p(eD_{p,h})\right\rceil,
 \qquad C=(2q+2)(q+2).
\]
An explicit sufficient index in the manuscript's convention is
\[
                       \nu=2Ch^2.                    \tag{B-RD17.4}
\]
In other words, the prime-to-$p$ manuscript character is determined
by $M[p^{\nu+1}]$ at this index. This bound is not asserted to be
sharp. For every prime and $h=1$, Proposition~\ref{prop:B-RD5} gives
the sharp value $\nu=0$; Proposition~\ref{prop:B-RD18} supplies
nonzero lower bounds for every admissible ramification in higher
height.
\end{theorem}
\begin{proof}
Choose $L$ as in Lemma~\ref{lem:B-RD17-degree}. Its absolute
ramification $e_L$ satisfies $e_L\leq eD_{p,h}$.
By \cite[Corollary~2 and the bound following Theorem~1, p.2 of the
author version]{VasiuZink2012Boundedness}, the generic Hom extension
exponent can be chosen to satisfy
\[
 c\leq(\log_p e_L+v_p(e_L)+2)(v_p(e_L)+2)
       \leq(2q+2)(q+2)=C.
\]
At level $n=2Ch^2+1$, Lemma~\ref{lem:B-RD17-common} applies.
Since $M(\nu)=M[p^{\nu+1}]$, this is precisely (B-RD17.4).
\end{proof}

\begin{remark}[The remaining sharp problem]\label{auto:B:5338697ffdba}
\status{Proved under stated hypotheses}

Theorem~\ref{thm:B-RD17} excludes the alternative that no finite
index exists. It does not determine the least index in height
$h\geq2$, and does not turn the manuscript monodromy condition into
potential semistability. The logarithmic cyclotomic extension in
Proposition~\ref{prop:B-RD4} remains a counterexample to that claim.
The refinement proof above works directly with good graded pieces
and therefore includes that extension and analogous nongood
extensions.
\end{remark}


\subsection{Stronger lower bounds from formal groups}

\begin{proposition}[B-RD18: every admissible ramification has a height-two counterexample]
\label{prop:B-RD18}
\status{Proved; independent construction and determinant audit passed.}
Let $p$ be any prime, let $e$ be a positive multiple of $p-1$, and
put $s=v_p(e/(p-1))$. Suppose $n\geq1$ satisfies
\[
 s\geq n-1,\qquad
 e\geq(p-1)p^{\max(2n-2,\,4n-4-s)}.                \tag{B-RD18.1}
\]
Then there are a complete discretely valued field $K$ with
algebraically closed residue field of characteristic $p$, absolute
ramification $e$, and $\mu_p\subset K$, and two height-two BT groups
$M,N$ over $K$ with
\[
 M[p^n]\simeq N[p^n],
\]
whose characters on prime-to-$p$ inertia differ. Both groups
individually acquire good reduction. The group $M$ can be taken to
have tame good-reduction defect $p+1$, and $N$ to be constant etale.
The same examples in every height $h\geq2$ are obtained by adding
constant etale groups.
\end{proposition}
\begin{proof}
Choose a height-two Honda formal group over
$k=\overline{\mathbf F}_p$. Its endomorphism ring contains
$W(\mathbf F_{p^2})$, so the Teichmuller units give a tame
order-$p+1$ automorphism $\alpha$ whose tangent eigenvalue $\zeta$
is primitive of that order. The two eigenvalues on its isocrystal
are $\zeta,\zeta^p=\zeta^{-1}$.

Apply the full-level Drinfeld deformation-ring construction in the
proof of Theorem~\ref{thm:B-RD12}, with $m=p+1$. Every ingredient
before the elliptic algebraization step involves only the formal
height-two group: the regular full-level ring, the two torsion
coordinates, averaging for the tame action, the product formula
for $p$, its valuation regions, and the equivariant discrete
valuation specialization. The condition $m\mid p+1$ holds, and its
arithmetic existence condition is exactly (B-RD18.1), because
$(p^2-1)/m=p-1$.

Stop at that formal-group construction and descend its generic BT
group along the cyclic tame extension. Over a characteristic-zero
field this is ordinary Galois descent for the compatible finite
etale groups at all levels. The full level structure is fixed by
the descent, so $M[p^n]$ is constant over $K$. The special descent
action is $\alpha$, hence is faithful of order $p+1$. Its trace at a
generator is $\zeta+\zeta^{-1}\ne2$.
The constant etale height-two group $N$ has the same constant
$p^n$-torsion and has inertia trace two. Adding constant etale
height $h-2$ adds $h-2$ to both traces and preserves their difference.

The proof also applies at $p=2$: the chosen order is three, so every
averaging denominator remains a unit; all Drinfeld product and
valuation formulas used in the construction hold for that prime.
Finally the cyclotomic containment can be checked directly, without
assuming a polarization of $M$. Over the good extension, the
rank-one crystalline representation $\det(T_pM)\otimes\mathbf Q_p$
has the single Hodge--Tate weight one. Kisin's
Corollary~2.2.6 realizes it by a height-one good group, which over
this strictly henselian complete base is multiplicative. Thus
$\det(T_pM)\chi_p^{-1}$ is a finite character. Its inertia type is
$\det(\alpha)=\zeta^{p+1}=1$, so it is unramified, hence trivial
because the residue field is algebraically closed. Therefore
$\det(T_pM)=\chi_p$. The full rational $p^n$-torsion of $M$ forces
$\chi_p\equiv1\pmod{p^n}$, proving $\mu_{p^n}\subset K$.
\end{proof}

\begin{corollary}[An explicit general-height lower bound]
\label{cor:B-RD18}
\status{Proved under stated hypotheses}
Let $h\geq2$, $\mu_p\subset K$, and define
\[
 L=\left\lfloor\log_p\!\left(\frac e{p-1}\right)\right\rfloor,
 \qquad
 N=1+\min\left\{s,\left\lfloor\frac L2\right\rfloor,
                         \left\lfloor\frac{L+s}4\right\rfloor\right\}.
\]
Any universal least index in B.4 satisfies $\nu\geq N$.
In particular $N\geq1$, so for every admissible $e$ and every
$h\geq2$, the first truncation $M(0)$ does not determine the
requested character. This is stronger than restricting the lower
bound to elliptic examples: the tame automorphism of a formal
height-two group need not arise from an automorphism of an
elliptic curve.
\end{corollary}
\begin{proof}
The integer equivalence between (B-RD18.1) and $n\leq N$ is the
same calculation as in Corollary~\ref{cor:B-RD13}, with
$m=p+1$. Failure at level $p^N=M(N-1)$ gives $\nu\geq N$.
\end{proof}

\begin{remark}[Finite arithmetic certificate]\label{auto:B:c9c8e4df645d}
\status{Proved under stated hypotheses}

The proof above is symbolic. As a separate finite check,
\path{code/verify_honda_bounds.py} enumerates the exact integer-pair
formula for $p=2,3,5,7,11,13,17,19$, levels $n=1,2,3$, and
$e\leq1200$, always with $m=p+1$. It verifies all 24 range tables,
3803 constructed pairs, and 11088 least-level formula comparisons
through $e\leq5000$. All checks pass; the recorded run took
$0.508$ seconds. This certificate verifies the integer arithmetic,
not the formal deformation theory.
\end{remark}


\subsection{A wild restriction and a sharp minimum in characteristic three}

\begin{lemma}[B-RD19: additive full odd torsion is potentially good]
\label{lem:B-RD19}
\status{Proved.}
Let $p$ be odd and let an elliptic curve $E/K$ have full rational
$p$-torsion. If its reduction is additive, it is potentially good.
Thus the additive existence question in B.3 is the potentially good
one, although its unrestricted good-forcing question also includes
multiplicative Tate curves.
\end{lemma}
\begin{proof}
The Weil pairing gives $\mu_p\subset K$. Suppose instead that
$v_K(j(E))<0$. The integral power series for $j(q)^{-1}$ starts with
$q$, so its compositional inverse gives $q\in K$, $v_K(q)>0$, with
$j(E_q)=j(E)$, where $E_q$ is the Tate curve
\cite[Theorem~1 and the $j$-expansion, pp.4--6]{Tate1995Uniformization}.
This value of $j$ is neither zero nor $1728$. Elliptic curves with
that same $j$ are quadratic twists, so $E=E_q\otimes\epsilon$ for
a quadratic character $\epsilon$. Tate uniformization gives
\[
       E[p]^{\mathrm{ss}}
        \simeq\overline\epsilon\,\overline\chi_p
                      \oplus\overline\epsilon.
\]
The left side is trivial and $\overline\chi_p=1$, so
$\overline\epsilon=1$. As $p$ is odd, reduction modulo $p$ is
injective on $\{1,-1\}$, hence $\epsilon=1$. The curve is therefore
split multiplicative, contrary to additive reduction. Thus $j(E)$
is integral, which is the elliptic criterion for potential good
reduction.
\end{proof}

\begin{lemma}[The maximum formal valuation is divisible by three]
\label{lem:B-RD20-wild}
\status{Proved; independent proof and source audits by the parent and obstructions/crystals agents passed.}
Suppose $p=3$, $E[3^n]\subset E(K)$, and $E$ has additive reduction
with wild semistability defect $m\in\{3,6,12\}$.
After extending the residue field algebraically, let $L/K$ be the
minimal good-reduction extension, with its normalized valuation
$v_L$. The maximum valuation $A$ of a primitive $3^n$-torsion
parameter in the good formal group is divisible by three.
\end{lemma}
\begin{proof}
First the discriminant of any Weierstrass equation for $E$ is a
cube in $K^\times$. The curve $E_0:y^2+y=x^3$ has discriminant
$-27=(-3)^3$ and full rational $3$-torsion over $K\supset\mu_3$,
as checked in Theorem~\ref{thm:B-RD2}. Choose a symplectic
isomorphism $E_0[3]\simeq E[3]$. The universal property of the
Hesse pencil makes $E$ $K$-isomorphic to a member of the Hesse
pencil of $E_0$ \cite[Theorem~1.1, p.2; proof, pp.7--8]
{AnemaTopTuijp2018Hesse}. Completing the square gives the short
equation $y^2=x^3+1/4$ with the same discriminant. Proposition~5.1
of that source, pp.4--5, computes the discriminant of every finite
smooth member as $\Delta(E_0)(\det A_t)^3$; the member at infinity
is $E_0$ itself. Thus every such discriminant is a cube.
A change of Weierstrass coordinates changes the discriminant by a
twelfth power. Hence $\delta=v_K(\Delta)$ is divisible by three.

Write the relation between a $K$-rational invariant differential and
an integral generator on the good model over $L$ as
$\omega_L=u\omega_K$. Since the good discriminant is a unit,
\[
                  v_L(u)=m\delta/12.
\]
This is an integer divisible by three, for each
$m=3,6,12$. Let $\sigma$ generate the wild subgroup $C_3$, and set
$F=L^{C_3}$. Choose $b_0\in F^\times$ with
$v_L(b_0)=v_L(u)$. Replacing the integral differential by
$(b_0/u)\omega_L$ makes it $C_3$-invariant without changing its
integrality, since $b_0/u$ is a unit. Choose a formal parameter with
this linear differential and odd under $[-1]$, which is possible
because two is a unit. Thus the wild descent series has linear
coefficient exactly one and only odd powers.

The special fibre is supersingular by Proposition~\ref{prop:B-RD3}.
On $y^2=x^3-x$ its order-three automorphisms are translations
$x\mapsto x+r$, $r\in\mathbf F_3^\times$. In the odd parameter
$t=-x/y$, their first nonidentity term has degree three and nonzero
coefficient. Therefore the wild descent series, or its inverse
according to the coordinate convention, has the form
\[
                   t\longmapsto t+c t^3+O(t^5),
                   \qquad c\in\mathcal O_L^\times.
\]
For a $K$-rational nonzero torsion point with parameter $x$ this gives
\[
                     v_L(\sigma x-x)=3v_L(x).       \tag{B-RD20.1}
\]

Let $b=v_L(\sigma\pi-\pi)-1$ be the positive lower ramification
break, for a uniformizer $\pi$ of $L$. The standard uniformizer
expansion, with coefficients in the fixed Cohen ring $W(k)$, gives
\[
 v_L(\sigma x-x)\geq v_L(x)+b,
 \quad\text{with equality if }3\nmid v_L(x).       \tag{B-RD20.2}
\]
Indeed the term of lowest valuation in $x$ gives equality when its
exponent is prime to three, and every higher term has strictly
larger valuation in its difference. If $B$ is the minimum primitive
valuation, (B-RD20.1)--(B-RD20.2) imply $b\leq2B$. If $A>B$ and
$3\nmid A$, the same equations give $b=2A>2B$, a contradiction.

It remains to exclude $A=B$ with $3\nmid A$. In that case $b=2A$.
For every primitive torsion parameter $x$, write
$x=a_x\pi^A+O(\pi^{A+1})$, with $a_x\in k^\times$. The first
coefficient of $\sigma x-x$, divided by $\pi^{3A}$, is a fixed
nonzero scalar times $a_x$; the first coefficient of the right
side of the descent equation is $\bar c a_x^3$. Thus all these
coefficients solve one equation
\[
                           \bar c X^3=dX,
                           \qquad d\ne0.
\]
Here $d=A\overline{(\sigma\pi-\pi)/\pi^{b+1}}\ne0$, since
$3\nmid A$. Its roots have dimension at most one over
$\mathbf F_3$. Choose a
primitive basis $P,Q$ of $E[3^n]$. Sending $(a,b)\in\mathbf F_3^2$
to the leading coefficient of $[\widetilde a]P+[\widetilde b]Q$
is well defined and linear: changing an integer lift by a multiple
of three does not affect the leading coefficient, because formal
addition is addition to first order and the residue characteristic
is three. This map is injective, since a nonzero combination is
primitive and cancellation would give valuation greater than the
maximum $A$. This requires dimension two, a contradiction. Hence
$3\mid A$ also in this case.
\end{proof}

\begin{proposition}[B-RD20: a stricter wild bound and an exact minimum]
\label{prop:B-RD20}
\status{Proved; independent proof and source audits by the parent and obstructions/crystals agents passed.}
For $p=3$ write $e=2u$ and $s=v_3(u)$. If full rational $3^n$-torsion
occurs with wild defect $m\in\{3,6,12\}$, necessarily
\[
 s\geq n-1,\qquad
 e\geq\frac8m\,3^{\max(2n-2,\,4n-4-s)}.        \tag{B-RD20.3}
\]
Allowing every perfect residue field and every defect, the least
absolute ramification index admitting an additive elliptic curve
with full rational $3^n$-torsion is exactly
\[
                            e_{\min}(n)=2\cdot3^{2n-2}.
                                                        \tag{B-RD20.4}
\]
In particular $e_{\min}(2)=18$: absolute indices six and twelve,
although compatible with the cyclotomic necessary condition,
do not admit an additive curve with full rational $9$-torsion.
\end{proposition}
\begin{proof}
The cyclotomic necessary condition gives $s\geq n-1$. The
valuation identity in the proof of Theorem~\ref{thm:B-RD12} uses
only the good height-two formal group, not the tame descent
congruences. In a valuation region
$3^{2r}B\leq A\leq3^{2r+2}B$, $0\leq r\leq n-1$, it therefore
still gives
\[
                  me=2\cdot3^{2n-r-2}(A+3^{2r+1}B).
                                                        \tag{B-RD20.5}
\]
The upper bound for $A$ is omitted when $r=n-1$.
By Lemma~\ref{lem:B-RD20-wild}, the parenthesis is divisible by
three. Since $v_3(m)=1$, this forces
$s\geq2n-r-2$, or $r\geq\max(0,2n-2-s)$.
The lower bound on $A$ and $B\geq1$ give
$me\geq8\cdot3^{2n+r-2}$. Substitute the least possible $r$ to
obtain (B-RD20.3).

To prove (B-RD20.4), tame defects satisfy the lower bound
$e\geq2\cdot3^{2n-2}$ by Theorem~\ref{thm:B-RD12} with $m\leq4$.
For a wild defect use $m\leq12$ in (B-RD20.3). If
$s\geq2n-2$, divisibility of $e$ by $2\cdot3^s$ already gives
the claimed bound. Otherwise $s\leq2n-3$, and the second exponent
in (B-RD20.3) is at least $2n-1$. Consequently
\[
 e\geq\frac23\,3^{2n-1}=2\cdot3^{2n-2}.
\]
Finally the tame construction with $m=4$ and
$e=2\cdot3^{2n-2}$ meets the criterion of
Theorem~\ref{thm:B-RD12}, so it attains this bound for every $n$.
Lemma~\ref{lem:B-RD19} covers every additive curve in the question.
\end{proof}

\begin{corollary}[B-RD21: the remaining wild question is one level]
\label{cor:B-RD21}
\status{Proved; independently audited by the parent agent; arithmetic certificate supplied.}
Fix $p=3$ and a positive even absolute ramification index $e$.
Let $N(e)$ be the largest $n\geq1$ for which an additive elliptic
curve with full rational $3^n$-torsion exists over some such field
with perfect residue field. Put
\[
 s=v_3(e/2),\qquad L=\lfloor\log_3(e/2)\rfloor,
\]
and define
\begin{align*}
 N_{\mathrm t}
   &=1+\min\left(s,\left\lfloor\frac L2\right\rfloor,
                         \left\lfloor\frac{L+s}{4}\right\rfloor\right),\\
 N_{\mathrm w}
   &=1+\min\left(s,\left\lfloor\frac{L+1}{2}\right\rfloor,
                         \left\lfloor\frac{L+s+1}{4}\right\rfloor\right).
\end{align*}
Then
\[
             N_{\mathrm t}\leq N(e)\leq N_{\mathrm w}
                             \leq N_{\mathrm t}+1.       \tag{B-RD21.1}
\]
In particular the full answer is already determined whenever
$N_{\mathrm t}=N_{\mathrm w}$. A gap occurs precisely when there is
an integer $k\geq1$ such that
\[
             k\leq s\leq2k-1,\qquad L=4k-s-1.          \tag{B-RD21.2}
\]
Then $N_{\mathrm t}=k$ and $N_{\mathrm w}=k+1$, so only existence
at level $n=k+1$ remains undecided. For full $9$-torsion, the only
indices not decided by these results are
\[
                              e=24,30,42,48.            \tag{B-RD21.3}
\]
At $e=24,30$ an example, if one exists, must have defect twelve;
at $e=42,48$ its possible wild defects are six and twelve.
\end{corollary}
\begin{proof}
The sharp tame criterion with defect four gives the lower bound
$N_{\mathrm t}$. The bound in Proposition~\ref{prop:B-RD20} is
weakest for defect twelve; its logarithmic inequalities are
$n-1\leq s$, $2n-2\leq L+1$, and $4n-4-s\leq L+1$.
They give the upper bound $N_{\mathrm w}$. Increasing $L$ by one
increases each displayed minimum by at most one.

Since $s\leq L$, a level $k+1$ admitted by the wild inequalities
but excluded by the tame ones must satisfy
\[
 k\leq s,\qquad \max(2k,4k-s)=L+1.
\]
The maximum cannot be $2k$ with $s\geq2k$, because that would
give $L=2k-1<s$. Thus $s\leq2k-1$ and $L=4k-s-1$.
Conversely these conditions make the two minima respectively
$k-1$ and $k$, which proves (B-RD21.2).
For $k=1$ we have $s=1,L=2$, so $e/2$ is a multiple of three
but not nine in the interval $[9,27)$; the four possibilities
are $12,15,21,24$. Substitute them into (B-RD20.3) for the
defect assertions. The finite integer certificate
\path{code/verify_p3_higher_wild_bounds.py} checks these formulae,
their extrema, and their compatibility with the unsimplified
valuation identity in the stated finite range.
\end{proof}


\subsection{An unramified height-two obstruction beyond the first level}

\begin{proposition}[B-RD22: four-torsion does not suffice at $p=2,e=1$]
\label{prop:B-RD22}
\status{Proved; exact finite certificate and independent parent and obstructions/crystals audits passed.}
For $p=2$, $e=1$, and every height $h\geq2$, the least index in
B.4 is at least two. In the manuscript's indexing, even
$M(1)=M[4]$ does not determine the character on inertia elements of
odd order. This already occurs over the unramified quadratic field
$K=\mathbf Q_2(\zeta_3)$.
\end{proposition}
\begin{proof}
Let
\[
 E/K:\quad y^2=x(x-1)(x+3),\qquad
 M=E[2^\infty],\qquad
 N=\mu_{2^\infty}\oplus\mathbf Q_2/\mathbf Z_2.
\]
The group $N$ extends to a split Barsotti--Tate group over
$\mathcal O_K$. Put $a=1+2\zeta_3$, so that $a^2=-3$, and choose
$i^2=-1$. The points
\[
       P=(a,ia(a-1)),\qquad Q=(3,6)
\]
satisfy $2P=(0,0)$ and $2Q=(1,0)$, and hence form a basis of
$E[4]$. Every element of $G_K$ fixes $a$ and sends $i$ to either
$i$ or $-i$. Consequently it sends $P$ to $\chi_2 P$ modulo four
and fixes $Q$. The equivalence between finite etale group schemes
over the characteristic-zero field $K$ and finite $G_K$-sets gives
an isomorphism of group schemes
\[
                     M[4]\simeq\mu_4\oplus\mathbf Z/4
                              =N[4].                 \tag{B-RD22.1}
\]

For the displayed integral Weierstrass equation the exact invariants
are
\[
 \Delta=2^8\,3^2,\qquad c_4=2^4\,13,\qquad
 j=\frac{2^4\,13^3}{9}.
\]
The discriminant valuation is eight, which is minimal since every
integral change to another Weierstrass equation changes it by a
multiple of twelve and a discriminant of an integral equation has
nonnegative valuation. In particular the curve is bad. Its integral
$j$ gives potential good reduction
\cite[Section~2, p.496]{SerreTate1968Good}. Since $v_K(j)=4>0$,
the good special fibre has $j=0$. In characteristic two its equation
has the form $y^2+a_3y=x^3+a_4x+a_6$, with $a_3\ne0$
\cite[Section~9.1, p.8]{AnemaTopTuijp2018Hesse}. Negation sends
$(x,y)$ to $(x,y+a_3)$, so there is no nonzero geometric
two-torsion point. Thus the good fibre is supersingular.

Pass to the completed maximal unramified extension and let $L/K$
be a minimal good-reduction extension, with inertia image of order
$m$. A change to a unit-discriminant equation over $L$ shows that
\[
                            12\mid8m,
\]
so $3\mid m$. The inertia image therefore contains an automorphism
of order three. On the supersingular special elliptic curve its
action on the rational Dieudonne module has eigenvalues $\zeta_3$
and $\zeta_3^{-1}$, hence trace $-1$: equivalently, use the standard
supersingular equation $y^2+y=x^3$ and its order-three automorphism
$(x,y)\mapsto(\zeta_3x,y)$; the two eigenvalues are interchanged by
Frobenius. The order-three trace of $N$, which has good reduction,
is $2$. Thus (B-RD22.1) does not determine the requested character.
Both groups satisfy the manuscript's monodromy condition.

For height $h>2$, add $h-2$ copies of the constant etale group
$\mathbf Q_2/\mathbf Z_2$ to both sides. Their equal traces add the
same constant, so the discrepancy remains. The entire finite
four-torsion calculation is checked with exact rational arithmetic
in $\mathbf Q(\sqrt{-3},i)$ by
\path{code/verify_p2_character_level2.py}, including all sixteen
points, all $256$ additions, and the Galois action.
\end{proof}


\subsection{The complete full-nine-torsion range}

\begin{theorem}[B-RD23: the four remaining full-nine cases exist]
\label{thm:B-RD23}
\status{Proved; independent parent and obstructions/crystals proof/source audits passed.}
For $p=3$, an additive elliptic curve with full rational $9$-torsion
exists over a complete discretely valued field of absolute
ramification index $e$ and perfect residue field if and only if
\[
                              6\mid e,\qquad e\geq18.
                                                        \tag{B-RD23.1}
\]
The examples can all be taken over finite extensions of $\mathbf Q_3$.
In particular all four indices $24,30,42,48$ left by
Corollary~\ref{cor:B-RD21} occur, with semistability defect twelve.
\end{theorem}
\begin{proof}
Necessity follows from the cyclotomic condition and
Proposition~\ref{prop:B-RD20}. We first make the example at $e=18$
explicit over a finite local field. Let $K_0=\mathbf Q_3(i)$,
$i^2=-1$, and $E_0:y^2=x^3-x$. This has good supersingular
reduction and the endomorphism
$[i]:(x,y)\mapsto(-x,iy)$, defined over $K_0$. Put
$L=K_0(E_0[9])$, the finite Galois division field.

Every primitive nine-torsion parameter $t_0=-x/y$ has
$v_3(t_0)=1/72$. Indeed, the unramified height-two formal group's
$[3]$-series has linear coefficient three, its first unit
coefficient is in degree nine, and every intervening coefficient
is divisible by three. A nonzero three-torsion parameter therefore
has valuation $1/8$. A preimage under $[3]$ must have smaller
valuation, on which the degree-nine term has strictly least
valuation; hence its valuation is $1/72$.

On the other hand, $E_0[9]$ is free of rank one over
$\mathbf Z_3[i]/9$. To see this directly, reduction modulo three
has dimension two over $\mathbf F_3$ and the operator $i$ has
irreducible polynomial $X^2+1$; it is therefore one-dimensional
over $\mathbf F_9$. Nakayama's lemma and the equal cardinalities
then give the rank-one assertion modulo nine. Galois commutes
with $i$, so its faithful image has order at most
$|(\mathbf Z_3[i]/9)^\times|=72$. The preceding valuation
denominator forces ramification degree at least $72$.
Consequently $L/K_0$ is totally ramified of degree $72$ and this
Galois image is the full unit group.

Let $\sigma\in\operatorname{Gal}(L/K_0)$ act on $E_0[9]$ by $[i]$,
and let $K=L^{\langle\sigma\rangle}$. Choose a primitive point
$Q\in E_0[9]$ and put $u=-x(Q)/y(Q)$. It is a uniformizer of
$L$, and the exact coordinate formula for $[i]$ gives
$\sigma u=iu$. Thus $a=u^4\in K$ has $v_K(a)=1$. The curve
\[
                         E:y^2=x^3-ax
\]
has full rational nine-torsion: its points obtained from
$P\in E_0[9]$ are $(u^2x(P),u^3y(P))$, and both coordinates are
fixed by $\sigma$. The absolute ramification index of $K$ is
$72/4=18$. The curve becomes $E_0$ over $L$, while
$v_K(\Delta(E))=v_K(64a^3)=3$ forces exact defect four. This gives
the promised finite local-field example at the endpoint.

We next give a uniform explicit construction for every multiple
$e=6d>18$.

Choose $i^2=-1$, a primitive ninth root of unity $\eta$, and put
$\zeta=\eta^3$. Let
\[
 K=\mathbf Q_3(i,\eta,\pi),\qquad \pi^d=\eta-1.
                                                        \tag{B-RD23.2}
\]
The field $\mathbf Q_3(i)$ is the unramified quadratic extension,
and $\eta-1$ is a uniformizer of its degree-six cyclotomic
extension. The final polynomial in (B-RD23.2) is Eisenstein.
Thus $K$ has absolute ramification $e$, residue field $\mathbf F_9$,
and uniformizer $\pi$. Write $v(\pi)=1$.

Put
\[
 h=1+\pi^2,\quad D=3(4h^3-1),\quad
 y=\frac{-3+\sqrt D}{2(1-h^3)}.                         \tag{B-RD23.3}
\]
Here $v(h^3-1)=6$, $v(4h^3-1)=6$, and $v(D)=e+6$.
Moreover $3/\pi^e\equiv-1\pmod\pi$, by the Eisenstein
cyclotomic equation for $\eta-1$. Consequently
$D/\pi^{e+6}\equiv-1\pmod\pi$. This unit has a square root in
$K$ lifting either $i$ or $-i$, by ordinary Hensel lifting, since
two is a unit. This specifies the square root in (B-RD23.3).
The quadratic formula gives
\[
       h^3=1+\frac3y+\frac3{y^2},\qquad
       r:=v(y)=\frac e2-3.                             \tag{B-RD23.4}
\]
In particular $r>0$. Define
\[
 \lambda=-\eta^2,\qquad c=\frac{\lambda}{\zeta-1},\qquad
 d_0=(1+y)^3-1=y^3h^3,\qquad
 t=1+\frac{1-\zeta}{d_0}.                             \tag{B-RD23.5}
\]
Since $(\zeta-1)^2=-3\zeta$, one has
$c^3=(1-\zeta)/9$. Therefore
\[
 \frac{t-1}{9}=\left(\frac{c}{yh}\right)^3,
 \qquad
 \frac{t-\zeta}{9}
   =\left(\frac{c(1+y)}{yh}\right)^3.                 \tag{B-RD23.6}
\]

We show that $(t-\zeta^2)/9$ is also a cube. Its ratio with the
first expression in (B-RD23.6) is
\[
              a=1-\zeta^2d_0=1+\lambda^3d_0.
\]
Use $\beta_0=1+\lambda y$ as a first approximation to a cube root.
Then
\[
 a-\beta_0^3
      =3(\lambda^3-\lambda)y
                   +3(\lambda^3-\lambda^2)y^2.
                                                        \tag{B-RD23.7}
\]
The first coefficient in parentheses has valuation $e/6$ and the
second is a unit. Since $r>e/6$, the valuation in (B-RD23.7) is
\[
             e+q,\qquad q=e/6+r=2e/3-3>e/2.
\]
The polynomial
\[
       \frac{(\beta_0+\pi^q Z)^3-a}{3\pi^q}
\]
has integral coefficients. Its linear coefficient is the unit
$\beta_0^2$, its quadratic coefficient is $\beta_0\pi^q$, and
its cubic coefficient has valuation $2q-e>0$. Its reduction is
therefore a degree-one polynomial, which has a simple root over
$\mathbf F_9$. Hensel lifting gives $\beta\in K$ with
$\beta^3=a$. This argument uses the displayed rescaling; it does
not invoke the generally inapplicable stronger inequality
$v(a-\beta_0^3)>2v(3)$.

Now use the original full-level modular-curve equations of
Fisher \cite[Lemma~4.2(i), p.8; Theorem~4.3, equation~(9), p.9]
{Fisher2015Nine}. They state that the family over $Y(3)$ is
\[
             E_t:\quad y_0^2+3tx_0y_0+(t^3-1)y_0=x_0^3,
                                                        \tag{B-RD23.8}
\]
and the cover $X(9)\to X(3)$ is given by the three equations
\[
     t-\theta=9(u+v\theta+w\theta^2)^3,
                     \qquad \theta=1,\zeta,\zeta^2.    \tag{B-RD23.9}
\]
The three cube roots supplied above are respectively
$c/(yh)$, $c(1+y)/(yh)$, and $c\beta/(yh)$. An invertible Fourier
matrix over $K$ gives $u,v,w$ with precisely these three values.
Thus (B-RD23.9) gives a $K$-point of $Y(9)$. By the moduli
interpretation in Section~2, p.4 of that source, it gives
$E_t[9]\simeq\mu_9\oplus\mathbf Z/9$. Since $\mu_9\subset K$,
the entire $9$-torsion is rational. The parameter is not a cusp,
because the valuation calculation below gives $v(t)<0$.

Finally, (B-RD23.4)--(B-RD23.5) give $v(t)=9-e<0$. The exact
invariants of (B-RD23.8) are
\[
 c_4=9t(t^3+8),\qquad \Delta=27(t^3-1)^3,\qquad
 j=\frac{27t^3(t^3+8)^3}{(t^3-1)^3}.
\]
Hence
\[
                   v(j)=27,\qquad v(\Delta)=81-6e
                                             \equiv9\pmod {12}.
\]
The integral $j$ gives potential good reduction
\cite[Section~2, p.496]{SerreTate1968Good}, and the discriminant
congruence excludes good reduction. Thus the reduction is additive.
An explicit integral equation is obtained by the substitution
$x_0=\pi^{2(9-e)}X$, $y_0=\pi^{3(9-e)}Y$: its two nonzero
coefficients have valuations $e$ and zero, and its discriminant
has valuation $6e-27$.

At any of $e=24,30,42,48$, the discriminant congruence forces the
minimal good inertia order $m$ to be divisible by four. Its possible
orders in characteristic three are therefore four and twelve.
The sharp tame full-nine criterion excludes four at these indices,
so $m=12$. All displayed finite polynomial identities, the Fourier
inversion, and the valuation inequalities at the four previously
undecided indices are checked in
\path{code/verify_full9_construction.py}. The square-root and cubic
Hensel steps are justified above, rather than treated as numerical
approximations.
\end{proof}

\begin{remark}[Passage to all higher levels]\label{auto:B:b02333f8ed15}
\status{Proved under stated hypotheses}

This first settles level $n=2$ of the wild higher-torsion question,
including every exceptional band with $k=1$ in Corollary~\ref{cor:B-RD21}.
The first higher case beyond this construction is full $27$-torsion
at $e=216$. It is subsequently realized by the wild trace-invariant
construction of Theorem~\ref{thm:B-RD26-wild}, which settles every
higher wild level and yields Corollary~\ref{cor:B-RD27-exact-three}.
\end{remark}


\subsection{An explicit finite algebraic test for elliptic torsion existence}

\begin{proposition}[A bounded test; the output at the first wild gap is not yet evaluated]
\label{prop:B-RD24-finite-test}
\status{Proved; complete proof and original-source audit independently passed by the parent, with an independent rerun of the exact polynomial certificate.  No numerical output for $(p,e,n)=(3,216,3)$ is asserted.}
For every odd prime $p$ and positive integers $e,n$, the existence of a complete discretely valued field of absolute ramification $e$, perfect residue characteristic $p$, and an additive elliptic curve with full rational $p^n$-torsion is equivalent to the finite polynomial test specified below.  The test uses only arithmetic in finite-length Witt vectors, polynomial division and discriminants over $\mathbf Q$, and effective elimination over $\overline{\mathbf F}_p$.  If the test is affirmative, an example exists over a finite extension of $\mathbf Q_p$.
\end{proposition}

\par\smallskip\noindent\textbf{Constants and the division polynomial.}
Put $m=p^n$, $N=(m^2-1)/2$, and $a=ne$.  Write
\[
 \psi_m(X)=mX^N+\sum_{i=0}^{N-1}d_iX^i
\]
for the usual odd division polynomial (whose recurrences and torsion-root interpretation are recorded in \cite[section~1.1.3, pp.~8--9]{Burhanuddin2007}) of
\[
 y^2+a_1xy+a_3y=x^3+a_2x^2+a_4x+a_6.
\]
The long-equation version used here can be specified without any division in its coefficient ring.  Put
\[
\begin{gathered}
 b_2=a_1^2+4a_2,\quad b_4=a_1a_3+2a_4,\quad
 b_6=a_3^2+4a_6,\\
 b_8=a_1^2a_6+4a_2a_6-a_1a_3a_4+a_2a_3^2-a_4^2,\qquad
 B(X)=4X^3+b_2X^2+2b_4X+b_6.
\end{gathered}
\]
For odd $j$ let $P_j=\psi_j$, and for even $j$ let
$P_j=\psi_j/\psi_2$.  Their initial values are
\[
\begin{aligned}
 P_0&=0,\quad P_1=P_2=1,\\
 P_3&=3X^4+b_2X^3+3b_4X^2+3b_6X+b_8,\\
 P_4&=2X^6+b_2X^5+5b_4X^4+10b_6X^3+10b_8X^2\\
    &\hspace{7mm}+(b_2b_8-b_4b_6)X+b_4b_8-b_6^2.
\end{aligned}
\]
For subsequent indices the recurrences are
\[
 P_{2r}=P_r(P_{r+2}P_{r-1}^2-P_{r-2}P_{r+1}^2)
\]
and
\[
 P_{2r+1}=
 \begin{cases}
 P_{r+2}P_r^3-B^2P_{r-1}P_{r+1}^3,&r\text{ odd},\\
 B^2P_{r+2}P_r^3-P_{r-1}P_{r+1}^3,&r\text{ even}.
 \end{cases}
\]
They follow from the usual recurrences and the identity $\psi_2^2=B(X)$.
Completing the square and then translating $X$ to remove its quadratic coefficient over $\mathbf Q$ identifies these initial values and recurrences with the short-equation division polynomials in the cited source; thus it also gives their stated torsion-root interpretation.  All the displayed operations are integral polynomial operations.  In particular the $d_i$ can be calculated by this finite recursion.  Define
\[
 c_m=\frac{\operatorname{disc}\bigl(\psi_m(X;0,0,0,-1,0)\bigr)}
                  {64^{\,N(N-1)/6}}\ \in\mathbf Q^\times,
 \qquad t_m=v_p(c_m).
\]
This definition is an exact rational computation, not a conjectured closed formula for the discriminant constant.
Let
\[
 D_0=\begin{cases}6e+11,&p=3,\\11,&p\geq5,\end{cases}
 \qquad
 D_F=a(N-1)(N-2)+e\,t_m+\frac{N(N-1)}6D_0,
\]
and take the deliberately generous positive integer
\[
 D=2|D_F|+a(4N+3)+ND_0+1,\qquad
 H=\max\{2D+1,\,3e+7,\,D_0+1\},\qquad q=\lceil H/e\rceil.
\]
The proof below shows that $D_F$ is in fact nonnegative and bounds the valuation of the monic polynomial discriminant used in the test.  Absolute values in the definition of $D$ merely make positivity manifest before this verification.

\par\smallskip\noindent\textbf{The finite coefficient ring.}
All quantified residue-field variables range over $k=\overline{\mathbf F}_p$.
Choose $e$ Witt vectors $c_0,\ldots,c_{e-1}\in W_q(k)$ with the zeroth Witt coordinate of $c_0$ nonzero, and form the finite-length ring
\[
 R=W_q(k)[\pi]\big/\left(\pi^e+pC(\pi)\right),
 \qquad C(\pi)=\sum_{i=0}^{e-1}c_i\pi^i.
\]
This ring is free of rank $e$ over $W_q(k)$, so all its operations are explicit polynomial operations in $eq$ residue-field coordinates.  Its maximal ideal is $(\pi)$ and its length is $eq\geq H$.  Membership in $\pi^rR$ is the finite polynomial condition
\[
                  z\in\pi^rR\quad\Longleftrightarrow\quad
                  (\exists w\in R)\ z=\pi^rw.
\]
Thus every condition below is a finite first-order polynomial formula over $k$.

\par\smallskip\noindent\textbf{The test.}
Ask whether there exist such $c_i$, coefficients $a_1,a_2,a_3,a_4,a_6\in R$, an integer $1\leq\delta\leq D_0$, and $N$ pairs $(X_i,Y_i)\in R^2$ satisfying the following conditions.

\begin{enumerate}
\item The universal invariants satisfy
\[
 \Delta\in\pi^\delta R\setminus\pi^{\delta+1}R,
 \qquad c_4^3\in\pi^\delta R.
\]
\item There are no $S,R_1,T_1\in R$ for which all five expressions
\[
\begin{aligned}
 U_1&=a_1+2S,\\
 U_2&=a_2-Sa_1+3r-S^2,\\
 U_3&=a_3+ra_1+2t,\\
 U_4&=a_4-Sa_3+2ra_2-(t+rS)a_1+3r^2-2St,\\
 U_6&=a_6+ra_4+r^2a_2+r^3-ta_3-rta_1-t^2,
\end{aligned}
\qquad r=\pi^{-e}R_1,\quad t=\pi^{-e}T_1
\]
satisfy $\pi^{3e}U_j\in\pi^{3e+j}R$, for $j=1,2,3,4,6$.
This notation contains no inversion of $\pi$ in $R$: first expand each $\pi^{3e}U_j$ and substitute $r=\pi^{-e}R_1,t=\pi^{-e}T_1$; every resulting exponent of $\pi$ is nonnegative.
\item Form the two integral polynomials
\[
\begin{aligned}
 F(X)&=X^N+\sum_{i=0}^{N-1}
       (-C(\pi))^n\pi^{a(N-i-1)}d_iX^i,\\
 G(X,Y)&=Y^2+(a_1\pi^aX+a_3\pi^{2a})Y\\
 &\hspace{9mm}-\bigl(\pi^aX^3+a_2\pi^{2a}X^2
                   +a_4\pi^{3a}X+a_6\pi^{4a}\bigr).
\end{aligned}
\]
For every $i$ require
\[
 F(X_i),G(X_i,Y_i)\in\pi^HR,\qquad
 F'(X_i)G_Y(X_i,Y_i)\notin\pi^{D+1}R.
\]
For $i\ne j$ require $X_i-X_j\notin\pi^{D+1}R$.
\end{enumerate}
Here $R_1$ in condition~2 is a variable, not the coefficient ring $R$.
The test returns \emph{yes} exactly when this finite formula is true in an algebraically closed field of characteristic $p$.

\begin{proof}
We give the bounds and the exact approximation argument, since these distinguish the test from an enumeration of unknown local fields or unknown infinite lifts.
Lemma~\ref{lem:B-RD19} proves that full rational odd-prime torsion excludes additive potentially multiplicative reduction.  Consequently every additive curve in the statement is potentially good, so the following integral-$j$ bounds cover the whole stated problem.

\emph{Minimal discriminants.}
Let an integral minimal equation have potentially good reduction, and put $\delta=v(\Delta)$.  Since $j=c_4^3/\Delta$ is integral, $v(c_4)\geq\delta/3$.  The identity $c_4^3-c_6^2=1728\Delta$ then gives $v(c_6)\geq\delta/2$.  Its short equation has coefficients $-c_4/48$ and $-c_6/864$.  If $p=3$ and $\delta\geq6e+12$, their valuations are at least $4$ and $6$; scaling $x=\pi^2x',y=\pi^3y'$ gives an integral equation with discriminant valuation $\delta-12$, a contradiction.  For $p\geq5$ the same argument applies to $\delta\geq12$.  Hence $\delta\leq D_0$.

An integral equation is nonminimal exactly when a change of variables with scaling factor $\pi$ produces another integral equation.  Indeed, a change with factor $\pi^b$, $b\geq1$, also gives one with factor $\pi$ using the same translations; unit scaling does not affect integrality.  The standard change-of-variables formulas are precisely the $U_j$ above.  The first condition gives $s\in\OO$ because $2$ is a unit.  The second gives $v(r)\geq-v(3)\geq-e$; the third then gives $v(t)\geq-e$.  Conversely these are all the possible translations.  Clearing denominators by $\pi^{3e}$ proves that condition~2 is exactly minimality and depends only on the displayed finite precision.

\emph{The two discriminant bounds.}
Over $\mathbf Q[a_1,a_2,a_3,a_4,a_6,\Delta^{-1}]$, the nonzero $m$-torsion modulo the free involution $P\mapsto-P$ is finite \'etale, with coordinate algebra defined by $\psi_m/m$.  Consequently $\operatorname{disc}(\psi_m)$ is a unit after inverting $\Delta$.  The universal Weierstrass discriminant is irreducible over $\mathbf Q$: after the invertible change to a short equation this is the irreducibility of $4A^3+27B^2$.  Its weighted degree is $12$, while the weighted degree of $\operatorname{disc}(\psi_m)$ is $2N(N-1)$.  It follows that
\[
           \operatorname{disc}(\psi_m)=c_m\Delta^{N(N-1)/6};
\]
specializing to $y^2=x^3-x$ proves exactly the formula defining $c_m$.

Over a lift to a complete DVR, $p/\pi^e=-C(\pi)^{-1}$, so $F(X)$ is exactly
\[
                 \frac{\pi^{aN}}{m}\psi_m(\pi^{-a}X).
\]
It is monic integral, and therefore
\[
 v(\operatorname{disc}F)
   =a(N-1)(N-2)+e\,t_m+\frac{N(N-1)}6\delta
   \leq D_F.
\]
The algebra
\[
                 \mathcal A=\OO[X,Y]/(F,G)
\]
is finite free of rank $2N=m^2-1$, with \'etale generic fibre equal to the coordinate algebra of $E[m]\setminus\{0\}$.
For completeness the needed second resultant has an elementary derivation.  Write
\[
 B(x)=4x^3+b_2x^2+2b_4x+b_6.
\]
After completing the square, at a two-torsion abscissa $\alpha$ the division-polynomial recurrences give
\[
 \psi_m(\alpha)=(-1)^{(m-1)/2}
         f'(\alpha)^{(m^2-1)/4},\qquad
 f(x)=x^3+(b_2/4)x^2+(b_4/2)x+b_6/4.
\]
The base cases are $\psi_1=1$ and $\psi_3(\alpha)=-f'(\alpha)^2$.
For $m=2r+1$ one of the two terms in the odd recurrence vanishes at $\alpha$, and the other proves the formula by induction.  Multiplying over the three roots of $f$ gives
\[
                 \operatorname{Res}(\psi_m,B)
                     =(-1)^{(m-1)/2}\Delta^{N/2}.
\]
The discriminant of the quadratic $G$ is $\pi^{4a}B(\pi^{-a}X)$.
The discriminant formula for a tower of finite free algebras, or its direct determinant proof after splitting $F$, now gives
\[
 v(\operatorname{disc}\mathcal A)
 =2v(\operatorname{disc}F)+a(4N-3)+(N/2)\delta<D.
\]
This proves the required uniform bound without using a formula for $c_m$ that has not been independently established.

\emph{Why approximate roots are exact enough.}
If $E[m]$ is rational, all $N$ roots of the monic $F$ and both roots of the corresponding monic quadratics $G$ lie in $\OO$.  At each point the Jacobian determinant of $(F,G)$ is $F'G_Y$.  The product of these determinants over all $2N$ points is, up to sign, the discriminant of $\mathcal A$.  All the factors are integral, so each has valuation less than $D$.  Distinct roots $X_i,X_j$ have $v(X_i-X_j)\leq v(\operatorname{disc}F)/2<D$, by the Vandermonde product.  Thus reduction modulo the stated precision supplies all the conditions.

Conversely, lift approximate pairs to $\OO$.  Their equation errors have valuation at least $H\geq2D+1$, whereas their Jacobian determinants have valuation at most $D$.  The following elementary multivariate Hensel argument applies to each pair.  If a polynomial map $f$ has integral Jacobian $J$ at $z$, determinant valuation $b\leq D$, and error valuation at least $H$, substitute $z'=z+\pi^{H-b}w$ and consider
\[
                \pi^{-H}\operatorname{adj}(J)
                              f(z+\pi^{H-b}w).
\]
This is an integral polynomial system.  Its linear Jacobian modulo $\pi$ is $(\det J/\pi^b)I$, which is invertible; all nonlinear coefficients are divisible by $\pi$, since $H-2b\geq1$.  The reduction is therefore a linear system with a unique solution, and successive coefficient lifting gives an exact zero.  Since $J$ is invertible over the fraction field, this is also a zero of the original system.  Its difference from $z$ has valuation at least $H-b\geq D+1$.  The distinctness conditions consequently retain $N$ distinct exact abscissas.  Each exact pair is a nonzero $m$-torsion point; its inverse is also rational.  These $2N$ points exhaust $E[m]\setminus\{0\}$.

The first two conditions give a minimal integral equation with positive discriminant valuation and integral $j$.  It is therefore additive and potentially good, and the torsion conclusion follows.

\emph{Why the fields and the decision procedure are bounded.}
By the Cohen structure theorem \cite[Theorem~10.160.8]{Stacks032A}, a mixed-characteristic complete DVR with perfect residue field contains its Witt coefficient ring.  It is generated over that ring by a uniformizer: reduction modulo $p$ and completeness give the basis $1,\pi,\ldots,\pi^{e-1}$, and the relation is Eisenstein of degree $e$.  Conversely every Eisenstein relation used in the test lifts to such a DVR, and $R$ is its quotient modulo $p^q$.  All conditions above are stable under arbitrary lifts because $eq\geq H$.

For the necessity direction, start with its arbitrary perfect residue field $k_0$ and pass to the completed maximal unramified extension.  The absolute ramification remains $e$, rational torsion remains rational, and the new residue field is an algebraic closure $\overline{k_0}$.  The identification of the absolute Galois group of this completed extension with the original inertia group is stated explicitly in \cite[section~5.1.5, pp.~156--157]{Fontaine1994Semi}.  Serre--Tate's good-reduction criterion \cite[Theorem~1, p.~493]{SerreTate1968Good} therefore shows that a bad curve remains bad.  Its integral $j$ remains integral, so it remains additive and potentially good.  Choose a new minimal integral equation after this extension; no invariance claim about the old equation's minimality is needed.

The field $\overline{k_0}$ need not be algebraic over $\mathbf F_p$.  The passage from that field to the field $k=\overline{\mathbf F}_p$ in the test is instead a logical transfer of a finite formula over $\mathbf F_p$.  Addition and multiplication on $W_q$ are universal polynomial operations in $q$ residue-field coordinates, and reduction by the monic Eisenstein polynomial gives polynomial operations on $R$.  Thus the entire displayed test, including its existential and universal quantifiers, is a finite first-order polynomial formula with coefficients in $\mathbf F_p$.  Elimination in algebraically closed fields makes its truth independent of the algebraically closed field of characteristic $p$ used to evaluate it.  In particular a witness over $\overline{k_0}$ gives an affirmative test over $\overline{\mathbf F}_p$, without identifying these two residue fields.

One can describe that elimination procedure directly.  Replace inequations by inverses, express Boolean formulas as constructible sets, and compute existential images by elimination of polynomial ideals.  On each irreducible component of the closure of an image, a Gr\"obner basis over its function field supplies a nonzero denominator whose complement lies in the image; repeat on the remaining closed subset.  Dimension strictly decreases, so this is finite.  Complements handle universal quantifiers.  All coefficients lie in $\mathbf F_p$, so a nonempty constructible output has a point over $\overline{\mathbf F}_p$ with finitely many coordinates in one finite field.  Lifting their Witt coordinates to that finite unramified extension of $\mathbf Q_p$ and the Eisenstein coefficients to finite Witt polynomials produces the asserted finite local-field example.
\end{proof}

\par\smallskip\noindent\textbf{Exact computational certificate.}
The program \path{code/verify_finite_elliptic_test.py} reconstructs the division polynomials and computes their exact discriminants and two-torsion resultants for nine specializations, including the degree-$364$ polynomial for $m=27$ on $y^2=x^3-x$.  It also checks $2000$ integer-bound instances.  The original run took $74.857496652$ seconds, and the parent's independent rerun took $74.239452977$ seconds; both passed.  The retained output records the large precision bound at $(3,216,3)$ and explicitly records that its elimination test was not evaluated.  These finite checks support the identities and arithmetic used in the proof; they do not substitute for the universal argument. The formerly open index $216$ is instead decided by the independent construction in Theorem~\ref{thm:B-RD26-wild}.

\par\smallskip\noindent\textbf{Scope.}
This bounded test itself has not been evaluated at $(e,n)=(216,3)$. The simplified arithmetic characterization is now proved separately in Theorem~\ref{thm:B-RD26-wild} and Corollary~\ref{cor:B-RD27-exact-three}, without executing this elimination test.  It also does not decide Part~B.4: arbitrary truncated BT data with a finite wild descent action are not thereby supplied with an infinite equivariant BT lift.


\subsection{Why a ramification-nine base change of the full-nine examples does not solve the next gap}

\begin{proposition}\label{prop:B-RD25-basechange-obstruction}
\status{Proved; independent complete proof audit and exact certificate rerun by the parent passed.}
Let $K$ have absolute ramification index $24$, residue characteristic $3$, and perfect residue field.  If an additive elliptic curve $E/K$ has full rational $9$-torsion, then no extension of $K$ of ramification index $9$ makes its full $27$-torsion rational.  More generally, over its minimal good-reduction field after completed unramified base change, the next full division field has ramification degree at least $27$.

This rules out the proposed construction of an $(e,n)=(216,3)$ example by a ramification-nine base change of an $e=24$ full-nine example.  It does not rule out a new curve defined directly over a field of absolute index $216$; such a curve is constructed in Theorem~\ref{thm:B-RD26-wild}.
\end{proposition}
\begin{proof}
Pass first to completed unramified bases.  The proved tame criterion and the wild bound of Proposition~\ref{prop:B-RD20} force the semistability defect to be $12$.  Let $L/K$ be its minimal good-reduction field, so $v_L(3)=288$.  The good special fibre is supersingular.  Write $t$ for an integral formal parameter of the good model.

The universal valuation identity (B.12.4), together with the wild divisibility of the maximum valuation, gives the unique primitive-nine extrema
\[
                            A=21,\qquad B=1.
\]
Explicitly the only possible region is $r=1$, and
$288=6(A+27B)$ with $A\geq9B$, $3\mid A$, $B\geq1$
gives $(A,B)=(21,1)$; region $r=0$ would require
$A+3B=16$ with $3\mid A$, and is impossible.
Choose a basis $P,Q$ of $E[9]$ with $v_L(t(P))=21$ and $v_L(t(Q))=1$.

Formula (B.12.3) gives the minimum nonzero three-torsion valuation $9$.
The points of larger valuation form a line $C\subset E[3]$ of order three: two independent such points would make every nonzero three-torsion point have larger valuation.  The constant-term product for $[3](T)/T$ gives their common valuation $c$:
\[
                         288=6\cdot9+2c,\qquad c=117.
\]
All these torsion points are $L$-rational.

Let $R$ be a point with $3R=Q$, and let $S$ be any point with $3S=P$.
The Weierstrass factorization of $[3](T)$ is a unit times the product of $T-t(U)$ over $U\in E[3]$.  A point of valuation at least $9$ has image under $[3]$ of valuation at least $81$.  A point of smaller positive valuation $v$ has image of valuation $9v$.  Consequently
\[
                      v_L(t(R))=\frac19,\qquad
                      v_L(t(S))=\frac73.
\]
The field $F=L(R)$ has degree at most $9$, because $E[3]$ is rational; its ramification degree is at least $9$ by the valuation of $t(R)$.  Hence it is a totally ramified Galois extension of degree $9$.
The map $\sigma\mapsto\sigma R-R$ identifies
$\Gamma=\operatorname{Gal}(F/L)$ with $E[3]\simeq(\mathbf Z/3)^2$.
Put $z=t(R)$; it is a uniformizer of $F$.
Integral formal addition gives
\[
 v_F(\sigma z-z)=
 \begin{cases}
 1053,&0\ne\sigma R-R\in C,\\
 81,&\sigma R-R\notin C.
 \end{cases}
\]
Here the valuation of a formal difference equals that of the translating torsion point: the coefficient of that point in the difference is a unit.

Let $H\subset\Gamma$ correspond to $C$.  For every $0\ne\sigma\in H$ the lower ramification number is $1052$.  The elementary ramification estimate proved in Lemma~\ref{lem:B-RD20-wild} says
\[
 v_F(\sigma x-x)\geq v_F(x)+1052\qquad(x\in\mathfrak m_F).
\]
If $S$ belonged to $E(F)$, then $v_F(t(S))=21$; a nonzero translation of $S$ by three-torsion would have difference valuation either $81$ or $1053$, both less than $21+1052=1073$.  Thus every element of $H$ fixes $S$.  Set $F_1=F^H$; then $S$ would be $F_1$-rational, and $v_{F_1}(t(S))=7$.

We compute the ramification number of $F_1/L$ directly, without an appeal to a normalization of the upper numbering.  Its uniformizer
\[
                             w=N_{F/F_1}(z)
\]
satisfies $w=z^3+\epsilon$ with $v_F(\epsilon)\geq1055$.
For $\tau\notin H$, put $\delta=\tau z-z$, so $v_F(\delta)=81$.
Since $v_F(3)=2592$,
\[
 v_F\bigl((z+\delta)^3-z^3\bigr)
   =\min\{2592+2+81,\ 2592+1+162,\ 243\}=243.
\]
The correction $\tau\epsilon-\epsilon$ has valuation at least $1055$ and cannot cancel this term.  Hence
$v_{F_1}(\tau w-w)=81$, so the nontrivial automorphisms of $F_1/L$ have lower ramification number $80$.

For an element of valuation prime to $3$, the elementary ramification estimate is an equality: expand in the uniformizer with Teichm\"uller coefficients; the first term has coefficient equal to its exponent modulo $3$, which is nonzero.  Applied to $t(S)$ of valuation $7$, it gives
\[
                         v_{F_1}(\tau t(S)-t(S))=7+80=87.
\]
But $\tau S-S$ is a nonzero three-torsion point, so its formal coordinate has valuation $3\cdot9=27$ or $3\cdot117=351$ in $F_1$.  Neither equals $87$, a contradiction.
It follows that $F$ does not contain full $27$-torsion.  The full division field over $L$ is a $3$-extension, since the action is trivial modulo $9$; it contains $F$, and its degree is therefore at least $27$.

Finally, an extension $K'/K$ of ramification index $9$, followed by completed unramified base change, has degree $9$.  The compositum with $L$ has degree at most $9$ over $L$.  If it contained all $27$-torsion it would contradict the preceding argument.  This proves the asserted obstruction.
\end{proof}

\begin{remark}[Finite arithmetic certificate]\label{auto:B:c83d5f781371}
\status{Proved under stated hypotheses}

The exact certificate \path{code/verify_level27_basechange_obstruction.py}
enumerates every feasible $(A,B,r)$, checks the three-torsion and
next-division valuations, and verifies both ramification contradictions.
Its original run returned PASS in $0.006479513998783659$ seconds;
the parent independently reran it with PASS in $0.033169171001645736$ seconds.
These checks supplement the formal-group and field arguments above.
\end{remark}


\subsection{Wild equivariant arcs and the exact characteristic-three range}

\begin{lemma}[A universal wild trace invariant]
\label{lem:B-RD26-trace-arc}
\status{Proved; independent complete proof and primary-source audit by the parent passed, with an independent exact certificate rerun.}
Let $n\geq1$, and let $A\geq1$ be divisible by three. Put
\[
 d=2\cdot3^{n-1}\left(A+3^n+
        \sum_{j=1}^{n-1}2\cdot3^{n-j-1}\min(A,3^{2j})\right).
                                                        \tag{B-RD26.1}
\]
There is a good supersingular elliptic curve over a complete DVR
$R$ with algebraically closed residue field of characteristic three,
absolute ramification index $d$, and full rational $3^n$-torsion,
whose primitive torsion parameters have extrema $A,1$.
Moreover $R$ and this marked curve admit faithful semilinear actions
of the following groups, fixing every marked torsion point:
\[
 C_3;\qquad C_6\ \text{if }A\text{ is odd};\qquad
 C_3\rtimes C_4\ \text{if }A\equiv1\pmod4.
\]
The action on the special fibre is the usual faithful subgroup of
$\operatorname{Aut}(Y^2=X^3-X)$, and all these actions are trivial
on the residue field. The order-four generator acts on $C_3$ by
inversion.
\end{lemma}

\begin{proof}
Let $k$ be an algebraic closure of $\mathbf F_3$, and work over
$W=W(k)$. The elliptic curve
\[
       \mathcal E_u:\quad Y^2=X^3+uX^2-X
                                                        \tag{B-RD26.2}
\]
over $W[[u]]$ is a universal framed deformation of its special
fibre $E_0$. Here the framing identifies the special fibre with
$E_0$, and isomorphisms respect that identification. To check the
parameter directly, complete the square in any deformation, move
the simple root lifting $X=0$ to zero, and take the unique fourth
root reducing to one to make the coefficient of $X$ equal to $-1$.
The only remaining coefficient is $u$. The simple-root and
fourth-root operations are unique because two and four are units.
This also proves the required uniqueness for coordinate changes
reducing to the identity. Its discriminant is $16(u^2+4)$, a unit.
Serre--Tate identifies this deformation with that of the height-two
formal group; the precise nilpotent-base theorem is
\cite[Theorem~1.2.1]{Katz1981SerreTate}.

Use the odd formal parameter $T=X^{-1/2}$ at infinity, with the
choice of sign for which
\[
 Y=-T^{-3}(1+uT^2-T^4)^{1/2},\qquad
 \omega=\frac{dX}{2Y}=
                  \frac{dT}{(1+uT^2-T^4)^{1/2}}.
\]
The coefficient of $T^3$ in $[3](T)$ is exactly $4u$.
Indeed the formal logarithm begins $T-uT^3/6$, so substitution in
$\log([3](T))=3\log(T)$ gives this coefficient. This identity is
an identity of integral coefficients, although it was calculated
after inverting three.

Let $R_n$ be the full Drinfeld-level deformation ring, and let
$x,y$ be the parameters of a universal basis of $3^n$-torsion.
It is a complete regular local ring with these two regular
parameters, by \cite[Theorem~2.1.2 and Proposition~2.2.3(i)]
{CarterStrauch2021LubinTate}. The universal family in (B-RD26.2)
pulls back to this ring. Changing the special-fibre framing gives
an action on $R_n$ and its family of
\[
 G=\langle\sigma,\tau\rangle=C_3\rtimes C_4,
\]
where on $E_0$ the first generator translates $X$ by one and the
second is the order-four automorphism with $T\mapsto iT$,
$i^2=-1$. This is the functorial action by precomposition of the
framing; its isomorphisms preserve the universal level structure.
In the explicit family, the second action is $u\mapsto-u$ and
$T\mapsto iT$ exactly. The first action reduces to
\[
              h(T)=\frac{T}{(1+T^2)^{1/2}}.
\]
The isomorphisms preserve the origin. Consequently, for $j=0,1,2$,
the difference between the series for $\sigma^j(T)$ and
\[
                 h_j(T)=\frac{T}{(1+jT^2)^{1/2}}
\]
belongs to $(3,u)T W[[u,T]]$. The factor $T$ here is essential:
it follows from preservation of the zero section, in addition to
agreement of the special-fibre series.

We record two bounds in $R_n$ needed to make the forthcoming
rational expressions integral on a chart. Set
\[
               M=3^{2n-2},\qquad H=6M.
\]
Every three-torsion coordinate belongs to
$\mathfrak m_{R_n}^{M}$, by iterating the prime factorization in
\cite[Proposition~2.2.3(vii)]{CarterStrauch2021LubinTate}.
Weierstrass preparation factors $[3](T)$ as a unit times
$T\prod_{0\ne Q\in E[3]}(T-t(Q))$. Its coefficient of $T^3$
is a sum of terms containing at least six of these coordinates.
The formula $[3]_{T^3}=4u$ therefore proves
\[
                    u\in\mathfrak m_{R_n}^{H}.          \tag{B-RD26.3}
\]
Likewise $3\in\mathfrak m_{R_n}^{8M}$. In particular $H\geq6$.
The special formal series is $[3](T)=-T^9$. For an explicit
algebraic check, the three-division polynomial of (B-RD26.2) is
$3X^4+4uX^3-6X^2-1$, and its tripling numerator reduces to $X^9$
modulo $(3,u)$. Substitution $X=T^{-2}$ with the square-root
branch whose leading coefficient is three gives the asserted
series. The sign is also checked directly by the characteristic-three
addition law: it gives $Y([3]P)=-Y^9$, using the polynomial identity
\[
 ((X^3-X)^2-1)(X-X^9)-(X^3-X)+(X^3-X)^5=0
                                      \quad\text{in }\mathbf F_3[X].
\]
Hence the successive coordinates $[3^j]y$ have weighted
initial forms $(-1)^jY^{3^{2j}}$ whenever $j<n$: the error ideal
$(3,u)$ has weight at least $H>M$, while these weights are at
most $M$.

Here is a precise weighted chart construction, including the
possible coefficients divisible by three in a Cohen presentation.
Give $x,y,3$ weights $A,1,d$. The factorization argument in the
proof of Theorem~\ref{thm:B-RD12}, which does not use tame descent
at this stage, gives lifts of all primitive torsion coordinates
with the prescribed weighted initial forms: a nonzero multiple
of $Y$, or the lower-weight term of $X+cY^{3^{2j}}$, or that
binomial at an equality of weights. Choose $c\in k^\times$ so
none of these finitely many initial forms vanishes at $(c,1)$.
In fact one may choose $c=i$ uniformly. All nonzero scalar
coefficients of the boundary binomials are in $\mathbf F_3^\times$
by the preceding explicit series and the linear terms of formal
addition and multiplication by units. The only excluded values
at a boundary are therefore $1,-1$, whereas $i^2=-1$ makes
$i\ne\pm1$. This optional choice fixes $\lambda=-1$ below.
The relation between three and the product of all primitive
coordinates supplies a generator of the Cohen kernel
\[
          R_n\simeq W[[X_0,Y_0]]/(3-P(X_0,Y_0)).        \tag{B-RD26.4}
\]
The series $P$ is a unit times the product of the chosen lifts.
It has weight at least $d$, and its weight-$d$ initial form,
evaluated at $(c,1)$, is a nonzero element $v_0\in k$.
Terms of a lifted factor divisible by three have weight at least
$d$, strictly greater than the weight of that factor. In the
whole product they consequently have weight strictly greater
than $d$. Thus the weight-$d$ part is independent of the variable
used below for $3/\pi^d$. This explicitly deals with low ordinary
degree monomials whose coefficients are divisible by three.
The generator in (B-RD26.4) has coefficient one at its linear
term $3$, so it generates the Cohen kernel, rather than merely
an ideal contained in it.

Adjoin the chart variables
\[
          x=\pi^A Z,\qquad y=\pi,\qquad3=\pi^d V,
                                                        \tag{B-RD26.5}
\]
where $Z,V$ are units, and complete at $Z=c,V=v_0,\pi=0$.
This chart has the two relations
\[
           3=\pi^d V,\qquad
           V=P(\pi^A Z,\pi)/\pi^d.                    \tag{B-RD26.6}
\]
The division in the second relation is legitimate: expand each
coefficient of $P$ in Teichmuller digits and replace each factor
$3^b$ by $\pi^{db}V^b$ using the first relation. The preceding
weighted bound gives a convergent series with nonnegative powers
of $\pi$. Its reduction is a polynomial in $Z$ alone, with value
$v_0$ at $c$. The second equation therefore has derivative one
with respect to $V$ modulo the maximal ideal. Equivalently, this
is the completed chart obtained by adjoining $x/y^A$ and $3/y^d$
to $R_n$ and removing the exceptional torsion. Before imposing
the next equation it is a complete regular local ring of
dimension two, with parameters $\pi,Z-c$. Formula (B-RD26.3)
gives $u\in\pi^H$ on this chart, while $3\in\pi^d$ and
$d\geq8M>H$.

The $C_3$ action already extends to this chart, before any trace
is formed. Indeed $\sigma x/x$ and $\sigma y/y$ are units of
$R_n$, and
\[
 \sigma Z=(\sigma x/x)(\sigma y/y)^{-A}Z,\qquad
 \sigma V=(\sigma y/y)^{-d}V.
\]
Both displayed multipliers reduce to one. The action therefore
preserves the centre $Z=c,V=v_0$ and extends continuously to its
completion. These formulas come from the actual ratios in the
fraction field, so they respect both chart relations.

In its fraction ring put
\[
 z_x=x^{-2},\quad z_y=y^{-2},\quad
 D_y=\operatorname{Tr}_{\langle\sigma\rangle}(z_y^2),
 \quad
 I=\frac{\operatorname{Tr}_{\langle\sigma\rangle}
                      (z_y^2(z_x-z_y))}{D_y},\quad
 U_y=N_{\langle\sigma\rangle}(y)^{-2}.                \tag{B-RD26.7}
\]
These are exact $C_3$-invariants. We verify all divisions and
error estimates, interpreting $O(\pi^b)$ inside the fraction
ring of this chart. The origin-preserving series above give
\[
 (\sigma^j x)^{-2}=x^{-2}+j+O(\pi^{H-2A}),\qquad
 (\sigma^j y)^{-2}=y^{-2}+j+O(\pi^{H-2}).
\]
It follows that
\[
 \begin{split}
 D_y&=3y^{-4}+6y^{-2}+5+O(\pi^{H-4})
                         =-1+O(\pi^{H-4}),\\
 I&=z_x-z_y+O(\pi^{H-2A-4}),\\
 U_y&=y^{-6}-y^{-2}+O(\pi^{H-6}).
 \end{split}                                           \tag{B-RD26.8}
\]
For the middle equality, subtract $(z_x-z_y)D_y$ before
estimating the numerator: each remaining factor has valuation
at least $H-2A-4$. For the last equality use
$\prod_{j=0}^2(y^{-2}+j)=y^{-6}+3y^{-4}+2y^{-2}$;
its difference from $y^{-6}-y^{-2}$ has a factor three.
Since $H>4$, $D_y$ is an actual unit of the chart. These
calculations also show that all the normalized expressions used
next are convergent integral series, not just valuation
estimates after an unspecified specialization.

Put $a=A/3$, and take $\lambda\in W^\times$ with residue $c^{-2}$.
Impose the exact invariant equation $I=\lambda U_y^a$.
After multiplication by $x^2$ this is an integral equation on
the chart, and (B-RD26.8) gives its explicit form
\[
 x^2(I-\lambda U_y^a)
   =1-\pi^{2A-2}Z^2-\lambda Z^2(1-\pi^4)^a
                                  +O(\pi^{H-4}).       \tag{B-RD26.9}
\]
Its reduction is $1-\bar\lambda Z^2$. At $Z=c$ its derivative
is $-2\bar\lambda c\ne0$. The formal implicit-function theorem,
applied together with the second equation of (B-RD26.6), now
eliminates $Z$ and $V$. The resulting ring $R$ is a complete DVR
with uniformizer $\pi$ and $3=\pi^d$ times a unit. This is a
simple-root argument over a complete local ring; no division by
three or averaging by the wild group order was used.
The choice of $c$ guarantees that every primitive coordinate has
exactly its assigned valuation, so the extrema are $A,1$.

Finally the chart and this branch are stable under the asserted
groups. Since every coordinate isomorphism preserves the
origin, $\sigma x/x$ and $\sigma y/y$ are units. They are both
one modulo $\pi$; hence $\sigma$ fixes the residue of $Z=x/y^A$
and of $V=3/y^d$. The rational equation is exactly
$C_3$-invariant, so its unique simple branch is stable.
For $\tau$, $x,y$ have eigenvalue $i$ exactly. Conjugation
permutes the three factors in the trace and norm, giving
\[
             \tau I=-I,\qquad \tau U_y=-U_y.
\]
If $A\equiv1\pmod4$, then $a=A/3$ is odd, the equation is
stable, and $\tau(Z)\equiv i^{1-A}c=c$. Also $4\mid d$ in
this case. For the order-two element $\tau^2$, stability of the
equation is automatic; its branch is stable when $A$ is odd,
and then $2\mid d$. The wild action is faithful on $R$:
\[
                 \sigma\pi=\pi-\tfrac12\pi^3
                                      +O(\pi^5),
\]
because the deformation error has order at least $H+1>3$.
The tame action has the stated faithful eigenvalue on $\pi$.
All actions fix $W$ and the residue field. The generic marked
elliptic curve and its level structure descend over
$\operatorname{Frac}(R)^G$ (or the indicated subgroup's fixed
field). This last assertion concerns the generic fibre; the
integral good model over $R$ is used for its faithful special
inertia action. This proves every assertion of the lemma.
\end{proof}

\begin{remark}[Exact finite certificate]\label{auto:B:10f3741cc5dc}
\status{Proved under stated hypotheses}

The accompanying program \path{code/verify_wild_trace_arcs.py}
checks the trace and norm identities, the special tripling
identity, the logarithmic coefficient $[3]_{T^3}=4u$, the
equivariance of the residual curve, and the exact integer
construction and least-index formula over a finite range.
It checks $66$ symbolic equivariance identities, $109727$ constructed
integer arcs, and $25000$ least-index comparisons through $e=50000$
and $n=8$. The final sign-expanded author run passed in
$2.5154165040003136$ seconds; the independent parent and
obstructions/crystals audits passed in $2.827041984997777$ and
$2.6579352160006238$ seconds respectively. Both independent
audits also checked the complete proof and original sources.
It does not computationally instantiate the universal formal
deformation or its descent; those are proved above.
\end{remark}

\begin{theorem}[B-RD26: the sharp wild higher-torsion criterion]
\label{thm:B-RD26-wild}
\status{Proved; independent parent proof, source, and exact arithmetic audits passed.}
Let $p=3$, $e$ be a positive even integer, $s=v_3(e/2)$, and
$m\in\{3,6,12\}$. An additive elliptic curve with full rational
$3^n$-torsion and exact semistability defect $m$ exists over a
complete DVR of absolute ramification index $e$ and perfect
residue field if and only if
\[
 s\geq n-1,\qquad
 e\geq\frac8m\,3^{\max(2n-2,\,4n-4-s)}.              \tag{B-RD26.10}
\]
The fields constructed here have algebraically closed residue
field. No finite-residue-field assertion is needed for this
existence criterion.
\end{theorem}

\begin{proof}
Necessity is Proposition~\ref{prop:B-RD20}. Suppose the two
inequalities hold. Start at $r_0=\max(0,2n-2-s)$ and increase
$r$ until $r=n-1$ or the inequality
\[
                 me\geq8\cdot3^{2n+r-2}
\]
would fail on increasing $r$ by one. Define
\[
       A=\frac{me}{2\cdot3^{2n-r-2}}-3^{2r+1},
       \qquad B=1.                                    \tag{B-RD26.11}
\]
Then $3\mid A$, $A\geq3^{2r}$, and, unless $r=n-1$,
$A<3^{2r+2}$. The divisibility by three follows from
$v_3(me/2)=s+1$ and $2n-r-2\leq s$.
The geometric-series identity (B-RD20.5) therefore says that
(B-RD26.1) equals $me$. For $m=6$, $A$ is odd; for $m=12$,
$A\equiv1\pmod4$. Apply Lemma~\ref{lem:B-RD26-trace-arc}
and the subgroup of order $m$. Its fixed field has absolute
index $me/m=e$. Galois descent is effective for the elliptic
curve, for instance using the invariant divisor $3[0]$.
The marked torsion points descend and are all rational.
The special-fibre action is faithful of order $m$, so the
good-reduction criterion gives exact semistability defect $m$.
In particular reduction is additive and potentially good.
\end{proof}

\begin{corollary}[B-RD27: the complete higher elliptic range]
\label{cor:B-RD27-exact-three}
\status{Proved; the full wild construction and integer criterion have passed independent audit.}
For $p=3$, write $e=2u$, $s=v_3(u)$ and
$L=\lfloor\log_3u\rfloor$. The largest level admitting an additive
elliptic curve with full rational $3^n$-torsion is
\[
 N(e)=1+\min\left(s,\left\lfloor\frac{L+1}{2}\right\rfloor,
                         \left\lfloor\frac{L+s+1}{4}\right\rfloor\right).
                                                        \tag{B-RD27.1}
\]
Equivalently, existence at level $n$ is exactly
\[
  v_3(e/2)\geq n-1,\qquad
  e\geq\frac23\,3^{\max(2n-2,4n-4-v_3(e/2))}.          \tag{B-RD27.2}
\]
Among potentially good elliptic curves, the least positive full
torsion level forcing good reduction is $N(e)+1$.
In particular full $27$-torsion with additive reduction exists
at $e=216$, with exact defect twelve. Together with the already
proved tame criterion, these formulas decide the elliptic
higher-torsion existence question at every odd prime.
\end{corollary}

\begin{proof}
Defect twelve gives the weakest wild inequality and is always
available. Its realized maximum is the upper bound $N_{\mathrm w}$
already calculated in Corollary~\ref{cor:B-RD21}. The tame maximum
is no larger. The final index statement uses the potentially good
hypothesis; multiplicative Tate examples retain the unrestricted
qualification in Proposition~\ref{prop:B-RD15}.
\end{proof}


\subsection{A sharper wild lower bound for the character problem}

\begin{proposition}[B-RD28: wild character ambiguity in every higher height]
\label{prop:B-RD28-wild-character}
\status{Proved; independent parent proof and indexing audit passed.}
Let $p=3$, let $e$ be a positive even integer, and put
\[
 s=v_3(e/2),\qquad L=\lfloor\log_3(e/2)\rfloor,\qquad
 N_{\mathrm w}=1+\min\left\{s,
     \left\lfloor\frac{L+1}{2}\right\rfloor,
     \left\lfloor\frac{L+s+1}{4}\right\rfloor\right\}.
\]
For every height $h\geq2$ there are two height-$h$ BT groups
$M,N$ over one complete discretely valued field $K$ of absolute
ramification $e$, with perfect residue field and $\mu_3\subset K$,
such that both groups acquire good reduction,
\[
                 M[3^{N_{\mathrm w}}]\simeq
                 N[3^{N_{\mathrm w}}],
\]
and their manuscript characters differ on a prime-to-three
inertia element. Consequently the least universal B.4 index,
if denoted by $\nu_{\min}(3,e,h)$, satisfies
\[
                       \nu_{\min}(3,e,h)\geq N_{\mathrm w}.
                                                        \tag{B-RD28.1}
\]
In particular the lower bounds at $(e,h)=(24,2)$ and $(216,2)$
are respectively two and three.
\end{proposition}

\begin{proof}
Apply Theorem~\ref{thm:B-RD26-wild} with $n=N_{\mathrm w}$ and
$m=12$. The construction gives an elliptic curve $E/K$ with
full rational $3^n$-torsion and faithful special inertia group
$C_3\rtimes C_4$ on its supersingular good reduction
$Y^2=X^3-X$. By the Weil pairing, $\mu_{3^n}\subset K$.
Set
\[
 M=E[3^\infty]\times(\mathbf Q_3/\mathbf Z_3)^{h-2},
 \qquad N=(\mathbf Q_3/\mathbf Z_3)^h.
\]
These generic BT groups have isomorphic constant $3^n$-torsion.
The second group has good reduction already over $\mathcal O_K$;
the first acquires good reduction over the degree-twelve field
of the construction. Thus both satisfy the manuscript monodromy
condition with a single good-after-extension graded piece.

For the order-four automorphism $\tau$ of the supersingular
elliptic reduction, $\tau^2=[-1]$. Its two eigenvalues on the
rational covariant Dieudonne module are $i,-i$; one may see the
opposite pair either from the degree-one determinant or from
Frobenius interchanging the two primitive fourth roots in
characteristic three. Hence its trace is zero. On $M$ the trace
is $h-2$, whereas on $N$ it is $h$. The period twist in the
corrected Fontaine dictionary has trivial finite inertia action,
so these are exactly the manuscript character values.
The characters therefore differ. Failure at
$3^n=M(n-1)$ gives $\nu_{\min}\geq n$, with precisely the indexing
specified in the manuscript.
\end{proof}

\begin{remark}[Scope of the lower bound]\label{auto:B:688dec0e6255}
\status{Proved under stated hypotheses}

This improves the tame Honda bound in
Corollary~\ref{cor:B-RD18} in the wild exceptional ranges.
It does not assert that the displayed lower bound is the least
universal index: the latter concerns arbitrary pairs of groups
with isomorphic truncations, and the comparison group need not
be constant. The independent characteristic-two example in
Proposition~\ref{prop:B-RD22} already illustrates that distinction.
The finite upper bound in Theorem~\ref{thm:B-RD17} continues to
apply to every pair satisfying the full filtration condition.
The numerical formula here is included among the $25000$ exact
least-index comparisons in the certificate for
Theorem~\ref{thm:B-RD26-wild}; no new unexecuted computation is
being used.
\end{remark}


\subsection{A level-nine character ambiguity at the smallest odd ramification}

\begin{lemma}[Cyclic isogenies and a Tate comparison]
\label{lem:B-RD29-isogeny-Tate}
\status{Proved under stated hypotheses}
Let $K$ be a complete discretely valued field of characteristic zero,
let $r\geq2$, and let $E/K$ have a $K$-rational point $P$ of exact
order $r$. Put $B=E/\langle P\rangle$. There is a split Tate curve
$T_q/K$, with a parameter $q\in K^\times$ of positive valuation,
such that $B[r]\simeq T_q[r]$ as finite group schemes over $K$.
The generic $p$-divisible group of a Tate curve satisfies the
manuscript monodromy condition: its two graded pieces are
$\mu_{p^\infty}$ and $\mathbf Q_p/\mathbf Z_p$.
\end{lemma}
\begin{proof}
The dual of the quotient isogeny is a morphism
$\widehat\phi:B\to E$ with kernel canonically Cartier dual to
$\langle P\rangle\simeq\mathbf Z/r$. Its restriction gives an
exact sequence of finite etale $K$-group schemes
\[
 0\longrightarrow\mu_r\longrightarrow B[r]
   \xrightarrow{\widehat\phi}\langle P\rangle
   \longrightarrow0.                                      \tag{B-RD29.1}
\]
It is exact on the right: an inverse image of $P$ under
$\widehat\phi$ is killed by $r$, since
$[r]=\phi\widehat\phi$ on $B$ and $\phi(P)=0$.
Identify the last group with $\mathbf Z/r$ using its generator $P$.
The sequence is killed by $r$, so it splits as a sequence of
abstract $\mathbf Z/r$-modules. Its Galois descent class lies in
\[
 H^1(K,\mu_r)=K^\times/(K^\times)^r.
\]
Choose a representative $q_0$ and multiply it by an $r$th power
of a sufficiently high power of a uniformizer to obtain $v_K(q)>0$.
The finite group associated to this Kummer class is exactly
$T_q[r]$. This is the explicit finite group scheme
$G_{K,q,r}$ in \cite[``Points of finite order'', pp.18--19]{Tate1995Uniformization};
its component over $a\in\mathbf Z/r$ is given by $z^r=q^a$.
That source proves both its extension-class description and its
identification with the torsion of the Tate curve. The latter
curve exists by its Theorem~1, p.6.

For $r=p^n$ these descriptions at every level yield
\[
 0\longrightarrow\mu_{p^\infty}\longrightarrow T_q[p^\infty]
   \longrightarrow\mathbf Q_p/\mathbf Z_p\longrightarrow0.
\]
Both graded groups have good models over $\mathcal O_K$.
Their finite inertia characters are trivial, so the manuscript
character of $T_q[p^\infty]$ is constantly two on inertia.
\end{proof}

\begin{proposition}[B-RD29: nine-torsion fails already at $p=3,e=2$]
\label{prop:B-RD29-p3-level-two}
\status{Proved; independent parent complete proof, original-source and exact-certificate audits passed.}
Let $K=\mathbf Q_3(\zeta_3)$, whose absolute ramification is two
and whose residue field is $\mathbf F_3$. Define
\[
 B:\quad y^2-xy+y=x^3-x^2-28x-53,
 \qquad q=2^5\,3^6=23328,
\]
and let $T_q/K$ be the split Tate curve with this parameter.
Then
\[
 B[9]\simeq T_q[9],                                      \tag{B-RD29.2}
\]
but the manuscript characters of $B[3^\infty]$ and
$T_q[3^\infty]$ differ on an inertia element of order two in
the finite inertia image. The first group is potentially good;
the second satisfies the filtration condition with good rank-one
graded pieces. Consequently
\[
                 \nu_{\min}(3,2,h)\geq2\qquad(h\geq2).
                                                        \tag{B-RD29.3}
\]
The same lower bound holds for every $e\equiv2\pmod4$.
\end{proposition}

\begin{proof}
We give the quotient, the Kummer class and the inertia calculation
explicitly. Start with
\[
 E:\quad y^2-3xy-12y=x^3-12x^2,
 \qquad P=(0,0).
\]
Exact addition gives
\[
 \begin{array}{c|rrrrrrrr}
 a&1&2&3&4&5&6&7&8\\\hline
 x(aP)&0&12&4&6&6&4&12&0\\
 y(aP)&0&48&8&18&12&16&0&12
 \end{array}
\]
and $9P=0$. In particular $P$ has exact order nine.
The quotient by this subgroup is
\[
 E':\quad y^2-3xy-12y=x^3-12x^2-2310x-36190.
                                                        \tag{B-RD29.4}
\]
Here is a direct verification of the quotient formula, so that
no equation is imported without checking its isogeny. In short
coordinates $X=x-13/4$, $Y=y-(3x+12)/2$, put
\[
 A=-219/16,\quad C=827/32,\quad F(X)=X^3+AX+C,
 \quad\mathcal R=\{-13/4,35/4,3/4,11/4\}.
\]
Define
\[
 R(X)=X+\sum_{a\in\mathcal R}
   \left(\frac{2(3a^2+A)}{X-a}
              +\frac{4F(a)}{(X-a)^2}\right).
\]
Exact polynomial arithmetic gives
\[
 F(X)R'(X)^2=R(X)^3-\frac{37179}{16}R(X)
                         -\frac{1397493}{32}.
                                                        \tag{B-RD29.5}
\]
The reduced numerator and denominator of $R$ have degrees nine
and eight, and the denominator has double zeros precisely at
$\mathcal R$. Thus $(X,Y)\mapsto(R(X),YR'(X))$ extends to a
pointed degree-nine morphism of smooth elliptic curves, hence an
isogeny, and its kernel consists of the nine displayed points.
The target short equation is exactly (B-RD29.4) in the original
long coordinates. Finally $x=9x_1+1$, $y=27y_1+21$ changes
(B-RD29.4) to the displayed equation for $B$. This identifies
$B$ with $E/\langle P\rangle$ over $\mathbf Q$.

To compute the extension class in (B-RD29.1), consider
\[
 f(x,y)=\frac{(x^2-y)^3(y-3x+4)}{(x-4)^3}.
                                                        \tag{B-RD29.6}
\]
The tangent and chord construction gives
\[
 \operatorname{div}\left(\frac{x^2-y}{x-4}\right)
       =3[P]-[3P]-2[0],
 \qquad
 \operatorname{div}(y-3x+4)=3[3P]-3[0].
\]
For the first equality, the usual two addition lines are $y=0$
and $y=4x$, and the equation of $E$ gives
$y(y-4x)=(x-12)(x^2-y)$. For the second, substitution of
$y=3x-4$ into the equation gives $-(x-4)^3$.
Consequently $\operatorname{div}(f)=9[P]-9[0]$.

Let $\delta:E(K)\to K^\times/(K^\times)^9$ be the connecting
homomorphism of the dual isogeny. The precise divisor description
of this map is \cite[Appendix~A.1, Lemmas~A.1--A.2, p.22]{AntieauAuel2021Descent}.
There is a constant $c$ such that $cf$ is normalized for that
map. Since $5P$ and $4P$ avoid both its zero and its pole,
homomorphicity cancels this constant and gives
\[
 \delta(P)=\frac{f(5P)}{f(4P)}
       =\frac{-3456}{2916}
       =-\frac{32}{27}
       \equiv2^5\,3^6\pmod{(K^\times)^9}.               \tag{B-RD29.7}
\]
The last ratio is the ninth power $(-3)^9$. One can also compare
this calculation directly with the $N=9$ family and Figure~9 on
pp.26--27 of that source, specialized at $\lambda=1$.
Thus (B-RD29.1) is the Kummer extension with the specified
$q=23328$, whose valuation over $K$ is twelve. The preceding
lemma proves (B-RD29.2) with this explicit Tate parameter.

The exact invariants of $B$ are
\[
 c_4=17\,3^4,\qquad \Delta=-2\,3^9,
 \qquad j=-\frac{17^3\,3^3}{2}.                          \tag{B-RD29.8}
\]
The integral $j$ gives potential good reduction by
\cite[Section~2, p.496]{SerreTate1968Good}.
Over the completed maximal unramified extension of $K$, take
an extension of minimal good-reduction degree $m$. Its inertia
acts faithfully on the good special elliptic curve.
The displayed equation has discriminant valuation eighteen.
A change to a unit-discriminant good equation over the extension
therefore implies
\[
                            12\mid18m,
\]
so $m$ is even. The inertia image contains an element of order
two. In residue characteristic three the only nontrivial
order-two origin-preserving elliptic automorphism is $[-1]$;
it acts as $-1$ on the rational Dieudonne module and has trace
$-2$. The Tate comparison has trace $2$, by its two good graded
pieces. Hence the requested characters differ.
The inertia calculation uses the actual finite descent action
of the good model and therefore has the conventions of
Proposition~\ref{prop:B-RD4-dictionary}.

For $e\equiv2\pmod4$, take an extension of $K$ generated by a
root of $X^{e/2}-(\zeta_3-1)$. This polynomial is Eisenstein,
so its absolute ramification is $e$ and residue field remains
$\mathbf F_3$. The isomorphism (B-RD29.2) survives base change.
The raw discriminant valuation becomes $9e$, and
$12\mid9em$ still forces $m$ even. Thus the same trace
comparison holds. Adding $h-2$ identical constant etale factors
to both groups gives every $h\geq2$. Finally level nine is
$M(1)$, so its failure proves $\nu_{\min}\geq2$.
\end{proof}

\begin{remark}[Exact certificate and remaining question]\label{auto:B:7f506124b28c}
\status{Proved under stated hypotheses}

The program \path{code/verify_p3_character_level2.py} verifies all
nine subgroup points and all $81$ additions, the degree-nine
isogeny identity, the change to $B$, both discriminants and
$j$-invariants, the two divisor identities, the Kummer ratio and
its ninth-power adjustment. Its recorded exact run passed in
$4.882327191000513$ seconds; the independent parent rerun passed
in $4.293367760001274$ seconds. The program verifies finite
algebraic assertions; the Tate extension dictionary and inertia
argument are proved above.
This proves that the full-torsion bound $N_{\mathrm w}=1$ at
$e=2$ is strictly smaller than the least B.4 index. Whether
$\nu=2$ already suffices for arbitrary height-two groups at
$(p,e)=(3,2)$ is settled affirmatively in Theorem~\ref{thm:B-RD42}.
\end{remark}


\subsection{A logarithmic universal upper bound}

\begin{proposition}[B-RD30: the sharper generic-Hom exponent]
\label{prop:B-RD30}
\status{Proved; independent proof and primary-source audit passed.}
Use the integer $D_{p,h}$ of Lemma~\ref{lem:B-RD17-degree} and put
\[
 c_{p,e,h}=\left\lfloor\log_p\left(\frac{peD_{p,h}}{p-1}\right)
                  \right\rfloor.
\]
Then a sufficient universal index in B.4 is
\[
                    \nu=2h^2c_{p,e,h}.                 \tag{B-RD30.1}
\]
In particular, the previously established quadratic logarithmic
bound in Theorem~\ref{thm:B-RD17} can be replaced by this bound.
This is an upper bound, not a formula for the least index.
\end{proposition}
\begin{proof}
Bondarko's original Theorems~A and~B
\cite[p.661]{Bondarko2006Generic} assert the following, for every
prime, over a complete discrete valuation ring of absolute
ramification $a$: every homomorphism between generic fibres of
finite flat commutative group schemes extends after multiplication
by $p^c$, where
\[
                  c=\lfloor\log_p(pa/(p-1))\rfloor.
\]
These theorems apply to arbitrary finite flat models, not just
truncated BT groups. The proof is in \S4, pp.681--684; the passage
to arbitrary residue fields is stated on pp.662 and 681. The
stronger value of this exponent is also explicitly confirmed in
\cite[p.2 of the author version]{VasiuZink2012Boundedness}.

Choose the common good-reduction field $L$ from
Lemma~\ref{lem:B-RD17-degree}. Then $e_L\leq eD_{p,h}$, so the
exponent $c$ in Lemma~\ref{lem:B-RD17-common} can be taken at most
$c_{p,e,h}$. Apply that lemma at
$n=2h^2c_{p,e,h}+1$. Its Zassenhaus pieces are finite flat and
therefore satisfy exactly the hypotheses just stated. Finally,
$M(\nu)=M[p^{\nu+1}]$ gives (B-RD30.1).
\end{proof}

\begin{remark}[A useful stronger comparison for actual formal BT groups]\label{auto:B:cb3f24f10664}
\status{Proved under stated hypotheses}

Bondarko's Theorem~E(1), pp.663--664, proved from Proposition~18
on p.685 of \cite{Bondarko2006Generic}, has a different conclusion:
if two finite-height formal groups over $\mathcal O_L$ have
isomorphic generic $p^{\ell+u}$-torsion and
$e_L\leq(p-1)p^u$, then their $p^\ell$-torsion is isomorphic over
$\mathcal O_L$, compatibly with the original generic
identification. Uniqueness of a generic-fibre extension makes this
comparison equivariant for any given descent actions. This
stronger statement will be used only for actual formal BT groups.
It is not being applied to the possibly non-BT Zassenhaus pieces.
\end{remark}


\subsection{The smallest odd-prime case reduces to two levels}

\begin{lemma}[Finite inertia in a height-two formal group at $3$]
\label{lem:B-RD31-finite}
\status{Proved under stated hypotheses}
Let $G/K$ have potentially good reduction to a one-dimensional
formal BT group of height two, with residue characteristic $3$.
After passage to the completed maximal unramified extension, it
has good reduction over a Galois extension of degree at most $12$.
Its faithful special inertia image is either cyclic of order
dividing $8$, or has a normal subgroup of order $3$ and order at
most $12$; in the latter case its tame quotient has order $1,2$,
or $4$. Two such groups have simultaneous good reduction over an
extension of ramification degree at most $36$.
\end{lemma}
\begin{proof}
We work with algebraically closed residue field. Lazard's
classification theorem \cite[Theorem~IV, pp.270--271]{Lazard1955Formal}
identifies the special formal group with the height-two Honda
group. Its rational Dieudonne module has a basis in which
\[
 F=\begin{pmatrix}0&3\\1&0\end{pmatrix}\sigma.
\]
This is also obtained directly from the Dieudonne module of the
supersingular elliptic curve $Y^2=X^3-X$; a change of basis over
$W(\bar k)[1/3]$ absorbs the sign in $F^2=-3$.
The endomorphisms commuting with $F$ are precisely
\[
 A(a,c)=\begin{pmatrix}a&3\sigma(c)\\c&\sigma(a)\end{pmatrix},
       \qquad a,c\in\mathbf Q_{3^2}.
                                                        \tag{B-RD31.1}
\]
Here $\mathbf Q_{3^2}=W(\mathbf F_9)[1/3]$ and $\sigma$ is its
nontrivial automorphism. The determinant is
$N(a)-3N(c)$, whose two nonzero terms have different parity of
$3$-adic valuation. Thus this algebra is a quaternion division
algebra. Its valuation ring consists of the matrices with
$a,c\in W(\mathbf F_9)$; its residue field is $\mathbf F_9$, and
the kernel of reduction on units is pro-$3$. These facts also
follow by writing the algebra as
$\mathbf Q_{3^2}\langle\Pi\rangle$, with
$\Pi^2=3$ and $\Pi a=\sigma(a)\Pi$.

Let $H$ be any finite subgroup of its units. Its reduction has
cyclic image in $\mathbf F_9^\times$. Its kernel $P$ is a finite
$3$-group. If $P\ne1$, choose a central element $z$ of order $3$.
Its centralizer is the quadratic field $\mathbf Q_3(z)$, so all
of $P$ lies in that field. Its finite $3$-power roots of unity
form $\mu_3$: an element of order $9$ would have minimal
polynomial of degree $6$, whereas every element of the
quaternion algebra satisfies a quadratic polynomial. Hence
$P=C_3$. The centralizer of $P$ in $H$ lies in the roots of unity
of $\mathbf Q_3(\zeta_3)$, namely $\mu_6$. Conjugation embeds
$H/C_H(P)$ in $\operatorname{Aut}(C_3)=C_2$. Consequently
$|H|\leq12$. Its tame quotient has order dividing $8$ and at
most $4$, as asserted. If $P=1$, reduction embeds $H$ in the
cyclic group of order $8$.

The faithful special inertia image gives a good-reduction
extension of that degree by exactly the descent and lattice
argument of Lemma~\ref{lem:B-RD17-degree}. In particular, the
good model remains a one-dimensional formal group, since its
special isocrystal remains of slope $1/2$.

For two groups, their joint wild image embeds in $P_1\times P_2$,
and the joint tame quotient is cyclic. If both wild images are
nontrivial, its order is at most $9\cdot4=36$. If only one is
nontrivial, it is at most $3\cdot8=24$. If neither is nontrivial,
it is at most $8$. Applying the same descent argument to the
joint image proves the last assertion.
\end{proof}

\begin{lemma}[A point bound for supersingular formal groups]
\label{lem:B-RD31-point}
\status{Proved under stated hypotheses}
Let $F$ be a one-dimensional formal group over a complete DVR
of mixed characteristic $(0,3)$, of absolute ramification $a$,
whose special fibre has height at least two. If $F$ has a
rational point of exact order $3^n$, then
\[
                         a\geq3^n.                     \tag{B-RD31.2}
\]
\end{lemma}
\begin{proof}
Normalize the valuation to give a uniformizer value $1$.
Write $[3](T)=3T+\sum_{j\geq2}b_jT^j$. For an invariant
differential $u(T)dT$, with $u(T)$ an integral unit series, the
identity $u([3](T))[3]'(T)=3u(T)$ shows that $j b_j$ is divisible by $3$.
Thus $b_2\in3\mathcal O$. Since the special fibre has height
at least two, $b_3,\ldots,b_8$ lie in the maximal ideal.
For $x$ of positive integral valuation $v$ this gives
\[
 v([3]x)\geq\min\{a+v,\,1+3v,\,9v\}.
                                                        \tag{B-RD31.3}
\]
A nonzero point $y$ of order $3$ satisfies
$v(y)\leq(a-1)/2$. Indeed, if $v(y)>(a-1)/2$, the linear term
$3y$ has strictly smaller valuation than all terms of degree
$2$ through $8$. It also has smaller valuation than all terms
of degree at least $9$: if $a\geq2$, then
$8v(y)>4(a-1)\geq a$, with strict inequality also at $a=2$;
if $a=1$, the integral inequality $v(y)\geq1$ gives it
directly. This would contradict $[3]y=0$.

For a point of exact order $3^n$, write
$v_i=v([3]^i x)$ for $0\leq i<n$. The final point satisfies
$v_{n-1}\leq(a-1)/2$. Each earlier $v_i$ also lies below this
bound: (B-RD31.3) strictly increases every positive valuation,
so a value already above that bound could never reach the
final one. Consequently $a\geq2v_i+1$, and
(B-RD31.3) gives $v_{i+1}\geq3v_i+1$.
Starting from $v_0\geq1$, induction yields
\[
 v_{n-1}\geq\frac{3^n-1}{2}.
\]
Combining the two inequalities gives $a\geq3^n$.
\end{proof}

\begin{proposition}[The mixed formal and filtered case is detected by $27$-torsion]
\label{prop:B-RD31-mixed}
\status{Proved under stated hypotheses}
Suppose $p=3$, $e_K=2$, and $\mu_3\subset K$.
Let $M$ be a height-two potentially good group with supersingular
special fibre. Let $N$ admit a filtration with height-one graded
pieces and with sum of their dimensions equal to $1$.
Then
\[
                         M[27]\not\simeq N[27].
                                                        \tag{B-RD31.4}
\]
\end{proposition}
\begin{proof}
Pass to the completed maximal unramified extension, which
preserves $e_K=2$ and the given torsion comparison. A potentially
good height-one character is then $\chi_3^b\epsilon$, where
$b\in\{0,1\}$ and $\epsilon$ has values in $\{1,-1\}$.
This is the height-one classification used in
Proposition~\ref{prop:B-RD5}. The determinant of $T_3M$ is
$\chi_3\epsilon_0$, with $\epsilon_0$ likewise quadratic;
this follows by taking the determinant in crystalline comparison,
as in Proposition~\ref{prop:B-RD18}.

An asserted isomorphism transports the first height-one
subgroup of $N[27]$ to a cyclic direct summand $C\subset M[27]$
whose character is $\chi_3^b\epsilon$. If $b=0$, twist $M$ by
$\epsilon$ to obtain a rational point of order $27$. If $b=1$,
form the BT quotient $M'=M/C$. The image of $M[27]$ in $M'[27]$
is cyclic of order $27$ and has character
$\epsilon_0\epsilon^{-1}$. Twisting $M'$ by this quadratic
character again gives a rational point of exact order $27$.
The quotient and twist remain potentially good, one-dimensional,
and supersingular. A finite flat closure of the kernel in a good
model justifies the assertion for the quotient; the twist becomes
trivial over a quadratic extension.

By Lemma~\ref{lem:B-RD31-finite}, the resulting group acquires
good formal reduction over an extension of ramification degree
at most $12$. Its absolute ramification is therefore at most
$24$. Its rational point extends to a point of this formal
group: finite flat torsion is proper over the DVR, and its
connected special fibre has only the identity geometric point.
Lemma~\ref{lem:B-RD31-point} would require absolute ramification
at least $27$, a contradiction.
\end{proof}

\begin{theorem}[B-RD31: three possibilities in the smallest odd-prime case]
\label{thm:B-RD31}
\status{Proved; two independent complete proof audits and primary-source checks passed.}
In B.4 one has
\[
                         2\leq\nu_{\min}(3,2,2)\leq4.
                                                        \tag{B-RD31.5}
\]
More precisely, $M[243]$ determines the prime-to-$3$ character
of every height-two group satisfying the manuscript monodromy
condition. Any counterexample to determination by $M[27]$ or
$M[81]$ must involve two groups which themselves have
potentially good supersingular reduction. Thus deciding the
two levels $27$ and $81$ in that restricted case decides
$\nu_{\min}(3,2,2)$.
\end{theorem}
\begin{proof}
The lower bound is Proposition~\ref{prop:B-RD29-p3-level-two}.
For the upper bound, let $M[243]\simeq N[243]$ and work after
completed maximal unramified base change. For a height-two
monodromy filtration, let $d\in\{0,1,2\}$ be the sum of the
dimensions of its good graded models. Its determinant character
on inertia is $\chi_3^d\epsilon$, with $\epsilon$ finite
quadratic. The image of $\chi_3$ modulo $9$ has order $3$:
$\mu_3\subset K$, whereas $e_K=2$ excludes $\mu_9\subset K$.
Equality of determinants modulo $9$ therefore forces equal $d$
and equal $\epsilon$ for $M$ and $N$.

If $d=0$, all good graded models are etale. Their special
inertia character is the Brauer lift of the semisimplification
of the generic representation modulo $3$, evaluated on
prime-to-$3$ elements. It is therefore determined already by
$M[3]$. If $d=2$, Cartier duality reduces to $d=0$; the
cyclotomic twist has trivial finite descent inertia. This
also treats extensions between the graded pieces, since
semisimplification and the character are additive.

Suppose $d=1$. Each group either has a filtration by height-one
pieces, or is itself potentially good and supersingular.
Indeed a potentially ordinary height-two graded piece has a
canonical connected--etale filtration on its good model;
uniqueness and descent of this filtration produce the claimed
height-one filtration over $K$.
For two groups with height-one filtrations, the graded generic
characters are $\chi_3^{b_i}\epsilon_i$. Since
$\chi_3\equiv1\pmod3$ on inertia, their finite characters
$\epsilon_i$ and hence their sum are determined by the
semisimplified $3$-torsion. A pair of different kinds is excluded
already at level $27$ by Proposition~\ref{prop:B-RD31-mixed}.

It remains to compare two potentially supersingular groups.
Lemma~\ref{lem:B-RD31-finite} gives a common good formal
extension $L$ with $e_L\leq72$. Since
\[
                         72\leq2\cdot3^4,
\]
Bondarko's Theorem~E(1), with $u=4$ and $\ell=1$, extends the
generic $3^5$-torsion comparison to an isomorphism between the
special $3$-torsion models, equivariantly for descent. The
special Dieudonne modules modulo $3$ consequently have the
same prime-to-$3$ Brauer character. A finite-order matrix of
order prime to $3$ on a free Witt module has eigenvalues which
are Teichmuller lifts of its residual eigenvalues. Thus this is
the required equality of characteristic-zero inertia
characters. The same reductions, without the final application
of Theorem~E, give the assertion about possible counterexamples
at levels $27$ and $81$.
\end{proof}


\begin{remark}[Exact auxiliary certificate and remaining gate]\label{auto:B:a26975327b26}
\status{Proved under stated hypotheses}

The program \path{code/verify_height_two_small_bounds.py} checks the
quaternion matrix identities, all $49$ joint-image numerical cases,
$16000$ formal-valuation inequalities and the $36$ determinant
comparisons used above. It passed in $0.05041721599991433$ seconds;
its JSON output is retained. These finite checks certify the
explicit algebra and inequalities, not the existence or
nonexistence of a counterexample at either remaining level.
At this intermediate stage neither $27$ nor $81$ was decided for arbitrary pairs
of potentially supersingular groups over fields of absolute
ramification $2$.
\end{remark}

\par\smallskip\noindent\textbf{Subsequent exact determination.}
The remaining small-height alternatives isolated here are settled in
Theorem~\ref{thm:B-RD42}: the exact manuscript index is
$\nu_{\min}(3,2,2)=2$. The intermediate bounds and conditional
reductions above are retained as steps in that proof.


\subsection{The remaining $81$-torsion comparison has only two inertia pairs}

Theorem~\ref{thm:B-RD34} below excludes both of the level-$81$
pairs isolated here and strengthens the level-$27$ restriction to
joint inertia order $18$ or $36$. The present proposition records
the group-theoretic reduction used in that proof.

\begin{lemma}[Norm and prime-to-$3$ characters in the quaternion algebra]
\label{lem:B-RD32-norm}
\status{Proved under stated hypotheses}
For the finite subgroups $H$ in Lemma~\ref{lem:B-RD31-finite},
the reduced norm is trivial whenever $3\mid |H|$.
If it is nontrivial, $H$ is cyclic of order $8$.
If it is trivial, the possible orders are
\[
                           1,2,3,4,6,12.
                                                        \tag{B-RD32.1}
\]
In the norm-one case the prime-to-$3$ character depends only on
the tame quotient order $t\in\{1,2,4\}$, and on a prime-to-$3$
element whose image in that quotient has order $r$ it is
\[
       r=1:\ 2,\qquad r=2:\ -2,\qquad r=4:\ 0.
                                                        \tag{B-RD32.2}
\]
Thus two norm-one special inertia images with the same tame
quotient order have identical prime-to-$3$ characters as
functions on inertia.
\end{lemma}
\begin{proof}
Retain the explicit quaternion algebra (B-RD31.1). A noncentral
element has characteristic polynomial of degree two, whose
constant term is its reduced norm. There is only one element
of order two, namely $-1$, because the algebra is a division
algebra and $(x-1)(x+1)=0$ forces $x=1$ or $-1$.

If the wild subgroup is nontrivial, it is $P=C_3$. Its
centralizer in $H$ is contained in
$\mu_6\subset\mathbf Q_3(\zeta_3)$, on which the reduced norm
is $1$. If $H$ is not equal to this centralizer, its quotient
by the centralizer has order two. The tame quotient of order
two cannot give a nonabelian group of order six, since its
involutions would be noncentral. Therefore the nonabelian
case has order twelve. Its Sylow $2$-subgroup injects under
reduction into $\mathbf F_9^\times$, so it is cyclic of order
four. A generator $j$ satisfies $j^2=-1$ and has reduced
characteristic polynomial $X^2+1$, hence reduced norm $1$.
The centralizer and $j$ generate $H$. This proves norm one
for every wild $H$.

If $H$ is tame, reduction embeds it in $\mathbf F_9^\times$,
so it is cyclic of order dividing eight. A generator lies in
a quadratic commutative subfield. For an element of order
eight this is the unramified quadratic field, and its reduced
norm is $\zeta_8^{1+3}=\zeta_8^4=-1$. For elements of order
dividing four the norm is $1$. This proves (B-RD32.1) and
the assertion about nontrivial norm.

In a norm-one subgroup, a prime-to-$3$ element has order
$1,2$, or $4$. Its trace is respectively $2,-2$, or $0$,
the last two facts following from $-1$ and $X^2+1$.
Its cyclic subgroup maps injectively to the tame quotient,
so its order equals the order of its image there.
Over the completed maximal unramified base, the absolute
tame inertia is procyclic and has a unique quotient of each
of the orders $1,2,4$. The order of the image of an element
in that quotient is independent of the choice of a generator.
Formula (B-RD32.2) therefore proves equality as functions on
inertia, even when the two wild good-reduction fields differ.
\end{proof}

\begin{proposition}[B-RD32: exact restriction of the two undecided tests]
\label{prop:B-RD32}
\status{Proved; parent and independent second proof audits passed.}
Keep $p=3$, $e_K=2$, $\mu_3\subset K$, and height two.
For a pair with isomorphic $27$-torsion and different
prime-to-$3$ characters, both groups must be potentially
supersingular, both finite special inertia images must have
trivial reduced norm, their tame quotient orders must differ,
and their joint faithful good-reduction image must have order
strictly greater than $9$.

For a pair with isomorphic $81$-torsion and different
prime-to-$3$ characters, after interchanging the pair their
individual faithful special inertia orders must be
\[
                             (12,3)\quad\hbox{or}\quad(12,6),
                                                        \tag{B-RD32.3}
\]
and their joint faithful special inertia image must have
order exactly $36$. In particular its wild subgroup has
order $9$ and the two wild order-three quotients are
independent.

Consequently $\nu_{\min}(3,2,2)\leq3$ if and only if
no $81$-torsion comparison of the two specified kinds has
different prime-to-$3$ characters. This is a restricted
remaining test, not a claim that either kind exists.
\end{proposition}
\begin{proof}
Theorem~\ref{thm:B-RD31} reduces either test to two whole
potentially supersingular groups. Its determinant argument
shows that their finite determinant characters agree.
Under crystalline comparison this finite determinant is
the reduced norm of the special quaternionic inertia
action: the comparison duality or period Tate twist does
not change finite descent inertia.

If this determinant is nontrivial, Lemma~\ref{lem:B-RD32-norm}
makes both special images cyclic of order eight and tame.
Their common good extension then has ramification at most
eight, hence absolute ramification at most sixteen.
Bondarko's Theorem~E(1), with $u=2$ and $\ell=1$, applies
since $16\leq2\cdot3^2$. Already an isomorphism on
$27$-torsion makes their special $3$-torsion equivariantly
isomorphic and their prime-to-$3$ characters equal.
Thus any counterexample has norm one on both sides.

Their tame quotient orders must differ by
Lemma~\ref{lem:B-RD32-norm}. Write $m$ for the order of
the joint faithful image. Its wild subgroup has order
$1,3$, or $9$, and its tame quotient has order $1,2$, or
$4$. In particular
\[
                m\in\{1,2,3,4,6,9,12,18,36\}.
                                                        \tag{B-RD32.4}
\]
For $m\leq9$, the common absolute ramification is at most
$18=2\cdot3^2$, so the same $u=2$ argument settles the
$27$-torsion comparison. This proves the first assertion.

At level $81$, take $u=3$ and $\ell=1$. The comparison is
settled whenever $2m\leq2\cdot3^3$, or $m\leq27$.
The only larger value in (B-RD32.4) is $36$. This forces
both wild images to have order three and the joint tame
quotient to have order four. One individual image must
therefore have order twelve. If both had that order,
their tame quotient orders and hence their prime-to-$3$
characters would agree. The other image instead has
tame quotient order one or two, and hence total order
three or six. Joint wild order nine is exactly the
asserted independence. Every failure of the unrestricted
$81$ test has now been reduced to (B-RD32.3); the converse
implication about the least index is immediate because
these specified pairs are themselves allowed in B.4.
\end{proof}

\begin{remark}[Exact auxiliary certificate]\label{auto:B:72e408a14989}
\status{Proved under stated hypotheses}

The retained program
\path{code/verify_height_two_character_gate.py}
checks the eight cyclotomic reduced norms and traces, all twenty-seven
numerical joint-image possibilities, and the two remaining level-$81$
pairs. Its exact run passed in $0.050076540999725694$ seconds.
The numerical level-$27$ list is an upper list of possibilities; no
existence claim is made for a listed group-theoretic combination.
\end{remark}

\par\smallskip\noindent\textbf{Subsequent exact determination.}
The remaining small-height alternatives isolated here are settled in
Theorem~\ref{thm:B-RD42}: the exact manuscript index is
$\nu_{\min}(3,2,2)=2$. The intermediate bounds and conditional
reductions above are retained as steps in that proof.


\subsection{A hypothetical comparison has a prescribed Frobenius graph}

\begin{proposition}[B-RD33: the graph at every smaller level]
\label{prop:B-RD33}
\status{Proved. Independent root and obstruction/crystal proof and primary-source
audits passed. The exact graph is a necessary condition and does not by itself
improve the numerical upper bound.}
Keep the completed maximal unramified base used in
Proposition~\ref{prop:B-RD32}. Let $M,N$ be height-two,
one-dimensional, potentially supersingular $3$-divisible groups
with trivial finite determinant characters and different tame
special-inertia quotient orders. Let $L/K$ be a common finite
Galois good-reduction extension, with good models
$\mathcal M,\mathcal N$ over $R=\mathcal O_L$ and algebraically
closed residue field $\bar k$. Write $M_0,N_0$ for their special
fibres. Suppose
\[
             f:M[3^n]\xrightarrow{\sim}N[3^n]
                                                        \tag{B-RD33.1}
\]
is an isomorphism over $K$.
For $1\le r\le n$, let $H_r$ be the scheme-theoretic closure of
the graph of $f|_{M[3^r]}$ in
$\mathcal M[3^r]\times_R\mathcal N[3^r]$.
Then, as closed subgroups of the indicated product,
\[
 (H_r)_{\bar k}
       =\ker(F_{M_0}^{r})\times\ker(F_{N_0}^{r}),
                                                        \tag{B-RD33.2}
\]
where $F_{M_0}^{r}:M_0\to M_0^{(3^r)}$ is the iterated
relative Frobenius, and likewise for $N_0$.
In particular a remaining comparison at level $27$ has graph
special fibre $\ker(F_{M_0}^{3})\times\ker(F_{N_0}^{3})$;
at level $81$ it has special fibre $M_0[9]\times N_0[9]$.

Put $H_0=0$ and $Q_r=H_r/H_{r-1}$. Every $Q_r$ is finite flat
of rank nine, its special fibre is $\alpha_3\times\alpha_3$,
and its generic fibre is the common representation $M[3]$.
Multiplication by $3$ induces maps
\[
                    Q_{r+1}\longrightarrow Q_r
                      \qquad(1\le r<n)
                                                        \tag{B-RD33.3}
\]
that are isomorphisms on generic fibres and zero on special
fibres. These statements are necessary conditions. No lift
of this prescribed chain with the required descent actions
is asserted here.
\end{proposition}
\begin{proof}
\par\smallskip\noindent\textbf{Perfect pairings on the full good models.}
The determinant argument of Theorem~\ref{thm:B-RD31}, with
trivial finite determinant and after the completed maximal
unramified base change, gives
\[
             \det T_3(M)=\mathbf Z_3(1),\qquad
             \det T_3(N)=\mathbf Z_3(1).
\]
Here equality means an isomorphism of rank-one Galois lattices.
Indeed the weight-zero quotient of the determinant by the
cyclotomic character has trivial finite inertia and is
unramified, hence trivial on this base. A choice of each
isomorphism supplies a perfect alternating pairing of Tate
lattices into $\mathbf Z_3(1)$, and therefore an isomorphism
of the generic $3$-divisible group with its Cartier dual.
Tate's integral full-faithfulness theorem extends this map
and its inverse over $R$; their composites remain identity
by generic-fibre faithfulness. Thus both good models have
perfect alternating pairings at every torsion level.
Alternation and Galois-descent compatibility also extend
uniquely from the generic fibre. The exact source is
\cite[Theorem~4 and Corollaries~1--2, printed p.~181]
{Tate1967PDivisible}; this argument uses full $3$-divisible
groups for extension, not finite-level full faithfulness.

Since a linear isomorphism of free rank-two
$\mathbf Z/3^n$-modules multiplies an alternating pairing by
its determinant, there is a unit
$\lambda\in(\mathbf Z/3^n)^\times$ such that
\[
            e_N(fx,fy)=e_M(x,y)^\lambda.
\]
Choose any lift of $\lambda$ to $\mathbf Z_3^\times$.
On each product $P_r=\mathcal M[3^r]\times\mathcal N[3^r]$
use the perfect alternating product pairing
$e_M\,e_N^{-\lambda^{-1}}$. The graph is maximal isotropic
on the generic fibre.

\par\smallskip\noindent\textbf{Closure and orthogonality commute with reduction.}
Each graph closure $H_r$ is finite flat over the DVR:
its coordinate algebra is the image of that of $P_r$ in
the generic graph algebra, and is finite and torsion-free.
It is a closed subgroup, of rank $3^{2r}$, stable under
$\operatorname{Gal}(L/K)$. The restricted pairing on
$H_r\times_RH_r$ is trivial, because this product is flat
and the pairing is trivial generically.

Finite-flat Cartier duality is exact, so the orthogonal
$H_r^\perp$ is finite flat of rank $3^{2r}$ and its formation
commutes with base change. There is a closed immersion
$H_r\subset H_r^\perp$. The coordinate-algebra surjection
is between free $R$-modules of equal rank; it is therefore
an isomorphism. Thus $H_r=H_r^\perp$, also on the special
fibre. Exact finite-flat duality and the rank formula are
recalled in \cite[\S1.2--1.3, printed pp.~159--160]
{Tate1967PDivisible}.

\par\smallskip\noindent\textbf{Tame idempotents split the graph.}
The norm-one possibilities in Lemma~\ref{lem:B-RD32-norm}
have tame orders $1,2,4$. Choose an element $\tau$ in a
Sylow $2$-subgroup of the joint finite descent image whose
image generates its tame quotient; Schur--Zassenhaus, or
the cyclic tame quotient and the odd wild kernel, supplies
such a lift. Its action on the special fibre is
$\bar k$-linear.

If one tame order is four and the other is one or two,
then $\tau^2$ acts as $[-1]$ on the first special group
and as $[1]$ on the second. Indeed the unique quaternionic
involution is $-1$, and rational Dieudonne faithfulness
implies the same equality of integral group endomorphisms.
Consequently $(1-\tau^2)/2$ is the projection onto the
first factor. If the two orders are one and two, use
$(1-\tau)/2$ instead. These are well-defined endomorphisms
of the $3$-primary group schemes. They preserve
$(H_r)_{\bar k}$ because that subgroup is stable under
$\tau$ and under scalar multiplication. Hence
\[
                 (H_r)_{\bar k}=J_M\times J_N
                                                        \tag{B-RD33.4}
\]
for closed subgroups $J_M\subset M_0[3^r]$ and
$J_N\subset N_0[3^r]$.

A one-dimensional connected formal group over $\bar k$
has a unique finite closed subgroup of any specified
$p$-power rank: its completed coordinate ring is
$\bar k[[T]]$, and the only ideal of colength $3^a$ is
$(T^{3^a})$. When it is a subgroup this is the kernel of
relative $3^a$-Frobenius. Thus
$J_M=\ker F_{M_0}^a$ and $J_N=\ker F_{N_0}^b$ for some
$a,b$, with $a+b=2r$.

The orthogonal of $\ker F_{M_0}^a$ in the perfectly
self-dual group $M_0[3^r]$ has rank $3^{2r-a}$.
Uniqueness makes it $\ker F_{M_0}^{2r-a}$, and similarly
for $N_0$. Since the product in (B-RD33.4) is its own
orthogonal, $a=r$ and $b=r$. This proves (B-RD33.2).
For $r=4$, uniqueness of the rank-$3^4$ subgroup identifies
$\ker F^4$ with the $9$-torsion of a height-two formal group;
no choice identifying a Frobenius twist with the group is
needed.

\par\smallskip\noindent\textbf{The successive finite-flat models.}
The graph closure in $P_r$ is also its closure in $P_n$,
since $P_r$ is a closed subgroup of $P_n$. Hence the $H_r$
form nested finite-flat closed subgroups, and their finite
flat quotients exist. Base change of the quotient sequence
is exact. Each special quotient
$\ker F^r/\ker F^{r-1}$ of a supersingular one-dimensional
group is $\alpha_3$: it is a rank-three subgroup of the
appropriate height-two Frobenius twist, whose only
rank-three subgroup is its Frobenius kernel and is local
with local Cartier dual. Thus $(Q_r)_{\bar k}=\alpha_3^2$.
The generic quotient is the usual quotient
$M[3^r]/M[3^{r-1}]\simeq M[3]$.

Multiplication by $3$ maps $H_{r+1}$ to $H_r$, by the
same generic-fibre closure argument, and maps $H_r$ to
$H_{r-1}$. It gives (B-RD33.3), generically an isomorphism.
On a height-two special formal group, multiplication by
$3$ maps $\ker F^{r+1}$ into $\ker F^{r-1}$: its kernel
has rank $3^2$ and its image in the displayed subgroup
has rank at most $3^{r-1}$, so uniqueness gives the
inclusion (with the case $r=1$ giving zero). The induced
map on the rank-three quotients is therefore zero.
This proves the last assertion.
\end{proof}

\begin{remark}[What the graph gate does not decide]\label{auto:B:18f97bffbcfd}
\status{Proved under stated hypotheses}

Bondarko's kernel theorem
\cite[Theorem~5, printed pp.~681--682]{Bondarko2006Generic}
says that the kernel of the special fibre of a generically
injective finite-flat map is killed by $F^u$ when
$e_L<2\cdot3^u$ (and also at equality for connected source).
Applied to the graph projections, (B-RD33.2) forces
$n\le u$. This recovers the proved numerical exclusions.
The additional content here is the exact special graph
and the entire chain (B-RD33.3). The refined logarithm argument
of Theorem~\ref{thm:B-RD34} subsequently excludes the
level-$81$ case. At level $27$, existence of a compatible
chain with joint inertia order $18$ or $36$ is not excluded by this proposition alone.
The remaining cases are excluded in Theorem~\ref{thm:B-RD42}.
\end{remark}


\subsection{A supersingular logarithm estimate excludes the remaining level $81$}

\begin{lemma}[A coefficientwise estimate for a product-type special fibre]
\label{lem:B-RD34-log}
\status{Proved under stated hypotheses}
Let $R$ be a complete discrete valuation ring of mixed characteristic
$(0,3)$, with perfect residue field, uniformizer $\pi$, and normalized
valuation $v(3)=a$. Let $G$ be a commutative formal group of finite height
and dimension $d$ over $R$. Suppose the coordinates of $[3]_{G_k}$ have
no monomial of total degree less than nine. Define
\[
 q_0=1,\qquad q_{i+1}=\min\{a+q_i,\,1+3q_i,\,9q_i\}.
                                                        \tag{B-RD34.1}
\]
If $q_u>a/2$, then for every $h\in tR[[t]]^d$ the normalized logarithm
satisfies
\[
       3^u\log_G(\pi h)\in\pi^{q_u}R[[t]]^d.
                                                        \tag{B-RD34.2}
\]
The hypothesis holds for a special fibre isomorphic to a product of
one-dimensional height-two formal groups, and is invariant under an
arbitrary integral change of formal coordinates.
\end{lemma}
\begin{proof}
Write the matrix of normalized invariant differentials as $U(X)$;
its entries are in $R[[X]]$ and $U(0)=1$. Multiplication by three gives
the exact differential identity
\[
     U([3]_G X)\,J_{[3]_G}(X)=3U(X).
                                                        \tag{B-RD34.3}
\]
Since $U([3]_GX)$ is invertible over $R[[X]]$, every entry of the
Jacobian $J_{[3]_G}$ is divisible by $3$. If a monomial $X^\alpha$ has
an exponent not divisible by three, its coefficient in $[3]_G$ is
therefore divisible by three. In particular this holds for every
monomial of total degree two. The linear part is $3X$. The special-fibre
hypothesis makes every coefficient of total degree three through eight
divisible by $\pi$, and all higher coefficients are integral. Thus an
augmentation vector whose coefficients have valuation at least $q\geq1$
is sent by $[3]_G$ to one whose coefficients have valuation at least
\[
                        \min\{a+q,1+3q,9q\}.
\]
Iteration starting with $\pi h$ proves the bound $q_i$ for
$[3]^i_G(\pi h)$.

For completeness this also controls the multivariable logarithm.
Write a coefficient of $X^\alpha$ in any coordinate of $\log_G$ as
$c_\alpha$, and put $m=|\alpha|\geq1$. Since $d\log_G=U(X)dX$,
\[
 v(c_\alpha)\geq
       -a\min_{j:\alpha_j>0}v_3(\alpha_j)
       \geq-a\lfloor\log_3m\rfloor.
                                                        \tag{B-RD34.4}
\]
For every integer $r\geq0$, $3^r-1\geq2r$. Consequently
$\lfloor\log_3m\rfloor\leq(m-1)/2$. If a vector $z$ has coefficient
valuations at least $q>a/2$, every coefficient of
$c_\alpha z^\alpha$ has valuation at least
$mq-a\lfloor\log_3m\rfloor\geq q$. Formal substitution is legitimate
because $z$ is an augmentation vector; each coefficient involves only
finitely many terms. Hence $\log_G(z)\in\pi^qR[[t]]^d$.
Apply this to $z=[3]^u_G(\pi h)$ and use
$\log_G([3]^u_G(\pi h))=3^u\log_G(\pi h)$.

For a one-dimensional height-two group in characteristic three,
multiplication by three starts in degree nine. This persists under
products. An integral coordinate change and its inverse have zero
constant terms and invertible linear parts. Substitution in either
direction preserves the ideal generated by monomials of total degree
at least nine, so the hypothesis is coordinate independent. No
coordinate choice on the mixed-characteristic quotient is assumed.
\end{proof}

\begin{lemma}[A refined kernel estimate]
\label{lem:B-RD34-kernel}
\status{Proved under stated hypotheses}
Keep $R,a$ as above. Let $f:F\to G$ be a homomorphism of finite-height
formal groups having no nonzero finite flat subgroup scheme in its
kernel, and with finite special-fibre kernel. This is absence of a
flat kernel; it does not require finite-level closed immersions and
does not preclude a nonzero kernel after reduction. Suppose $G_k$
satisfies the hypothesis of
Lemma~\ref{lem:B-RD34-log}. If $q_u>a/2$, then
\[
                         F^u\ker(f_k)=0,
                                                        \tag{B-RD34.5}
\]
where $F^u$ denotes iterated relative Frobenius. In particular one can
take $u=2$ when $a\leq24$, and $u=3$ when $a=72$.
\end{lemma}
\begin{proof}
We give the change to the original Cartier-module proof explicitly.
Use Bondarko's invariant module $D_G\subset L[[\Delta]]^d$, where
$L=\operatorname{Frac}(R)$, and put
\[
 \mathcal R=R[[\Delta]][1/3],\qquad
 J=\pi^{\lceil a/2\rceil}R.
\]
Here $\mathcal R$ is precisely the ring of series with coefficient
valuations bounded below. Let $G_\pi$ be the formal group law
$\pi^{-1}G(\pi X,\pi Y)$, and let $\mathcal C_G$ denote Bondarko's
completed module $M_G$. We use the exact identities
\[
 \mathcal C_G=\Omega\,\pi D_{G_\pi},\qquad
 J[[\Delta]]^d\subset D_G,\qquad
 E(G_k)=D_G/(D_G\cap\mathcal R^d).
                                                        \tag{B-RD34.6}
\]
They are respectively Proposition~10(4), Proposition~9, and Lemma~3
of \cite[printed pp.~674--675]{Bondarko2006Generic}. The canonical
embedding $D_G\to\mathcal C_G\subset L\{\{\Delta\}\}^d$ agrees
with the literal series on every bounded coefficient projection;
this is Proposition~11(1)--(2), printed p.~675.

By the definition of $D_{G_\pi}$, an element of $\pi D_{G_\pi}$,
viewed as a $3$-typical series in $t$, has the form
$\log_G(\pi h)$ for $h\in tR[[t]]^d$.
Lemma~\ref{lem:B-RD34-log} therefore gives
\[
             3^u\mathcal C_G
                  \subset\pi^{q_u}R\{\{\Delta\}\}^d
                  \subset J\{\{\Delta\}\}^d.
                                                        \tag{B-RD34.7}
\]
Indeed $\Omega$ consists of two-sided series with integral Witt
coefficients, tending to zero towards negative powers. Its action
preserves every coefficient valuation ideal; the formula for this
action is displayed immediately before Proposition~10 in the cited
source. Thus the estimate on $\pi D_{G_\pi}$ passes to its
$\Omega$-span. This replaces the general logarithm bound in
Proposition~11(5).

Here is the needed replacement for Proposition~12(1)(ii). If
$z\in D_G$, $A$ is an integral matrix, and $Az\in\mathcal R^b$, then
\[
                     \mathbf f^uAz\in J[[\Delta]]^b.
                                                        \tag{B-RD34.8}
\]
To see this, (B-RD34.7) and the bounded-projection property give
$3^uAz\in J[[\Delta]]^b$. The Cartier Frobenius satisfies
\[
        \mathbf fz=3(z-z(0))\Delta^{-1},\qquad z(0)\in R^d.
\]
Each $\mathbf f^jz$ again lies in $D_G$, so its constant term is
integral. Induction gives
\[
   \mathbf f^uAz\equiv3^uAz\Delta^{-u}
                        \pmod{3R((\Delta))^b}.
\]
Since $3R\subset J$, the right side and the error lie in
$J((\Delta))^b$. The left side has no negative powers, proving
(B-RD34.8). This is equation~(8) and the proof of Proposition~12,
printed p.~676, with the improved bound inserted.

We now spell out the final step of Proposition~16, printed p.~681.
Let $B=df(0)$, $I'=B(L^{\dim F})$, and $I=I'\cap R^d$.
Choose an integral direct-sum complement $R^d=I\oplus C$ and
the associated projection $\operatorname{pr}:R^d\to C$.
Propositions~6 and~8 identify
\[
                     B D_F=D_G\cap I'[[\Delta]].
                                                        \tag{B-RD34.9}
\]
Indeed the formal group belonging to the right side factors $f$;
the resulting map from $F$ is an isogeny, and generic injectivity
forces its finite flat kernel to vanish, so it is an isomorphism.
The same reduction identity used in Proposition~16 then embeds
$E(F_k)$ in $E(G_k)$.

Take $z\in D_G$ whose reduction $\bar z$ belongs to
$\mathbf Q_3E(F_k)\cap E(G_k)$. By (B-RD34.6), after multiplying
by a power of three, $z$ differs from an element of $BD_F$ by a
bounded-coefficient series. Hence
$\operatorname{pr}(z)\in C\otimes_R\mathcal R$.
Formula (B-RD34.8) shows
\[
 c=\mathbf f^u\operatorname{pr}(z)\in J C[[\Delta]]
                                                \subset D_G.
\]
This element is bounded, so reduces to zero in $E(G_k)$. Moreover
$\mathbf f^uz-c$ lies in the right side of (B-RD34.9). Therefore
$\mathbf f^u\bar z\in E(F_k)$. The saturation quotient
\[
 (\mathbf Q_3E(F_k)\cap E(G_k))/E(F_k)
\]
is the Dieudonne module of $\ker(f_k)$ by Theorem~2(8), equation~(6),
printed p.~668 of the same source. The operator $\mathbf f$
corresponds to relative Frobenius, as stated immediately after that
equation. Full faithfulness proves (B-RD34.5).

For the numerical assertions, if $a\leq2$ then $q_2=2a+1$;
if $3\leq a\leq9$ then $q_2=a+4$; and if $9\leq a\leq24$
then $q_2=13$. Each is greater than $a/2$. At $a=72$ the successive
values are $1,4,13,40$, and $40>36$.
\end{proof}

\begin{theorem}[B-RD34: the height-two least index has only two possibilities]
\label{thm:B-RD34}
\status{Proved. Independent parent and obstruction/crystal proof and original-source
audits passed; exact arithmetic checks also passed.}
For the character question B.4 with $p=3$, absolute ramification
$e_K=2$, $\mu_3\subset K$, and height two,
\[
                         2\leq\nu_{\min}(3,2,2)\leq3.
                                                        \tag{B-RD34.10}
\]
Thus an isomorphism of $81$-torsion always determines the requested
prime-to-$3$ character. A failure at level $27$, if one exists,
must compare two whole potentially supersingular groups with
norm-one finite special inertia and different tame quotient orders;
the order of their joint faithful good-reduction image is $18$ or
$36$.
\end{theorem}
\begin{proof}
The lower bound is the explicit level-nine comparison in
Proposition~\ref{prop:B-RD29-p3-level-two}. The reductions in
Theorem~\ref{thm:B-RD31} and Proposition~\ref{prop:B-RD32} show
that only two whole potentially supersingular groups need be
considered. At level $81$, their joint image would have order
$36$, so their common good-reduction DVR $R$ has $a=72$.
After completed unramified base change, form the graph closure
$H_n\subset P=\mathcal M\times\mathcal N$ at level $3^n$.
Proposition~\ref{prop:B-RD33} gives
\[
       (H_n)_k=\ker F_{M_0}^n\times\ker F_{N_0}^n.
                                                        \tag{B-RD34.11}
\]
The quotient $Q=P/H_n$ is a $3$-divisible group over $R$.
Its special fibre is the product of the corresponding Frobenius
twists of $M_0,N_0$, because quotient formation commutes with base
change. In particular it is a product of one-dimensional
height-two formal groups. Use the homomorphism of formal groups
\[
 g:P\longrightarrow Q\times\mathcal M,
                 \qquad z\longmapsto(qz,\operatorname{pr}_{\mathcal M}z).
                                                        \tag{B-RD34.12}
\]
The generic graph of an isomorphism meets $0\times N$ trivially,
so $g$ has no nonzero finite flat subgroup scheme in its kernel.
By (B-RD34.11), its special kernel is
$0\times\ker F_{N_0}^{n}$. Its target special fibre is a product
of three height-two groups, so Lemma~\ref{lem:B-RD34-kernel}
applies. For $n=4,a=72$ it says that $F^3$ kills
$\ker F_{N_0}^4$. This is impossible: on a one-dimensional formal
group that Frobenius image has rank three. This excludes every
remaining level-$81$ case and proves the upper bound.

For a hypothetical level-$27$ comparison, the same construction
with $n=3$ and $a\leq24$ would force $F^2\ker F_{N_0}^3=0$,
again contradicting its rank-three image. The norm-one joint-image
list in (B-RD32.4) therefore leaves only $18$ and $36$. No
existence claim for a remaining pair, and no decision between
least indices two and three, is made by this intermediate bound.
The final decision $\nu_{\min}(3,2,2)=2$ is proved in
Theorem~\ref{thm:B-RD42}.
\end{proof}

\begin{remark}[Exact arithmetic certificate]\label{auto:B:2b832ed8e5d4}
\status{Proved under stated hypotheses}

The retained program \path{code/verify_supersingular_log_bound.py}
checks $282$ multivariable monomials, $18\,432$ coefficient-valuation
cases, and $40\,000$ logarithm-degree cases. It also verifies every
ramification value $1\leq a\leq24$, the sequence $(1,4,13,40)$ at
$a=72$, the remaining joint orders $18,36$, and both rank-three
Frobenius images. Its exact author run passed in
$0.05011515400110511$ seconds. These are arithmetic checks of the
proof; no existence of a remaining level-$27$ comparison is asserted.
\end{remark}


\subsection{The remaining character comparison can be realized by elliptic curves}

\begin{lemma}[Integral conjugacy of the finite norm-one action]
\label{lem:B-RD35-conjugacy}
\status{Proved under stated hypotheses}
Let $k$ be algebraically closed of characteristic three, let
$E_0/k$ be $y^2=x^3-x$, and write $G_0=E_0[3^\infty]$.
Every finite reduced-norm-one subgroup of $\operatorname{Aut}(G_0)$
is conjugate by an automorphism of $G_0$ into
$\operatorname{Aut}(E_0)$. In particular this is integral conjugacy,
not only conjugacy by a quasi-isogeny.
\end{lemma}
\begin{proof}
Use the quaternion division algebra $D$ and its valuation ring
$\mathcal O_D=\End(G_0)$ from (B-RD31.1), normalizing the division
valuation by $v_D(3)=2$. Write
\[
 \zeta:(x,y)\longmapsto(x+1,y),\qquad
 j:(x,y)\longmapsto(-x,iy),\quad i^2=-1.
\]
These automorphisms have
\[
 \zeta^3=1,\quad j^2=-1,\quad
 j\zeta j^{-1}=\zeta^{-1},\qquad
 H_0=\langle\zeta,j\rangle\quad\hbox{of order twelve}.
                                                        \tag{B-RD35.1}
\]
Here $-1$ denotes the elliptic involution, not the identity.
Every origin-preserving automorphism of this curve has the form
$(x,y)\mapsto(u^2x+r,u^3y)$, with $u^4=1$ and $r^3-r=0$:
the pole orders of $x,y$ give the usual Weierstrass change of
variables, and comparison of the $y$-terms eliminates additions
to $y$ since $2$ is invertible. Substitution then gives the two
displayed conditions. Thus $H_0=\operatorname{Aut}(E_0)$.
Put $E=\mathbf Q_3(\zeta)\subset D$. We have
\[
 v_D(\zeta-1)=1,\qquad
 (\zeta-1)j(\zeta-1)^{-1}=-\zeta j.
                                                        \tag{B-RD35.2}
\]
The first equality follows from
$(\zeta-1)(\zeta^{-1}-1)=3$. Consequently $\zeta-1$
normalizes $H_0$ and permits adjustment of a conjugator's valuation
by any integer.

By Lemma~\ref{lem:B-RD32-norm}, a finite norm-one group $H$ is
cyclic of order $1,2,3,4$, or $6$, or is the nonabelian group of
order twelve described there. We give conjugators rather than
appealing only to rational Skolem--Noether.

If $z\in H$ has order three, the element $d=z+\zeta+1$
satisfies $dz=\zeta d$. If $d=0$, then $z=\zeta^{-1}$,
and $j$ instead conjugates $z$ to $\zeta$. Otherwise $d$ is
invertible because $D$ is a division algebra. Left multiplication
by a power of $\zeta-1$, which commutes with $\zeta$, makes $d$
a valuation-zero element. It is then an element of
$\mathcal O_D^\times=\operatorname{Aut}(G_0)$.
This handles the cyclic groups of order three and six; the unique
involution is $-1$.

For a generator $t$ of a cyclic group of order four, $t^2=-1$.
Unless $t=-j$, the nonzero element $d=t+j$ satisfies $dt=jd$.
Multiplication on the left by a suitable power of $3$ makes its
valuation either zero or one. In the latter case multiply once by
$q=\zeta-\zeta^{-1}$, which has valuation one and conjugates
$j$ to $-j$, and then divide by $3$. The resulting conjugator is
a unit and sends $\langle t\rangle$ onto $\langle j\rangle$.
The case $t=-j$ already has the desired subgroup.

Finally suppose $|H|=12$. First use the integral conjugacy just
proved to put its order-three subgroup inside $\langle\zeta\rangle$.
Its order-four generator $t$ conjugates $\zeta$ to $\zeta^{-1}$,
so $t=bj$ for some $b\in E^\times$. The equality $t^2=-1$ gives
$b\bar b=1$. If $b\ne-1$, take $c=1+b$; then $c/\bar c=b$.
If $b=-1$, take $c=\zeta-\zeta^{-1}$ instead, with the same
identity. Thus $c^{-1}tc=j$, while conjugation by $c$ fixes
$\zeta$. Multiplying this rational conjugator on the left by a
power of $\zeta-1$ makes it a unit without changing the target
group $H_0$, by (B-RD35.2). This proves integral conjugacy in
every case.
\end{proof}

\begin{proposition}[B-RD35: elliptic realization with its exact Tate lattice]
\label{prop:B-RD35}
\status{Proved. Independent parent and obstruction/crystal proof, original-source,
and exact-code audits passed.}
Let $K$ be a complete discretely valued field of mixed characteristic
$(0,3)$ with algebraically closed residue field. Let $M/K$ be a
height-two $3$-divisible group that has good reduction to a
one-dimensional supersingular formal group over a finite extension.
Assume its finite special-inertia action has reduced norm one.
Then there is an elliptic curve $E/K$ and an isomorphism of
$3$-divisible groups
\[
                             E[3^\infty]\simeq M.
                                                        \tag{B-RD35.3}
\]
Thus the comparison is integral on the full Tate lattice. It is not
merely an equality of rational representations or of finite inertia
characters.

In particular, after the completed unramified base change already
used in B-RD31--B-RD34, every remaining level-$27$ pair in the
$p=3,e_K=2,h=2$ problem is a pair of elliptic curves with
isomorphic $27$-torsion. The two curves have different tame special
inertia orders, norm-one special inertia, and joint faithful image
of order $18$ or $36$.
\end{proposition}
\begin{proof}
Choose a finite Galois good-reduction extension $L/K$, with
$R=\mathcal O_L$ and group $\Gamma$. Its residue extension is
trivial. Let $\mathcal G/R$ be the good model of $M_L$.
Tate's integral full-faithfulness theorem extends the generic
semilinear descent isomorphisms, and their inverses, to
$\mathcal G$; uniqueness preserves the cocycle relations
\cite[Theorem~4 and Corollaries~1--2, p.~181]{Tate1967PDivisible}.
Reduction gives the specified finite action of $\Gamma$ on
$\mathcal G_k$.

Lazard's height classification identifies $\mathcal G_k$ with
$G_0=E_0[3^\infty]$
\cite[Theorem~IV and its proof, pp.~270--271]{Lazard1955Formal}.
By Lemma~\ref{lem:B-RD35-conjugacy}, this identification can be
chosen so that the reduced action is the $3$-divisible-group action
of a homomorphism $\Gamma\to\operatorname{Aut}(E_0)$.
The homomorphism, including its labelling by elements of $\Gamma$,
is retained; no alteration of the inertia representation is made.

For $R_n=R/\pi^nR$, the prime three is nilpotent. Serre--Tate
equivalence gives a compatible elliptic deformation $\mathcal E_n/R_n$
of $E_0$, together with an isomorphism
$\mathcal E_n[3^\infty]\simeq\mathcal G_{R_n}$ respecting the
chosen special-fibre identification
\cite[Theorem~1.2.1, printed p.~143]{Katz1981SerreTate}.
Compatibility follows from uniqueness in that equivalence.
Each identity section is an effective relative Cartier divisor,
and the line bundles $\mathcal O_{\mathcal E_n}(3\cdot0)$ are
compatible and relatively ample. Grothendieck existence therefore
algebraizes the formal system to a projective curve
$\mathcal E/R$; the identity and group operations algebraize by
formal full faithfulness. The formal flatness and smoothness
criteria give smoothness along the closed fibre. Since $R$ is
local and $\mathcal E$ is proper, the closed nonsmooth locus
cannot avoid the closed fibre unless it is empty. Thus
$\mathcal E$ is an elliptic scheme over $R$. The exact existence
and full-faithfulness passages used here are the same as in C.1,
namely \cite[Theorems~5.4.1 and~5.4.5]{EGAIII1}.
For every $r$, the compatible finite-level Hopf algebra
isomorphisms modulo $\pi^n$ pass to an isomorphism over $R$,
because the coordinate algebras of $\mathcal G[3^r]$ and
$\mathcal E[3^r]$ are finite free and $R$ is complete.
They commute with the transitions in $r$. Hence
\[
                    \mathcal E[3^\infty]\simeq\mathcal G
                                                        \tag{B-RD35.4}
\]
integrally over $R$.

For every $\gamma\in\Gamma$, the specified special automorphism
of $E_0$ and the corresponding semilinear isomorphism of
$\mathcal G$ lift, by Serre--Tate full faithfulness on each $R_n$,
to a semilinear isomorphism of $\mathcal E_n$. The lifts are
compatible in $n$, and their cocycle relations hold by the same
uniqueness. Proper formal full faithfulness algebraizes these
isomorphisms and relations over $R$.

Restrict this descent datum to the generic elliptic curve
$\mathcal E_L$. Finite Galois descent is effective: the identity
section and its canonical ample line bundle
$\mathcal O(3\cdot0)$ are preserved, giving polarized projective
descent. We obtain an elliptic curve $E/K$. The isomorphism
(B-RD35.4) respects the descent data, so its generic restriction
descends to (B-RD35.3). Only the generic elliptic curve is
descended here; the ramified good model is not claimed to descend
to an elliptic scheme over $\mathcal O_K$.

The last assertion now follows from Theorem~\ref{thm:B-RD34}.
Conversely, such elliptic pairs are a subclass of the original
$3$-divisible-group question. Thus existence of a remaining
comparison after this unramified completion is equivalent to
existence in the elliptic subclass, with the same exact torsion
isomorphism and the same finite inertia characters.
\end{proof}

\begin{remark}[An explicit bounded-dimensional parameterization]\label{auto:B:665ea52b93eb}
\status{Proved under stated hypotheses}

Since $2$ is invertible, every curve in this reduction has an
integral Weierstrass equation
\[
        y^2=x^3+a_2x^2+a_4x+a_6,\qquad a_2,a_4,a_6\in\mathcal O_K,
                                                        \tag{B-RD35.5}
\]
with nonzero discriminant, after scaling. Hence the remaining
question concerns two triples of integral coefficients and an
isomorphism of their rank-$3^6$ generic group schemes of
$27$-torsion. The proposition justifies this reduction for every
allowed potentially supersingular BT pair; it does not evaluate
the resulting congruence problem or assert that such a pair exists.
\end{remark}

\begin{remark}[Exact quaternion checks]\label{auto:B:4074b39ba0d6}
\status{Proved under stated hypotheses}

The retained program \path{code/verify_elliptic_inertia_conjugacy.py}
works in the rational quaternion algebra
$\mathbf Q(\zeta_3)\oplus\mathbf Q(\zeta_3)j$.
It verifies the displayed normalizer and conjugator identities,
all $144$ products of the twelve automorphisms, their element-order
distribution, and the integral-normalization algorithm on eighty
exact rational conjugates. The author run passed in
$0.15318646800005808$ seconds. These finite checks illustrate the
explicit algebraic proof of all cases; they do not replace its
classification or its valuation argument.
\end{remark}


\subsection{A coefficient and discriminant form of the remaining level-$27$ test}

\begin{proposition}[B-RD36: the tame character is read from the discriminant]
\label{prop:B-RD36}
\status{Proved. Complete root and independent proof, original-source, and exact-certificate audits passed.}
Let $K$ have algebraically closed residue field of characteristic three,
and let $E/K$ be an elliptic curve with potentially supersingular
reduction. For any Weierstrass equation put
$\delta=v_K(\Delta)$. The order $t$ of its tame special-inertia
quotient is
\[
                  t=\frac4{\gcd(4,\delta)}.
                                                        \tag{B-RD36.1}
\]
Here $\gcd(4,0)=4$, and changing the equation changes $\delta$ by
a multiple of twelve. Hence the three possibilities are
\[
 \begin{array}{c|ccc}
 \delta\pmod4&0&2&1\text{ or }3\\ \hline
 t&1&2&4.
 \end{array}
                                                        \tag{B-RD36.2}
\]
Together with (B-RD32.2), this determines the requested
prime-to-$3$ character as a function on inertia.

For absolute ramification $e_K=2$, the unresolved alternative
$\nu_{\min}(3,2,2)=3$ is consequently equivalent to the existence,
for some algebraically closed $k$ of characteristic three, of two
triples $(a_2,a_4,a_6)$ in
\[
 R_k=W(k)[\pi]/(\pi^2-3),\qquad K_k=\operatorname{Frac}(R_k),
                                                        \tag{B-RD36.3}
\]
such that the integral equations
$y^2=x^3+a_2x^2+a_4x+a_6$ have nonzero discriminants, satisfy
\[
        0\leq v(\Delta)\leq23,\qquad
        3v(c_4)>v(\Delta),
                                                        \tag{B-RD36.4}
\]
have isomorphic generic $27$-torsion group schemes, and their
discriminant valuations belong to different columns of
(B-RD36.2). The convention is $v(0)=+\infty$.
This is a six-coefficient elliptic test. It is not asserted here
to be a finite residue-field elimination that has been evaluated.
\end{proposition}
\begin{proof}
Take a minimal faithful Galois good-reduction extension $L/K$.
Its group has order $m=3^wt$, with $w\in\{0,1\}$ and
$t\in\{1,2,4\}$, by the norm-one classification and the
canonical alternating pairing of an elliptic curve. A change
from the given Weierstrass equation to a good equation over
$\mathcal O_L$ has scaling factor $u$, with
\[
                  r=v_L(u)=m\delta/12\in\mathbf Z.
                                                        \tag{B-RD36.5}
\]
Choose a tame element $\tau$ of order $t$. Its action on a
uniformizer has residual multiplier a primitive $t$-th root
$\xi$. The residual multiplier of $\tau(u)/u$ is therefore
$\xi^r$. Up to inversion, it is the multiplier on the invariant
differential of the good special elliptic curve. By the explicit
automorphisms (B-RD35.1), the tangent character has precisely the
wild order-three group as its kernel. Thus the multiplier of
$\tau$ has order exactly $t$, and $\gcd(r,t)=1$.

If $t=1$, equation (B-RD36.5) gives $4\mid\delta$.
If $t>1$, $r$ is odd; taking two-adic valuations in
$12r=3^wt\delta$ gives
$2=v_2(t)+v_2(\delta)$. This proves (B-RD36.1)--(B-RD36.2).
The character identification follows from
Lemma~\ref{lem:B-RD32-norm}.

Every complete mixed-characteristic DVR of absolute ramification
two with algebraically closed residue field $k$ contains its Witt
coefficient ring. Write $3=\varpi^2v$ for a uniformizer $\varpi$
and a unit $v$. Since $k$ is algebraically closed and $2$ is
invertible, Hensel lifting gives a square root of $v$. Replacing
$\varpi$ by $\varpi\sqrt v$ gives (B-RD36.3); the degree-two
coefficient-ring description follows from the Cohen structure
argument used in Proposition~\ref{prop:B-RD24-finite-test}.
Such a field contains $\mu_3$, since its coefficient ring contains
a square root of $-1$.

By Proposition~\ref{prop:B-RD35}, every remaining BT comparison is
an elliptic comparison over one of these fields. Choose a minimal
integral Weierstrass equation and complete the square, which is an
integral change since $2$ is a unit. The discriminant bound in
Proposition~\ref{prop:B-RD24-finite-test}, with $p=3,e=2$,
gives $v(\Delta)\leq6e+11=23$.

We justify the strict $j$-condition in (B-RD36.4). Potential good
reduction is equivalent to integral $j$
\cite[Section~2, p.~496]{SerreTate1968Good}. On a good equation
in characteristic three the Cartier operator on
$\omega=dx/y$, for $y^2=x^3+a_2x^2+a_4x+a_6$, is
\[
                 C(\omega)=a_2^{1/3}\omega.
\]
Indeed $\omega=(x^3+a_2x^2+a_4x+a_6)\,dx/y^3$, and Cartier
selects the $x^2dx$ coefficient. Also $c_4\equiv a_2^2\pmod3$.
Thus a good elliptic special fibre is supersingular exactly when
its $j$-invariant is zero. Since $j=c_4^3/\Delta$, potential
supersingularity is exactly $3v(c_4)>v(\Delta)$, including $j=0$.

Conversely any pair in the stated coefficient test is potentially
supersingular and gives two allowed height-two BT groups. Its
different discriminant columns give different prime-to-$3$
characters, while its isomorphic $27$-torsion gives a failure at
manuscript index two. The proved upper bound in
Theorem~\ref{thm:B-RD34} then makes the least index three.
If no such pair exists, B-RD35 and B-RD34 instead make the least
index two. No existence or nonexistence of such a pair is inferred
from this reformulation.
\end{proof}

\par\smallskip\noindent\textbf{Exact finite certificate.}
The retained program
\path{code/verify_discriminant_tame_character.py} checks all twenty-four
displayed discriminant valuations against the necessary scaling equation
$12r=m\delta$ and the tame coprimality condition. It also checks 4,221
changes $\delta\mapsto\delta+12z$, 977 instances of the preceding crude
minimal-discriminant bound, and the exact polynomial identities
$c_4^3-c_6^2=1728\Delta$, $c_4\bmod3=a_2^2$, and the Cartier coefficient
$a_2$. The author run passed in $0.04732804800005397$ seconds;
\path{code/verify_discriminant_tame_character.output.json} retains its
full output. These finite checks certify the displayed arithmetic and
identities. They assert neither existence of all necessary scaling cases
nor existence of a level-$27$ comparison pair. Independent reruns by the
root reviewer and the second reviewer passed in respectively
$0.04524370399667532$ and $0.061048945997754345$ seconds; the author output
was preserved.

\par\smallskip\noindent\textbf{Subsequent exact determination.}
The remaining small-height alternatives isolated here are settled in
Theorem~\ref{thm:B-RD42}: the exact manuscript index is
$\nu_{\min}(3,2,2)=2$. The intermediate bounds and conditional
reductions above are retained as steps in that proof.


\subsection{Scalar residual representations are excluded by a unit obstruction}

\begin{lemma}[The five Hesse valuation patterns at absolute ramification two]
\label{lem:B-RD37-patterns}
\status{Proved. Complete root and independent proof, original-source, and exact-certificate audits passed.}
Let $K$ have algebraically closed residue field of characteristic three,
absolute ramification index two, and normalized valuation $v$.
Fix a primitive cube root $\zeta\in K$. For
\[
 E_t:\quad y^2+3txy+(t^3-1)y=x^3,\qquad t^3\ne1,
                                                        \tag{B-RD37.1}
\]
suppose $E_t$ is potentially supersingular. The ordered vector
$\boldsymbol a(t)=(v(t-1),v(t-\zeta),v(t-\zeta^2))$ has one
of the following five forms; the indicated permutations are allowed.
\[
\begin{array}{c|c|c}
\text{form}&\boldsymbol a(t)&\text{tame special-inertia order}\\\hline
\mathrm A&(-1,-1,-1)&4\\
\mathrm B&(0,0,0)&2\\
\mathrm C&(1,1,1)&4\\
\mathrm D&\text{a permutation of }(2,1,1)&2\\
\mathrm E&\text{a permutation of }(3,1,1)&4
\end{array}                                             \tag{B-RD37.2}
\]
The two sets of residue vectors modulo three in the last-column
classes two and four are disjoint, and remain disjoint after negating
either vector and permuting its coordinates.
\end{lemma}
\begin{proof}
The exact invariants are
\[
 c_4=9t(t^3+8),\qquad \Delta(E_t)=27(t^3-1)^3,\qquad
 j(E_t)=\frac{27t^3(t^3+8)^3}{(t^3-1)^3}.
                                                        \tag{B-RD37.3}
\]
If $v(t)=-a<0$, then $v(j)=6-3a$. Potential supersingularity
requires $v(j)>0$, so $a=1$, giving form A. If $t$ is integral
and its residue is not one, all three entries are zero, giving B.
This includes $t=0$, with the usual convention $v(0)=+\infty$.

It remains that $\overline t=1$. The three roots of unity have
pairwise difference valuation one, since
$(\zeta-1)^2=-3\zeta$. Therefore every entry is at least one,
and at most one is greater than one. Write the vector, after
permutation, as $(r,1,1)$ with $r\ge1$. Then
\[
 v(t^3-1)=r+2,
 \qquad t^3+8=(t^3-1)+9,
 \qquad v(9)=4.
\]
For $r=1$ the numerator factor also has valuation three and
$v(j)=6$. For $r=2$ its valuation is at least four and
$v(j)\ge6$, including $j=0$. For $r\ge3$ its valuation is
exactly four, and $v(j)=12-3r$; positivity permits only $r=3$.
This proves the list.

Put $S=\sum a_i(t)$. The displayed discriminant has valuation
$\delta=6+3S$, respectively $-3,6,15,18,21$ in the five forms.
Proposition~\ref{prop:B-RD36} gives the last column; negative raw
discriminant valuations are permitted since the same proposition
applies to any Weierstrass equation. Modulo three the order-two
vectors are $(0,0,0)$ and permutations of $(2,1,1)$. The
order-four vectors are $(2,2,2)$, $(1,1,1)$ and permutations of
$(0,1,1)$. Negating the second order-two form gives permutations
of $(1,2,2)$, which occur in neither order-four list. The asserted
disjointness follows, also for arbitrary coordinate permutations.
\end{proof}

\begin{lemma}[Cubes close to one at absolute ramification two]
\label{lem:B-RD37-units}
\status{Proved. Complete root and independent proof, original-source, and exact-certificate audits passed.}
If $u\in(K^\times)^3$ and $v(u-1)>0$, then $v(u-1)\ge3$.
\end{lemma}
\begin{proof}
Write $u=x^3$. Then $x$ is a unit and its residue has cube one.
In characteristic three this forces $\overline x=1$. Write
$x=1+w$, with $v(w)\ge1$. Each term of
$3w+3w^2+w^3$ has valuation at least three, proving the claim.
\end{proof}

\begin{proposition}[B-RD37: scalar residual representations are excluded]
\label{prop:B-RD37}
\status{Proved. Complete root and independent proof, original-source, and exact-certificate audits passed.}
Let $K$ be as in Lemma~\ref{lem:B-RD37-patterns}. Let $E,E'/K$
be potentially supersingular elliptic curves with
$E[9]\simeq E'[9]$. If their common residual representation on
three-torsion is scalar, then their finite special-inertia characters
agree on all inertia elements of order prime to three.

Consequently the common residual image in the remaining level-$27$
comparison of Theorem~\ref{thm:B-RD34} is non-scalar. After a
simultaneous quadratic twist it is either the nontrivial unipotent
cyclic group $C_3$ or the irreducible tame cyclic group $C_4$.
This excludes every scalar residual representation from both
remaining joint-inertia orders eighteen and thirty-six.
\end{proposition}
\begin{proof}
First suppose the three-torsion is rational. The full-level moduli
description of $Y(3)$ and its elliptic family, in
\citet[Section~2, p.~4; Lemma~4.2(i), p.~8]{Fisher2015Nine},
allows us to write $E$ as (B-RD37.1), with parameter $t_E\in K$.
An arbitrary isomorphism of nine-torsion has a Weil-pairing multiplier
in $(\mathbf Z/9)^\times$. Every such multiplier is plus or minus
a square. Composing with a scalar automorphism of the target torsion
therefore makes the isomorphism either symplectic or antisymplectic.
These are exactly the direct and reverse nine-congruences in
\citet[Introduction, p.~1; Section~1, first paragraph, p.~2]
{Fisher2015Nine}.

The precise original equations of
\citet[Theorem~4.3, equations~(10)--(11), pp.~9--10]
{Fisher2015Nine}, with its parameter $\Delta=1$, give a parameter
$t\in K$ representing $E'$ by the same family (B-RD37.1).
For every $\theta\in\{1,\zeta,\zeta^2\}$ they impose respectively
\[
 \frac{t-\theta}{t_E-\theta}\in(K^\times)^3
       \quad\text{in the direct case},
 \qquad
 \frac{(t-\theta^{-1})(t_E-\theta)}3\in(K^\times)^3
       \quad\text{in the reverse case}.
                                                        \tag{B-RD37.4}
\]
These are necessary equations for the actual nine-congruence, not
an inference from equality of the mod-three representations.
In particular the two valuation vectors agree modulo three in the
direct case. In the reverse case their coordinate sums are two
modulo three, with the permutation $\theta\mapsto\theta^{-1}$:
\[
       a_i(t)+a_{i^{-1}}(t_E)\equiv v(3)=2\pmod3.
                                                        \tag{B-RD37.5}
\]
The factor three in the reverse equation is essential; it is
present in the original printed equation (11).
Lemma~\ref{lem:B-RD37-patterns} excludes different tame columns
in the direct case. Equal tame quotient orders give equal trace
functions by Lemma~\ref{lem:B-RD32-norm}.

Suppose now that the congruence is reverse and the tame orders
are different. Interchanging the curves if necessary, take $E$
to have tame order two. Comparing the five forms in (B-RD37.5)
permits only B/A or D/E, with the first letter denoting the
source $t_E$. We exclude both by their unit parts; the valuation
congruence alone does not exclude them.

In case B/A, $v(t_E-1)=0$ and $v(t)=-1$. Put
\[
 R_\theta=\frac{(t-\theta^{-1})(t_E-\theta)}3.
\]
Every $R_\theta$ is a cube by (B-RD37.4), whereas
\[
 \frac{R_\zeta}{R_1}
  =\left(1+\frac{1-\zeta^2}{t-1}\right)
   \left(1+\frac{1-\zeta}{t_E-1}\right).
                                                        \tag{B-RD37.6}
\]
The two deviations have valuations two and one. The deviation of
this product from one therefore has valuation exactly one, because
its three summands have distinct valuations $2,1,3$. This contradicts
Lemma~\ref{lem:B-RD37-units}.

In case D/E, let $\theta_0$ be the source root with valuation two.
Equation (B-RD37.5) makes $\theta_0^{-1}$ the target root with
valuation three. Replace $t_E,t,\theta$ algebraically by
$t_E/\theta_0,t\theta_0,\theta/\theta_0$; the products in
(B-RD37.4) are unchanged. Thus the exceptional root is one.
Write $t_E=1+a$, $t=1+b$, with $v(a)=2$, $v(b)=3$, and put
$\lambda=\zeta-1$, $\lambda_2=\zeta^2-1$. Their exact product
is $\lambda\lambda_2=3$. The reverse cube condition at $\zeta$
becomes
\[
 R_\zeta=\frac{(a-\lambda)(b-\lambda_2)}3
      =1-\frac a\lambda-\frac b{\lambda_2}+\frac{ab}3.
                                                        \tag{B-RD37.7}
\]
The three nonconstant terms have valuations $1,2,3$, respectively.
Thus $v(R_\zeta-1)=1$, again contradicting the cube lemma.
No reverse opposite-character case remains. This proves the
proposition when the three-torsion is rational, for arbitrary
Weil-pairing multiplier.

For a scalar residual representation, write it as
$\varepsilon\operatorname{id}$ with
$\varepsilon:G_K\to\mathbf F_3^\times=\{\pm1\}$.
Use the same quadratic character to twist both elliptic curves.
Their three-torsion becomes rational and the nine-torsion
isomorphism persists. Potential supersingularity also persists.
Over a common good-reduction field the new descent action is the
old action multiplied by the scalar automorphism
$[\varepsilon(\sigma)]$. Hence both special-inertia characters
are multiplied by the same function $\varepsilon$; equality after
twisting is equivalent to equality before twisting. This proves
the assertion in the scalar case as well. The Weil-pairing
multiplier is preserved by the simultaneous twist, although the
preceding argument treats both its square classes.

For completeness, after the completed maximal unramified extension
the entire Galois group is inertia. The common residual image lies
in $\operatorname{SL}_2(\mathbf F_3)$ because $\mu_3\subset K$.
Its wild subgroup is a normal three-subgroup and its quotient is
cyclic of order prime to three. If the wild subgroup is nontrivial,
it has order three and its normalizer in
$\operatorname{SL}_2(\mathbf F_3)$ is $C_6$; the image is $C_3$
or $C_6$. A simultaneous quadratic twist removes the central
minus sign in the latter case. If the wild subgroup is trivial,
the image is cyclic of order $1,2$ or $4$. The first two are
scalar and have just been excluded. An order-four generator has
irreducible characteristic polynomial $X^2+1$ over $\mathbf F_3$.
These elementary subgroup assertions are also exhaustively checked
by the retained finite certificate. Thus the two stated
residual forms remain as possible tests; their existence as
opposite-character level-$27$ pairs is not asserted.
\end{proof}

\begin{remark}[Compatibility with the known level-nine ambiguity]
\label{rem:B-RD37-known-example}
\status{Proved under stated hypotheses}
The pair in Proposition~\ref{prop:B-RD29-p3-level-two} consists of
a potentially supersingular elliptic group and a Tate group with
good rank-one graded pieces. Moreover its common three-torsion is
non-scalar. Indeed its Kummer extension parameter is
$q=2^5 3^6$, whose cube class is nontrivial even after completed
unramified base change. If $2$ were a cube, so would be $-2$.
A putative unit $x$ with $x^3=-2$ has residue one; writing
$x=1+w$ gives $v(w)\ge1$ and
\[
 v\bigl((1+w)^3-1\bigr)
   =v(3w+3w^2+w^3)\ge3,
\]
whereas $v(-2-1)=v(-3)=2$. This is impossible.
Thus the example has nontrivial unipotent residual image $C_3$,
and supplies no exception to Proposition~\ref{prop:B-RD37}.
\end{remark}

\par\smallskip\noindent\textbf{Exact finite certificate and correction record.}
The program \path{code/verify_scalar_residual_patterns.py} checks
the nine ordered valuation patterns, all 240 opposite-column
comparisons allowing arbitrary coordinate permutations, every
nine-torsion multiplier, both displayed Weierstrass invariant
identities, both reverse-unit identities and their distinct valuation
terms, and all fifteen subgroups of
$\operatorname{SL}_2(\mathbf F_3)$. Its run passed in
$0.0822893020012998$ seconds, with the complete output retained
beside the program. In particular it uses the reverse valuation
shift \emph{two}. An earlier unpromoted argument omitted the
factor three in Fisher's original equation (11); a visual source
check caught this error. That argument is retained as withdrawn
in \path{failed_candidates/B37_scalar_exclusion_missing_factor3.withdrawn.tex}.
The current proof adds the exact unit obstructions (B-RD37.6)--
(B-RD37.7), which exclude the two reverse valuation patterns left
by the corrected equation. It thereby excludes scalar residual
level-$27$ pairs, while leaving the non-scalar cases and the least
index unresolved. Independent author-output-preserving reruns by the root
reviewer and a second reviewer passed in $0.06197141099983128$ and
$0.09311267000157386$ seconds, respectively.


\subsection{Irreducible tame residual image leaves only joint inertia eighteen}

\begin{proposition}[B-RD38: the cyclic-four residual case]
\label{prop:B-RD38}
\status{Proved. Complete root and independent proof, original-source, and exact-certificate audits passed.}
Let $K$ be a complete field of mixed characteristic $(0,3)$ with
algebraically closed residue field and absolute ramification index two.
Suppose an elliptic curve $E/K$ has residual image
$\rho_{E,3}(G_K)=C_4\subset\operatorname{SL}_2(\mathbf F_3)$.
Then
\[
                  v_K(j(E))=6,
       \qquad     6\mid v_K(\Delta(E)),
                                                        \tag{B-RD38.1}
\]
where the discriminant is that of any Weierstrass equation.
In particular $E$ is potentially supersingular and its tame
special-inertia quotient has order one or two, never four.

Consequently, in the remaining level-$27$ test of
Theorem~\ref{thm:B-RD34}, a common irreducible $C_4$ residual image
can occur with different prime-to-three characters only if the two
individual faithful special-inertia orders are $3$ and $6$, and the
joint faithful order is $18$. Every possible failure with joint order
$36$ has nontrivial unipotent residual image $C_3$ after the common
quadratic twist in Proposition~\ref{prop:B-RD37}.
\end{proposition}
\begin{proof}
Write $L=K(E[3])$. Since its Galois group has order four, $L/K$
is totally tamely ramified and cyclic. Let $F$ be its quadratic
subfield and let $\sigma$ be the nontrivial element of
$\operatorname{Gal}(F/K)$. The residual image of $G_F$ is the
scalar subgroup $\{\pm I\}$, with nontrivial character given by
$L/F$.

We first identify the projective level-three parameter carefully.
The full-level modular curve $X(3)$ carries the action of
$\operatorname{SL}_2(\mathbf F_3)/\{\pm I\}=A_4$, with quotient
the $j$-line, as in \citet[Section~2, p.~4]{Fisher2015Nine}.
Use the parameter and universal elliptic family in
\citet[Lemma~4.2(i), p.~8]{Fisher2015Nine}:
\[
 E_t:\quad y^2+3txy+(t^3-1)y=x^3,
 \qquad
 j(t)=\frac{27t^3(t^3+8)^3}{(t^3-1)^3}.
                                                        \tag{B-RD38.2}
\]
The nontrivial projective residual image is a double transposition
in $A_4$. After changing the symplectic basis, it acts by
\[
                  h(t)=\frac{t+2}{t-1}.
                                                        \tag{B-RD38.3}
\]
Here no transformation formula is being assumed without a check:
$h^2=1$, $j\circ h=j$, and $h$ interchanges the cusp pairs
$\{\infty,1\}$ and $\{\zeta,\zeta^2\}$, where $\zeta\in K$
is a primitive cube root of unity. Thus $h$ is one of the three
double transpositions, which are conjugate under $A_4$.

A symplectic basis of $E[3]$ over $L$ gives a noncuspidal point
of $X(3)$. The central scalar $-I$ acts trivially on this modular
curve, so its coordinate $t$ belongs to $F$, and
\[
                         \sigma(t)=h(t).
                                                        \tag{B-RD38.4}
\]
This permits $t\in K$ when the point is fixed by $h$; the ensuing
norm identity still holds. Put $z=t-1$. Equation (B-RD38.4)
becomes $\sigma(z)=3/z$, whence $N_{F/K}(z)=3$.

Choose $\pi\in K$ with $\pi^2=3$, as in
Proposition~\ref{prop:B-RD36}. With normalized valuations,
$v_F(3)=4$ and $v_F(\pi)=2$. Since $F/K$ is totally ramified,
$v_K(N_{F/K}(z))=v_F(z)$, so $v_F(z)=2$.
The unit $u=z/\pi$ has norm one. The residue action of
$\sigma$ is trivial, and therefore $\overline u^2=1$:
\[
                         \overline u=1\quad\hbox{or}\quad-1.
                                                        \tag{B-RD38.5}
\]
On the other hand, $(\zeta-1)^2=-3\zeta$ gives
\[
             \left(\frac{\zeta-1}{\pi}\right)^2=-\zeta.
\]
Its residual square is $-1$, so its residue is neither $1$
nor $-1$. The same holds for $(\zeta^2-1)/\pi$.
It follows that the three quantities $t-1,t-\zeta,t-\zeta^2$
all have $F$-valuation two. Thus
\[
 v_F(t^3-1)=6,\qquad v_F(t^3+8)=6,
 \qquad v_F(\Delta(E_t))=30,
 \qquad v_F(j(t))=12.                                  \tag{B-RD38.6}
\]
For the second equality use $t^3+8=(t^3-1)+9$ and $v_F(9)=8$.
Since $j(t)=j(E)\in K$, the last equality gives $v_K(j(E))=6$.

It remains to compare the discriminants of the actual elliptic
curves, rather than just their common modular point. Over $F$,
twist $E_F$ by the quadratic character of $L/F$. This kills its
scalar residual character and produces a full rational level-three
basis. Its point in $X(3)$ is exactly the point $t$ above: the
quadratic twist multiplies the basis descent by a central sign,
which is invisible in $X(3)$. The full-level moduli description
therefore identifies this twist with $E_t$ over $F$.

The quadratic extension $L/F$ is tame and ramified. It has the
form $F(\sqrt d)$ with $v_F(d)$ odd; every unit is a square by
Hensel lifting over the algebraically closed residue field.
A quadratic twist changes the discriminant by $d^6$, up to a
twelfth power from a Weierstrass coordinate change. Hence, writing
$\delta=v_K(\Delta(E))$, equation (B-RD38.6) gives
\[
                   30\equiv2\delta+6\pmod {12},
       \qquad\text{and thus}\qquad 6\mid\delta.
                                                        \tag{B-RD38.7}
\]
This comparison also covers the special values of $j$: the
full-level structure eliminates elliptic automorphisms fixing the
basis, so the identified twist is the same marked object, not
merely a curve with the same $j$-invariant.

The strict positivity of $v_K(j)$ gives potential supersingularity
by Proposition~\ref{prop:B-RD36}, whose discriminant formula now
makes the tame special-inertia order one or two. In a remaining
level-$27$ failure, Theorem~\ref{thm:B-RD34} permits only joint
orders eighteen or thirty-six. Joint order thirty-six would require
tame order four on one side, so it is excluded here. Joint order
eighteen with distinct tame characters forces tame orders one and
two and independent wild order-three images. The individual
faithful orders are therefore three and six. Finally apply the two
residual alternatives in Proposition~\ref{prop:B-RD37}.
\end{proof}

\par\smallskip\noindent\textbf{Exact finite certificate.}
The retained program \path{code/verify_c4_residual_gate.py} checks
the modular involution, its cusp permutation, the exact $c_4$ and
discriminant transformation factors, the two normalized residues in
$\mathbf F_9$, the quadratic-twist discriminant identity, and the
valuation congruence (B-RD38.7). The author run passed in
$0.07729700100026093$ seconds; its complete output is retained
beside the program.
These checks establish the finite identities used above; neither
existence nor nonexistence of a remaining $C_3$ or $C_4$ level-$27$
comparison is inferred. Independent author-output-preserving reruns by
the root and second reviewers passed in $0.1112131780027994$ and
$0.10795447899727151$ seconds, respectively.

\par\smallskip\noindent\textbf{Subsequent exact determination.}
The remaining small-height alternatives isolated here are settled in
Theorem~\ref{thm:B-RD42}: the exact manuscript index is
$\nu_{\min}(3,2,2)=2$. The intermediate bounds and conditional
reductions above are retained as steps in that proof.


\subsection{Unipotent residual image leaves only the order-twelve versus order-six pair}

\begin{lemma}[The two unit depths over absolute ramification two]
\label{lem:B-RD39-units}
\status{Proved. Complete root and independent proof, original-source, and exact-certificate audits passed.}
Let $K$ be complete of mixed characteristic $(0,3)$, with algebraically
closed residue field, and normalize $v_K(3)=2$. Choose $\pi^2=3$.
Every unit in $1+\pi^3\mathcal O_K$ is a cube. A noncube unit $D$
therefore has a well-defined cube-approximation depth $s\in\{1,2\}$:
for every unit $t$ with $v_K(t^3-D)>0$, one has
$v_K(t^3-D)=s$. The inverse unit $D^{-1}$ has the same depth.

If $L/K$ is totally ramified of degree three, and a cube
$u\in(L^\times)^3$ satisfies $0<v_L(u-1)<9$, then
$v_L(u-1)$ is either three or six.
\end{lemma}
\begin{proof}
For $b\in\mathcal O_K$, the equation
\[
       (1+\pi w)^3=1+\pi^3b
       \quad\Longleftrightarrow\quad
                     w^3+w+\pi w^2=b
                                                        \tag{B-RD39.1}
\]
has a solution in $\mathcal O_K$: its reduction has a root over the
algebraically closed residue field and its derivative is a unit.
This proves the first assertion. Choose a unit $t$ whose residue is
the unique cube root of $\overline D$. The positive valuation of
$t^3-D$ must be one or two, since valuation at least three would
make $D/t^3$, and hence $D$, a cube. Any two such $t$ differ by a
multiple of $\pi$, and their cubes differ by a multiple of $\pi^3$.
The positive valuation is therefore independent of $t$.
Replacing $t$ and $D$ by their inverses preserves it.

For the last assertion write $u=x^3$. Its residue is one, so
$x=1+w$ with $r=v_L(w)\ge1$. Since $v_L(3)=6$, the three terms in
$u-1=3w+3w^2+w^3$ have valuations $6+r,6+2r,3r$.
For $r=1,2$ their minimum is uniquely $3r$. For $r\ge3$ every term
has valuation at least nine. This proves the claim.
\end{proof}

\begin{proposition}[B-RD39: the cyclic-three residual case]
\label{prop:B-RD39}
\status{Proved. Complete root and independent proof, original-source, and exact-certificate audits passed.}
Let $K$ be as in Lemma~\ref{lem:B-RD39-units}. Suppose two potentially
supersingular elliptic curves $E,E'/K$ have isomorphic nine-torsion,
common nontrivial unipotent residual image $C_3$, and different
prime-to-three special-inertia characters. Then their common
three-division field is generated by a cube root of a noncube
\emph{unit} of $K$. Their tame special-inertia quotient orders are
two and four, and their $j$-valuations, in some order, are
\[
                            3\quad\hbox{and}\quad b\ge6,
                                                        \tag{B-RD39.2}
\]
where $b=+\infty$ is allowed.

Consequently every remaining level-$27$ comparison with this residual
image has individual faithful special-inertia orders $12$ and $6$,
joint faithful order $36$, and independent wild order-three images.
Combined with Proposition~\ref{prop:B-RD38}, after the simultaneous
quadratic normalization of Proposition~\ref{prop:B-RD37}, the entire
remaining level-$27$ test has the following necessary alternatives:
\[
\begin{array}{c|c|c|c}
\text{residual image}&\text{individual special orders}&
 \text{joint order}&\text{$j$-valuations}\\\hline
C_4&(3,6)&18&(6,6)\\
C_3&(12,6)&36&(3,b),\ b\ge6.
\end{array}                                             \tag{B-RD39.3}
\]
The table asserts necessary restrictions, not existence of either kind
of comparison. Undoing the common quadratic twist may interchange
the tame orders one and two; it leaves order four and both
$j$-invariants unchanged.
\end{proposition}
\begin{proof}
A nontrivial unipotent residual image fixes a nonzero three-torsion
point. By \citet[Lemmas~4.1--4.2, p.~8]{Fisher2015Nine}, write
\[
 E=E_{t_E,D}:\quad y^2+3t_Exy+(t_E^3-D)y=x^3,
 \qquad L=K(E[3])=K(\sqrt[3]D).
                                                        \tag{B-RD39.4}
\]
Concretely, in the rational-three-point form
$y^2+a_1xy+a_3y=x^3$, take $t_E=a_1/3$ and
$D=t_E^3-a_3$. Its actual discriminant is $(3a_3)^3D$,
so its cube class is precisely the parameter class $D$.
The element $D$ is not a cube. A Weierstrass scaling replaces
$(t_E,D)$ by $(t_E/c,D/c^3)$, so we normalize
$d=v_K(D)\in\{0,1,2\}$. The common residual field $L/K$ is
cyclic and totally ramified of degree three. Fix $\theta\in L$
with $\theta^3=D$; then $v_L(\theta)=d$ and
$v_L(\zeta-1)=3$.

An arbitrary nine-congruence can be made direct or reverse by a scalar
change, as in Proposition~\ref{prop:B-RD37}. Fisher's exact
Theorem~4.3, equations~(10)--(11), pp.~9--10, supplies a parameter
$t\in K$ for $E'$ and requires
\[
\begin{array}{ll}
 \displaystyle\frac{t-\theta}{t_E-\theta}\in(L^\times)^3,
   &E'=E_{t,D}\quad\text{(direct)},\\[6pt]
 \displaystyle\frac{(t-\theta^{-1})(t_E-\theta)}3
                 \in(L^\times)^3,
   &E'=E_{t,D^{-1}}\quad\text{(reverse)}.
\end{array}                                             \tag{B-RD39.5}
\]
Each condition also holds with $\theta$ replaced by $\zeta\theta$
or $\zeta^2\theta$. The reverse denominator three is retained.

Here are all the required valuations, with no minimal-model
assumption. For $r=v_K(t)$ and $S=v_K(t^3-D)$, the exact invariants are
\[
 c_4(E_{t,D})=9t(t^3+8D),\qquad
 \Delta(E_{t,D})=27D(t^3-D)^3.
                                                        \tag{B-RD39.6}
\]
If $d=1$ or $2$, the valuations $3r$ and $d$ are unequal.
Potential supersingularity requires $r\ge-1$, and
\[
             S=\min(3r,d),\quad
             v_K(j)=6+3r-d,\quad
             \delta=6+d+3S.
                                                        \tag{B-RD39.7}
\]
The value $t=0$ belongs to the last row below and has $j=0$.
By Proposition~\ref{prop:B-RD36}, the tame quotient is read from
$\delta$ modulo four. The full direct-source tables are
\[
\begin{array}{c|ccc|ccc}
 &\multicolumn{3}{c|}{d=1}&\multicolumn{3}{c}{d=2}\\
r&S&\delta&\text{tame}&S&\delta&\text{tame}\\\hline
-1&-3&-2&2&-3&-1&4\\
0&0&7&4&0&8&1\\
\ge1&1&10&2&2&14&2
\end{array}                                             \tag{B-RD39.8}
\]
For the reverse target, whose parameter is $D^{-1}$, put
$S'=v_K(t^3-D^{-1})$. The analogous complete table is
\[
\begin{array}{c|ccc|ccc}
 &\multicolumn{3}{c|}{d=1}&\multicolumn{3}{c}{d=2}\\
r&S'&\delta'&\text{tame}&S'&\delta'&\text{tame}\\\hline
-2&-6&-13&4&-6&-14&2\\
-1&-3&-4&1&-3&-5&4\\
\ge0&-1&2&2&-2&-2&2
\end{array}                                             \tag{B-RD39.9}
\]

Take norms in (B-RD39.5). Since a norm of a cube is a cube in $K$,
the necessary congruences are
\[
    S-S_E\equiv0\pmod3\quad\text{(direct)},\qquad
    S'+S_E\equiv0\pmod3\quad\text{(reverse)}.
                                                        \tag{B-RD39.10}
\]
In the reverse case the norm of the denominator is $27$, whose
$K$-valuation is six, hence zero modulo three.

Suppose first $d=1$ or $2$. In the direct case, the only
opposite-tame pairs compatible with (B-RD39.10) have
$\{S,S_E\}=\{0,-3\}$. Interchange the pair if needed so that
$S_E=-3$, $S=0$. The quotient of the cube conditions at
$\zeta\theta$ and $\theta$ is
\[
 \frac{1+(1-\zeta)\theta/(t-\theta)}
      {1+(1-\zeta)\theta/(t_E-\theta)}.
                                                        \tag{B-RD39.11}
\]
The two displayed deviations have $L$-valuations $3+d$ and $6+d$.
The quotient is congruent to one with deviation valuation exactly
$3+d$, which is four or five. It cannot be a cube by
Lemma~\ref{lem:B-RD39-units}.

In the reverse case, compare the two tables using (B-RD39.10).
All opposite-tame pairs and the two deviation valuations in their
cube quotient are as follows:
\[
\begin{array}{c|rr|rr}
d&S_E&S'&3-d-S'&3+d-S_E\\\hline
1&-3&-6&8&7\\
1&-3&-3&5&7\\
1&0&-3&5&4\\
2&-3&-6&7&8\\
2&0&-6&7&5\\
2&0&-3&4&5
\end{array}                                             \tag{B-RD39.12}
\]
Indeed, on dividing the condition at $\zeta\theta$ by that at
$\theta$, one obtains the product
\[
 \left(1+\frac{(1-\zeta^{-1})\theta^{-1}}
                      {t-\theta^{-1}}\right)
 \left(1+\frac{(1-\zeta)\theta}{t_E-\theta}\right).
                                                        \tag{B-RD39.13}
\]
The valuations of its two deviations are the last two columns.
They are unequal; their minimum is four, five, or seven. The cross
term has strictly larger valuation. Thus the deviation of this
required cube from one has valuation four, five, or seven, again
contrary to Lemma~\ref{lem:B-RD39-units}. This excludes $d=1,2$
for either orientation.

We are left with the noncube unit $D$, of depth $s=1$ or $2$ in
Lemma~\ref{lem:B-RD39-units}. Its inverse has the same depth.
For either of the parameters $D,D^{-1}$, potential
supersingularity and the unit-depth lemma give exactly the following
possibilities:
\[
\begin{array}{c|c|c|c}
S&\delta&\text{tame}&v_K(j)\\\hline
-3&-3&4&3\\
0&6&2&\ge6\\
1\ (s=1)&9&4&6\\
2\ (s=2)&12&1&6
\end{array}                                             \tag{B-RD39.14}
\]
For $S=-3$ one has $v_K(t)=-1$. For $S=0$ the parameter is
integral and its residue does not equal the cube root of the unit
parameter; the zero value of $t$ is allowed. If $S>0$, $t$ is a
unit and $v_K(t^3+8D)=S$, since $S< v_K(9D)=4$.
These observations and (B-RD39.6) verify the table.

The direct norm congruence in (B-RD39.10) allows different tame
columns only for $\{S,S_E\}=\{-3,0\}$. The reverse norm
congruence likewise allows different columns only for that pair:
both curves have the same positive depth $s$, and $2s$ is not
divisible by three. Formula (B-RD39.14) now proves (B-RD39.2) and
the tame orders two and four.

At level $27$, Theorem~\ref{thm:B-RD34} restricts the joint faithful
order to eighteen or thirty-six. A tame quotient of order four
rules out eighteen, so the joint order is thirty-six. Its wild
subgroup has order nine; the two individual wild quotients have
order three and are independent. Their distinct tame orders four
and two therefore give individual orders twelve and six. The
other row of (B-RD39.3) is Proposition~\ref{prop:B-RD38}.
\end{proof}

\par\smallskip\noindent\textbf{Exact certificate and scope.}
The program \path{code/verify_c3_residual_gate.py} checks the
Weierstrass invariants, the exact Hensel polynomial, both source
and reverse-target tables, every opposite-tame norm-congruence
case, and all the required forbidden cube deviations. Its output
records the author runtime $0.053551534001599066$ seconds. Independent
author-output-preserving reruns by the root and second reviewers passed
in $0.045766205999825615$ and $0.05912446099682711$ seconds.
These calculations narrow the remaining test;
they do not construct or exclude either final row of (B-RD39.3).


\subsection{The remaining cyclic-four comparison must reverse the pairing}

\begin{proposition}[B-RD40: no direct nine-congruence in the cyclic-four branch]
\label{prop:B-RD40}
\status{Proved. Complete root and independent proof, original-source, and exact-certificate audits passed.}
Let $K$ be as in Proposition~\ref{prop:B-RD38}, and let $E,E'/K$
have common irreducible cyclic-four residual representation. If there
is a symplectic isomorphism $E[9]\simeq E'[9]$, their requested
prime-to-three special-inertia characters agree.

Consequently any opposite-character level-$27$ comparison in the
$C_4$ row of (B-RD39.3) has pairing multiplier congruent to $-1$
modulo three. Every such comparison can be made antisymplectic by a
scalar change; none can be made symplectic. This is an orientation
restriction, not an exclusion of the entire $C_4$ row.
\end{proposition}
\begin{proof}
Put $L=K(E[3])=K(E'[3])$ and let $F$ be its quadratic subfield.
Use the given symplectic identification to choose symplectic bases of
both three-torsion groups with the same actual $C_4$ Galois matrix.
As in Proposition~\ref{prop:B-RD38}, choose the bases so that its
projective action on $X(3)$ is
\[
 h(t)=\frac{t+2}{t-1},\qquad \sigma(t)=h(t),
                                                        \tag{B-RD40.1}
\]
where $\sigma\in\operatorname{Gal}(F/K)$ is nontrivial.
Write $t_E,t'$ for the two parameters, $z=t_E-1$, $z'=t'-1$,
and choose $\pi^2=3$ in $K$. Thus
\[
 u=z/\pi,\quad u'=z'/\pi,\qquad
 N_{F/K}(u)=N_{F/K}(u')=1,\qquad
 \overline u,\overline{u'}\in\{1,-1\}.
                                                        \tag{B-RD40.2}
\]

\par\smallskip\noindent\textbf{The actual marked differential cocycle.}
We determine the lift of (B-RD40.1), rather than using only its
projective action. Let $\zeta$ be a primitive cube root and
$c=\zeta-\zeta^2$, so $c^2=-3$. On the Hesse pencil
\[
 F_t(x,y,z)=x^3+y^3+z^3-3txyz
\]
use the matrix
\[
 M=\begin{pmatrix}1&1&1\\1&\zeta&\zeta^2\\
                         1&\zeta^2&\zeta\end{pmatrix}.
\]
Direct polynomial calculation gives
\[
 \det M=-3c,\qquad
 F_t(MX)=3(1-t)F_{h(t)}(X).
                                                        \tag{B-RD40.3}
\]
For the Poincare residue differential
$\eta_t=\operatorname{Res}(\Omega/F_t)$, where
$\Omega=x\,dy\wedge dz-y\,dx\wedge dz+z\,dx\wedge dy$, this gives
\[
                  M^*\eta_t=\kappa(t)\eta_{h(t)},
        \qquad    \kappa(t)=\frac c{t-1},
        \qquad    \kappa(t)\kappa(h(t))=-1.
                                                        \tag{B-RD40.4}
\]
The two Hesse translation generators are the cyclic coordinate
permutation and $\operatorname{diag}(1,\zeta,\zeta^2)$.
Neither generator has a fixed point on a smooth member of the
pencil: the diagonal generator fixes only coordinate vertices, and
the eigenvectors of the cyclic permutation lie on the cubic only
when $t=1$. Thus both are translations; they generate all nine
three-torsion translations. Conjugation by $M$ interchanges them
with one inverted, so its marked action is an order-four symplectic matrix, whose square is
$-I$. Moving the image flex back to the chosen origin is a
translation and does not change the invariant differential.
The possible central sign is fixed by the actual common marked
matrix on $E[3]$ and $E'[3]$, hence is the same on both sides.
The use of these marked Hesse translations agrees with the full-level
moduli convention in \citet[Section~2, pp.~4--5]{Fisher2015Nine}.

We also check that passing to the Weierstrass family used in
Fisher's cube equations introduces no parameter-dependent factor.
On the plane cubic set
\[
              X=z,\qquad Y=x-y,\qquad Z=x+y+tz.
\]
Its equation becomes
\[
 3ZY^2=4(t^3-1)X^3-9t^2X^2Z+6tXZ^2-Z^3.
                                                        \tag{B-RD40.5}
\]
Put $a=4(t^3-1)/3$ and take affine coordinates $aX/Z,aY/Z$,
followed by the usual $x$-translation to remove the square term.
The resulting short Weierstrass coefficients are
\[
 A_H=-\frac{t(t^3+8)}3,\qquad
 B_H=\frac{2(t^6-20t^3-8)}{27}.
\]
The displayed family $E_t:y^2+3txy+(t^3-1)y=x^3$ instead has
\[
 A_E=-\frac{3t(t^3+8)}{16},\qquad
 B_E=-\frac{t^6-20t^3-8}{32}.
\]
They are related by the constant Weierstrass scaling $2/c$.
The determinant of the preceding plane coordinate change is two;
tracking the residue differential through (B-RD40.5) gives
\[
                         \eta_t=\frac c3\,\omega_t,
                                                        \tag{B-RD40.6}
\]
where $\omega_t$ is the invariant differential of $E_t$ under this
identification. Thus (B-RD40.4) is also the exact marked differential
factor for this Weierstrass family. This computation supplies the
normalization behind the constant twist mentioned in
\citet[Lemma~4.2(i), p.~8]{Fisher2015Nine}.

\par\smallskip\noindent\textbf{Different tame characters force opposite unit residues.}
Take marked isomorphisms $E_L\simeq E_{t_E,L}$ and
$E'_L\simeq E_{t',L}$, and let their pullbacks of the universal
invariant differentials be $\alpha\omega_E$ and
$\alpha'\omega_{E'}$, respectively. For a generator
$\tau\in\operatorname{Gal}(L/K)$, equation (B-RD40.4) and the
common actual marked matrix show that
\[
       \frac{\tau(\alpha'/\alpha)}{\alpha'/\alpha}
                =\frac{z'}z
          \quad\hbox{or its inverse},
                                                        \tag{B-RD40.7}
\]
according to the simultaneous choice of generator and direction of
the universal isomorphism. Either choice has the same residue for
the argument below. A common central minus sign cancels in the ratio.

Let $\delta,\delta'$ be discriminant valuations over $K$, and
put $n=v_L(\alpha)$, $n'=v_L(\alpha')$. The invariant discriminant
section $\Delta\omega^{12}$ and the valuation
$v_L(\Delta(E_t))=60$ from Proposition~\ref{prop:B-RD38} give
\[
                 4\delta=60+12n,\qquad
                 4\delta'=60+12n'.
                                                        \tag{B-RD40.8}
\]
If the tame characters differ, their orders are one and two, so
$\delta'-\delta\equiv2\pmod4$. Hence
$n'-n\equiv2\pmod4$. In a tame cyclic extension of degree four,
the residue of $\tau(\beta)/\beta$ for an element of valuation
$m$ is $\xi^m$, where $\xi$ is a primitive fourth root.
Apply this to $\beta=\alpha'/\alpha$. Equation (B-RD40.7) forces
\[
                           \overline{u'}=-\overline u.
                                                        \tag{B-RD40.9}
\]

\par\smallskip\noindent\textbf{The direct nine-congruence contradicts a norm cube.}
Over $F$, twist both curves by the same quadratic character of $L/F$.
Their three-torsion becomes rational, they become the marked curves
$E_{t_E},E_{t'}$, and the symplectic nine-congruence persists.
Fisher's direct equation (10) therefore makes
\[
       \frac{t'-\theta}{t_E-\theta}\in(F^\times)^3
          \qquad(\theta=1,\zeta,\zeta^2).
                                                        \tag{B-RD40.10}
\]
At $\theta=1$ this says $u'/u$ is a cube. Its negative is a cube
congruent to one by (B-RD40.9). Since $v_F(3)=4$, the expansion
$(1+w)^3-1=3w+3w^2+w^3$ shows that any such cube has deviation
valuation at least three. Thus $v_F(u'+u)\ge3$.

Set
\[
 T=z+3/z=\pi(u+u^{-1}),\qquad
 T'=z'+3/z'=\pi(u'+u'^{-1})\in K.
\]
Both $T$ and $T'$ have $K$-valuation one. Moreover
\[
       T'+T=\pi(u'+u)\bigl(1+(u'u)^{-1}\bigr).
\]
The second factor has positive $F$-valuation by (B-RD40.9), and
$v_F(\pi)=2$, so $v_K(T'+T)\ge3$.
For $\lambda=\zeta-1$, the exact relation
$\lambda^2=-3\zeta$ gives
\[
 N_{F/K}(t_E-\zeta)
       =(z-\lambda)(3/z-\lambda)=-\lambda(T+3).
\]
Taking the norm of the cube in (B-RD40.10) at $\theta=\zeta$
makes
\[
                 q=\frac{T'+3}{T+3}\in(K^\times)^3.
\]
But
\[
              -q-1=-\frac{T'+T+6}{T+3}
\]
has $K$-valuation exactly one: its numerator has valuation two
and its denominator valuation one. This contradicts
Lemma~\ref{lem:B-RD37-units}, since $-q$ is a cube congruent to one.
The assumed different tame characters are impossible.

Finally, a unit of $\mathbf Z/27$ congruent to one modulo three is
a square, by Hensel lifting for $X^2-a$. A torsion isomorphism with
such a multiplier could be rescaled to be symplectic, contradicting
the result just proved on nine-torsion. Every remaining multiplier
is congruent to $-1$, and its negative is a square, giving the stated
antisymplectic normalization.
\end{proof}

\par\smallskip\noindent\textbf{Exact certificate and scope.}
The program \path{code/verify_c4_orientation_gate.py} checks all
Fourier, cubic, differential-scaling and short-Weierstrass identities,
the exact norm identity, the valuation inequalities, and every
unit multiplier modulo twenty-seven. Its retained output records
the author runtime $0.1905790330056334$ seconds. An initial
certificate used structural equality for two differently factored
SymPy expressions; that failed version and traceback are retained in
\path{failed_candidates/B40_certificate_structural_equality_v1.py} and
its adjacent output. The corrected test proves their exact difference
is zero. No antisymplectic level-$27$ pair is constructed or
excluded by this proposition alone; it is excluded in
Theorem~\ref{thm:B-RD42}. Independent author-output-preserving reruns by the
root and second reviewers passed in $0.1678152059976128$ and
$0.1795537689977209$ seconds, respectively.


\subsection{A three-isogeny eliminates the unipotent residual branch}

\begin{lemma}[The normalized three-isogeny and its invariants]
\label{lem:B-RD41-isogeny}
\status{Proved; independently checked by the root and obstruction agents, including the complete proof and exact certificate.}
Over a characteristic-zero field let
\[
 E:\ y^2+a_1xy+a_3y=x^3,\qquad a_3(a_1^3-27a_3)\ne0.
\]
The quotient by the subgroup generated by $(0,0)$ is
\[
 E_1:\ y^2+a_1xy+a_3y
       =x^3-5a_1a_3x-a_3(a_1^3+7a_3).
                                                        \tag{B-RD41.1}
\]
Its invariants are
\[
 c_4(E_1)=a_1(a_1^3+216a_3),\qquad
 \Delta(E_1)=a_3(a_1^3-27a_3)^3.
                                                        \tag{B-RD41.2}
\]
For the Hesse-form coefficients $a_1=3t$, $a_3=t^3-D$, this gives
\[
          j(E_1)=\frac{27t^3(9t^3-8D)^3}{D^3(t^3-D)}.
                                                        \tag{B-RD41.3}
\]
\end{lemma}
\begin{proof}
The point $(0,0)$ has order three, and its inverse is $(0,-a_3)$,
as in \citet[Lemma~4.1, p.~8]{Fisher2015Nine}. Here the quotient
can be checked directly. On the open set $x\ne0$ set
\[
\begin{split}
 X&=x+\frac{a_1a_3}{x}+\frac{a_3^2}{x^2},\\
 Y&=y-\frac{a_3(a_1x+2a_3)y}{x^3}
        -\frac{a_3(a_1x+a_3)^2}{x^3}.
\end{split}                                             \tag{B-RD41.4}
\]
Substitution into (B-RD41.1), followed by reduction using the source
equation, gives zero identically. Also
\[
          2Y+a_1X+a_3=(2y+a_1x+a_3)\frac{dX}{dx},
\]
so the invariant differential is normalized. These are universal
polynomial identities after clearing denominators, checked by the
retained exact certificate.

The rational map extends to the smooth projective curves. The
function $X$ has pole divisor exactly
$2(O+(0,0)+(0,-a_3))$: at the two finite points $x$ is a uniformizer
and the term $a_3^2/x^2$ has nonzero leading coefficient. At $O$,
the leading term is $x$. Therefore the morphism has degree three
and sends all three displayed points to the origin. An origin-preserving
morphism of elliptic curves is a group homomorphism (apply rigidity
to the difference of the two addition maps). Its kernel is exactly
the indicated subgroup, so this is the quotient isogeny.
The standard Weierstrass invariant polynomials applied to
(B-RD41.1) give (B-RD41.2); substituting the Hesse-form coefficients
gives (B-RD41.3).
\end{proof}

\begin{proposition}[B-RD41: no level-twenty-seven failure with unipotent residual image]
\label{prop:B-RD41}
\status{Proved; independently checked by the root and obstruction agents, including the complete proof and exact certificate.}
Let $K$ be complete of mixed characteristic $(0,3)$, with algebraically
closed residue field and absolute ramification index two. No two
potentially supersingular elliptic curves with isomorphic $27$-torsion,
different prime-to-three special-inertia characters, and common
residual image $C_3$ or $C_6$ exist.

Thus any still-hypothetical height-two comparison in B.4 with
$p=3,e_K=2$ is restricted to the cyclic-four row:
\[
 \rho_3(G_K)=C_4,\qquad
 (|I_E|,|I_{E'}|)=(3,6),\qquad
 |I_{E,E'}|=18,\qquad
 v_K(j(E))=v_K(j(E'))=6.
                                                        \tag{B-RD41.5}
\]
Here $I_E,I_{E'}$ denote faithful special-inertia images, and the
pair is unordered. This proposition alone leaves least indices two
and three; Theorem~\ref{thm:B-RD42} subsequently excludes
(B-RD41.5) and proves that the least index is two.
\end{proposition}
\begin{proof}
If the common residual image is $C_6$, apply the same quadratic
twist to both curves to remove its central sign. The $27$-torsion
isomorphism persists, both $j$-invariants are unchanged, and
inequality of the two requested characters persists because both
are multiplied by the same quadratic character. We may therefore
assume the common image is the nontrivial unipotent group $C_3$.

By Proposition~\ref{prop:B-RD39}, the given comparison on nine-torsion
has the unit-parameter Hesse forms described there. After possibly
interchanging the curves, the first has parameter valuation $r=-1$
and $S=v_K(t^3-D)=-3$, and the second has $r\ge0$ and $S=0$.
The parameters $D$ and, in the reverse case, $D^{-1}$ are units.
The tame special-inertia orders are four and two.

Let $C\subset E[3]$ and $C'\subset E'[3]$ be their unique invariant
lines. These are the subgroups generated by the rational points
$(0,0)$ in the respective Hesse forms. Uniqueness follows because
a nontrivial unipotent matrix in dimension two has a unique invariant
line. The $27$-torsion isomorphism $f$ therefore carries $C$ onto $C'$.
Write $\varphi:E\to E_1=E/C$ and
$\varphi':E'\to E'_1=E'/C'$.

We spell out the level change. The inverse image of $E_1[9]$ under
$\varphi$ is
\[
 T=\{P\in E[27]:[9]P\in C\}.
\]
Indeed $[9]\varphi(P)=0$ is equivalent to $[9]P\in C$, which
also forces $[27]P=0$. Thus $E_1[9]=T/C$ as finite group schemes.
The same construction on the other side gives $T'/C'$, and $f$
induces an isomorphism
\[
                             E_1[9]\simeq E'_1[9]
                                                        \tag{B-RD41.6}
\]
over $K$. Both quotient curves are potentially supersingular.
An isogeny induces an isomorphism of rational Tate modules and of
their potential good-reduction Dieudonne modules. In particular the
special-inertia characters and their tame orders remain four and two.

Formula (B-RD41.3) gives a different pair of $j$-valuations. If
$r=-1$ and $D$ is a unit, then
$v_K(9t^3)=1$ and $v_K(8D)=0$, so
\[
                        v_K(j(E_1))=6.
\]
For the other curve, $r\ge0$ and $S=0$, while its unit parameter
is either $D$ or $D^{-1}$. Again $9t^3-8D$ (or its inverse-parameter
version) is a unit, so
\[
                       v_K(j(E'_1))=6+3r\ge6,
                                                        \tag{B-RD41.7}
\]
with $+\infty$ allowed for $t=0$.

Now classify the common residual representation in (B-RD41.6).
Since $\mu_3\subset K$, it lies in
$\operatorname{SL}_2(\mathbf F_3)$; the elementary inertia
classification in Proposition~\ref{prop:B-RD37} gives scalar image,
cyclic-four image, or cyclic-three/six image. A scalar image is
impossible by Proposition~\ref{prop:B-RD37}, since the characters
are different. Cyclic-four image is impossible by
Proposition~\ref{prop:B-RD38}, since one tame special-inertia order
is four. In the remaining case, a common quadratic twist gives
nontrivial unipotent residual image. It leaves both valuations in
(B-RD41.7) unchanged and preserves inequality of the characters.
Proposition~\ref{prop:B-RD39} would then require one $j$-valuation
to be three, contradicting (B-RD41.7). This excludes (B-RD41.6)
and hence the proposed original $27$-torsion comparison.

The reductions in Theorem~\ref{thm:B-RD34} and
Propositions~\ref{prop:B-RD35}, \ref{prop:B-RD37} identify every
remaining BT comparison with one of these elliptic residual cases.
Only the cyclic-four case survives, and
Proposition~\ref{prop:B-RD38} gives exactly (B-RD41.5).
\end{proof}

\par\smallskip\noindent\textbf{Exact certificate.}
The program \path{code/verify_c3_isogeny_exclusion.py} checks the
universal rational-map substitution and normalized differential,
both target invariant polynomials, the specialized quotient
$j$-formula, and all required valuation cases. It also models the
finite module passage from $27$-torsion to quotient nine-torsion,
with the actual invariant line and subgroup ranks. The retained
output records the author runtime $0.08725505400070688$ seconds.
It does not evaluate the remaining
cyclic-four antisymplectic branch.


\subsection{The reverse cyclic-four case and the exact smallest height-two index}

\begin{lemma}[The two cyclic-four wild families and their cubic fields]
\label{lem:B-RD42-normal}
\status{Proved. Complete root and independent original-source, proof and exact-certificate audits passed.}
Let $K$ be complete of mixed characteristic $(0,3)$, with algebraically
closed residue field and $v_K(3)=2$. Fix $\pi^2=3$ and $i^2=-1$ in $K$.
A potentially supersingular elliptic curve with irreducible cyclic-four
three-torsion representation and faithful special inertia of order six
is $K$-isomorphic to
\[
 M_b:\ y^2=x^3-3x+3b,\qquad b\in\OO_K^\times.
                                                        \tag{B-RD42.1}
\]
A curve with the same residual condition and faithful special inertia
of order three is the quadratic $\pi$-twist $M_b^{(\pi)}$ of such a curve.
The minimal discriminant valuations are respectively six and twelve.

Put $L_b=K(r)$, where $r^3-3r+3b=0$. This is a cyclic cubic extension,
and $M_b^{(\pi)}$ has good reduction over $L_b$, with faithful special
inertia of order three. The curve $M_b$ has good reduction over
$L_b(\sqrt\pi)$, with faithful special inertia of order six. Moreover
\[
 L_b=L_{b'}\quad\Longleftrightarrow\quad
                   \overline {b'}=\overline b\ \hbox{or}\ -\overline b.
                                                        \tag{B-RD42.2}
\]
\end{lemma}
\begin{proof}
The discriminant bound in Proposition~\ref{prop:B-RD36} and
Proposition~\ref{prop:B-RD38} give
$0\le\delta\le23$, $6\mid\delta$, and $v_K(j)=6$.
Thus $v_K(c_4)=(\delta+6)/3$. Passing to the short equation
$y^2=x^3+Ax+B$, with $A=-c_4/48$ and $B=-c_6/864$, gives an
integral equation when $\delta=6,12,18$. Indeed
$c_6^2=c_4^3-1728\Delta$ gives respectively
$v_K(c_6)\ge6,9,12$, so the pairs of coefficient valuations satisfy
\[
 (v_K(A),v_K(B))=(2,\ge0),\ (4,\ge3),\ (6,\ge6).
\]
The last equation scales down by $x=\pi^2X,y=\pi^3Y$, contradicting
minimality. The tame-order formula of Proposition~\ref{prop:B-RD36}
therefore forces $\delta=6$ for inertia order six and $\delta=12$
for order three; $\delta=0$ would mean good reduction.

In the first case write $A=3a$, $a$ a unit. The four three-torsion
lines have $x$-coordinates given by
\[
                     x^4+6ax^2+4Bx-3a^2=0.
                                                        \tag{B-RD42.3}
\]
If $v_K(B)=0$, its Newton polygon has one root of valuation two
and three of valuation zero. If $v_K(B)=1$, it has one root of
valuation one and three of valuation $1/3$. In each case the unique
root valuation gives a Galois-stable three-torsion line, contrary
to irreducibility of the cyclic-four representation. Hence
$v_K(B)\ge2$. If $v_K(B)\ge3$, the $\pi$-twist has coefficient
valuations at least $(4,6)$; scaling by $\pi$ makes its discriminant
a unit. The original curve would acquire good reduction over a
quadratic extension, contrary to its wild inertia of order three.
Thus $B=3b_0$ with $b_0$ a unit. A fourth root of $-a$ exists in $K$
by Hensel's lemma and normalizes $a$ to $-1$, giving (B-RD42.1).

For the order-three curve, twist its integral short equation by
$\pi$ and scale by $\pi$. The result has discriminant valuation six,
coefficient valuations $(2,\ge0)$, cyclic-four residual image, and
faithful special inertia of order six. (The unique nontrivial tame
quadratic character multiplies its central tame action by $-1$.)
The preceding argument applies. Twisting back proves the assertion.

For completeness the good fields can be exhibited directly.
The Newton polygon of $r^3-3r+3b$ gives $v_K(r)=2/3$, so its
root field has degree three. Its discriminant is
$27(4-9b^2)$, a square in $K$; it is therefore cyclic. All three
pairwise root differences have valuation one, since Galois acts
transitively on the three unordered pairs and the discriminant has
valuation six. The $\pi$-twist of (B-RD42.1) has equation
$y^2=x^3-9x+\pi^5b$. The substitution
$x=\pi^2X+\pi r$, $y=\pi^3Y$ gives
\[
              Y^2=X^3+\pi rX^2+(r^2-1)X.
                                                        \tag{B-RD42.4}
\]
Its reduction is $Y^2=X^3-X$, a good supersingular curve.
A nontrivial $\sigma\in\operatorname{Gal}(L_b/K)$ changes $X$ by
$(r-\sigma r)/\pi$, whose residue is nonzero, and fixes $Y$.
Thus its action on the special curve is nontrivial of order three.
Untwisting requires $\sqrt\pi$; the added quadratic action is
$[-1]$, proving both faithful inertia assertions. In particular
these displayed good fields are the minimal good fields.

One can also check the ramification directly: $r^2/\pi$ is a
uniformizer of $L_b$, and
\[
 v_{L_b}\bigl(\sigma(r^2/\pi)-r^2/\pi\bigr)=2.
\]
The cubic lower break is one. This observation is not needed for
the equality criterion below.

Choose the square root $v_b=\sqrt{1-9b^2/4}$ congruent to one, and put
\[
                 D_b=v_b+\frac{3i}{2}b.
                                                        \tag{B-RD42.5}
\]
If $u^3=iD_b$, then $u^3+u^{-3}=-3b$, so $u+u^{-1}$ is a root
of the cubic. A cube root of $i$ lies in $K$, and
\[
                    D_b\equiv1+\frac{3i}{2}b\pmod{\pi^3}.
                                                        \tag{B-RD42.6}
\]
This element is not a cube: its deviation from one has valuation
two, whereas a cube congruent to one has deviation at least three.
Consequently $K(\sqrt[3]{D_b})=L_b$. By
Lemma~\ref{lem:B-RD39-units}, every element of $1+\pi^3\OO_K$ is a
cube. The Kummer classes of $D_b$ and $D_{b'}$ therefore generate
the same line in $K^\times/(K^\times)^3$ exactly when
$\overline{b'}=\pm\overline b$. This proves (B-RD42.2).
\end{proof}

\begin{lemma}[The reverse marked cocycle identifies the cubic good fields]
\label{lem:B-RD42-reverse}
\status{Proved. Complete root and independent original-source, proof and exact-certificate audits passed.}
Suppose that $M_b[9]\simeq M_{b'}^{(\pi)}[9]$ is antisymplectic,
and that their common residual representation is irreducible cyclic
of order four. Then $L_b=L_{b'}$.
\end{lemma}
\begin{proof}
Let $L=K(M_b[3])$ and let $F$ be its quadratic subfield. The residual
isomorphism gives the same fields on the other side. Choose symplectic
three-torsion markings with a fixed Weil-pairing value $\zeta$. Since
the given isomorphism reverses the pairing, its matrix can be taken
to be $R=\operatorname{diag}(1,-1)$ between these markings. Choose
the cyclic-four generator as
$S=\left(\begin{smallmatrix}0&-1\\1&0\end{smallmatrix}\right)$
on the first side. On the second side its matrix is
$RSR^{-1}=S^{-1}=-S$. This is precisely Fisher's reverse-marking
involution $\iota=\operatorname{diag}(1,-1)$; his inverse-transpose
twist dictionary is \citet[Section~3, p.~7, and Lemma~3.2, pp.~7--8]
{Fisher2015Nine}.

The projective action on the two Hesse parameters is the same
involution $h(t)=(t+2)/(t-1)$. Write $z=t-1$, $z'=t'-1$ and
$u=z/\pi$, $u'=z'/\pi$. As in Proposition~\ref{prop:B-RD40},
\[
 N_{F/K}(z)=N_{F/K}(z')=3,\qquad
 N_{F/K}(u)=N_{F/K}(u')=1,\qquad
 \overline u,\overline{u'}\in\{1,-1\}.
                                                        \tag{B-RD42.7}
\]
The actual Fourier lift for $S$ has differential multiplier
$\kappa(t)=c/(t-1)$, where $c=\zeta-\zeta^2$. The lift for $S^{-1}$
is its central negative: its multiplier is $-\kappa(t)$.
Indeed $S^{-1}=-S$, the central element is elliptic inversion,
and the exact normalization and identity
$\kappa(t)\kappa(h(t))=-1$ were proved in
Proposition~\ref{prop:B-RD40}. Thus no change of plane-residue
normalization or parameter-dependent factor is being suppressed.

Let $\alpha,\alpha'$ compare the respective universal Hesse
invariant differentials with invariant differentials over $K$.
The analogue of (B-RD40.7), now with opposite actual lifts, is
\[
       \frac{\tau(\alpha'/\alpha)}{\alpha'/\alpha}
                 =-\frac{z'}z\quad\hbox{or its inverse},
                                                        \tag{B-RD42.8}
\]
for a generator $\tau$ of $\operatorname{Gal}(L/K)$.
The discriminant valuations are six and twelve, so the invariant
section $\Delta\omega^{12}$ gives
$v_L(\alpha')-v_L(\alpha)=2$. Tame leading coefficients therefore
make the residue of the left side equal to $-1$. It follows that
\[
                            \overline{u'}=\overline u.
                                                        \tag{B-RD42.9}
\]

Over $F$ twist both curves by the same quadratic character of $L/F$.
Their three-torsion becomes rational, the Hesse curves are $E_t,E_{t'}$,
and the reverse nine-congruence persists. Fisher's reverse equation
(11), with its essential numerator three, gives at $\theta=1$
\[
                          u'u=\frac{(t'-1)(t-1)}3
                                    \in(F^\times)^3.
                                                        \tag{B-RD42.10}
\]
Here we use exactly the full-level marking convention of
\citet[Section~2, pp.~4--5; Theorem~4.3 and its proof, pp.~9--10]
{Fisher2015Nine}. By (B-RD42.9), this cube is congruent to one.
Since $v_F(3)=4$, expansion of $(1+w)^3$ shows
$v_F(u'u-1)\ge3$. Set
\[
 T=z+3/z=\pi(u+u^{-1}),\qquad
 T'=z'+3/z'=\pi(u'+u'^{-1})\in K.
\]
The exact identity
\[
       T'-T=\pi(u'-u^{-1})\left(1-\frac u{u'}\right)
                                                        \tag{B-RD42.11}
\]
and (B-RD42.9) give $v_F(T'-T)\ge2+3+1=6$, hence
\[
                              v_K(T'-T)\ge3.
                                                        \tag{B-RD42.12}
\]

We recover the residues of the normal-form parameters from these
traces. Eliminating $z^2-Tz+3=0$ from the Hesse $j$-formula gives
\[
 j(E_t)=27\left(\frac{T(T+4)}{T+3}\right)^3,
 \qquad
 j(M_b)=\frac{6912}{4-9b^2}.
                                                        \tag{B-RD42.13}
\]
These are exact rational identities, and twisting does not change
$j$. Put $H=T(T+4)/(T+3)$. If $\overline u=\epsilon\in\{1,-1\}$,
then the identity
\[
 u+u^{-1}-2\epsilon=\frac{(u-\epsilon)^2}{u}
\]
and $v_F(\pi)=2$ give $v_F(T-2\epsilon\pi)\ge4$. Since
$T\in K$, we may write $T=2\epsilon\pi+\pi^2a$ with
$a\in\OO_K$. Also
\[
 H-4=\frac{T^2-12}{T+3}\in\pi^2\OO_K,\qquad
 H^3=\frac{64}{1-9b^2/4}.
\]
Writing $H=4+\pi^2d$ and comparing the latter identity modulo
$\pi^5$ gives $d\equiv b^2\pmod\pi$. The former identity gives
$d\equiv2a\pmod\pi$. Consequently
\[
                T\equiv2\epsilon\pi+\frac{\pi^2}{2}b^2
                                      \pmod{\pi^3}.
                                                        \tag{B-RD42.14}
\]
The same formula holds for $T'$ and $b'$ with the same $\epsilon$
by (B-RD42.9). Equation (B-RD42.12) therefore gives
$\overline{b'}^{\,2}=\overline b^{\,2}$. In characteristic three
this says $\overline{b'}=\pm\overline b$, and
Lemma~\ref{lem:B-RD42-normal} proves $L_b=L_{b'}$.
\end{proof}

\begin{theorem}[B-RD42: the exact least height-two index at $p=3,e=2$]
\label{thm:B-RD42}
\status{Proved. Complete root and independent original-source, proof and exact-certificate audits passed.}
For the full class of height-two groups in B.4 satisfying the manuscript
monodromy hypothesis,
\[
                              \boxed{\nu_{\min}(3,2,2)=2}.
                                                        \tag{B-RD42.15}
\]
Equivalently, $27$-torsion determines the requested character on
prime-to-three inertia elements, while nine-torsion does not always do so.
This is a sharp small-height case; no general least-index formula is
asserted for other heights or ramification indices.
\end{theorem}
\begin{proof}
The lower bound is the explicit nine-torsion example of
Proposition~\ref{prop:B-RD29-p3-level-two}. The reduction and bounds in
Theorem~\ref{thm:B-RD34} and Proposition~\ref{prop:B-RD35} reduce
any putative failure at $27$ to two potentially supersingular elliptic
curves after completed maximal unramified base change. These steps
already cover arbitrary perfect residue fields and all the allowed
height-two monodromy filtrations. The completed maximal unramified
base change preserves inertia and absolute ramification; its algebraically
closed residue field makes every unit a square, so $\pi^2=3$ and
$\mu_3\subset K$ can be arranged without a further ramified extension.

Propositions~\ref{prop:B-RD37} and \ref{prop:B-RD41} exclude the
scalar and unipotent residual cases. By
Propositions~\ref{prop:B-RD38} and \ref{prop:B-RD40}, the sole remaining
case is an antisymplectic comparison with cyclic-four residual image,
individual faithful special-inertia orders three and six, and joint
faithful inertia order eighteen. Lemma~\ref{lem:B-RD42-normal}
writes this pair, after interchanging it, as $M_b,M_{b'}^{(\pi)}$.
Restriction of the normalized antisymplectic $27$-isomorphism to
nine-torsion satisfies Lemma~\ref{lem:B-RD42-reverse}; hence
$L_b=L_{b'}$. Both curves then have good reduction over the same
extension $L_b(\sqrt\pi)$ of degree six. Their joint faithful special
inertia has order at most six, contradicting the required order
eighteen. This excludes the last case and proves (B-RD42.15).
\end{proof}

\par\smallskip\noindent\textbf{Exact certificate and scope.}
The program \path{code/verify_c4_final_exclusion.py} checks the universal
normal-form and good-model invariants, the Cardano and Kummer identities,
the trace elimination and truncated expansions, the opposite marked
matrices, and the integer valuation cases used in the proof. Its output
is retained separately, with author runtime $0.1415409550027107$ seconds.
The initial structural-equality check and traceback are retained in
\path{failed_candidates/B42_certificate_structural_equality_v1.py} and
its adjacent output; the corrected certificate tests exact polynomial
differences. The preceding proof uses the algebraic-closure
residue hypothesis only in the established reduction and the explicitly
stated Hensel and unit-cube steps. The finite-elimination procedure of
B-RD24 remains unevaluated; it is not used to obtain this exact index.

\par\smallskip\noindent\textbf{Independent audits.}
The complete proof and exact certificate were independently checked by
two reviewers. Their author-output-preserving reruns passed in
$0.11300166400178568$ and $0.10002119700220646$ seconds, respectively.
The reverse-marking convention was checked against Fisher's original
Section~3 and the full proof of Theorem~4.3. This establishes only the
explicit small instance (B-RD42.15); the general least-index problem
in B.4 remains open.


\subsection{An explicit eight-torsion ambiguity in the unramified dyadic case}

\begin{proposition}[B-RD43: even eight-torsion does not suffice at $p=2,e=1$]
\label{prop:B-RD43}
\status{Proved. Complete root and independent original-source, proof and exact-certificate audits passed.}
Let $K=\mathbf Q_2(\omega)$, where $\omega^2+\omega+1=0$, and put
$\lambda=1+2\omega$, so $\lambda^2=-3$. Define
\[
 E:\ y^2+(5-\lambda)xy-4\lambda y=x^3-2\lambda x^2,
 \qquad P=(0,0),\qquad B=E/\langle P\rangle,
 \qquad q=-\frac{4\lambda}{3}.
                                                        \tag{B-RD43.1}
\]
Then $P$ has exact order eight, $v_K(q)=2>0$, and the split Tate
curve $T_q$ satisfies
\[
                             B[8]\simeq T_q[8].
                                                        \tag{B-RD43.2}
\]
The requested characters of $B[2^\infty]$ and $T_q[2^\infty]$
differ on an inertia element whose image has order three in the
finite special-inertia quotient. Consequently
\[
                       \nu_{\min}(2,1,h)\ge3\qquad(h\ge2).
                                                        \tag{B-RD43.3}
\]
The same lower bound holds for every absolute ramification index
$e$ not divisible by three. No assertion of sufficiency of
sixteen-torsion is made.
\end{proposition}
\begin{proof}
The field $K$ is unramified quadratic over $\mathbf Q_2$, since
$X^2+X+1$ is separable and irreducible over $\mathbf F_2$.
The displayed equation is integral. Exact addition gives
\[
\begin{array}{c|c|c}
 n&x(nP)&y(nP)\\\hline
 1&0&0\\
 2&2\lambda&-6-6\lambda\\
 3&-6+2\lambda&12\\
 4&(-3+\lambda)/2&3\\
 5&-6+2\lambda&12-12\lambda\\
 6&2\lambda&0\\
 7&0&4\lambda
\end{array}
\qquad 8P=0.
                                                        \tag{B-RD43.4}
\]
In particular $P$ has exact order eight. This curve arises from
the $X_1(8)$ Tate normal form in
\citet[Appendix~A.2, pp.~26--27, Figure~8]{AntieauAuel2021Descent}
by taking its parameter to be $\lambda$ and scaling
$x=X/4$, $y=Y/8$. The displayed addition certificate verifies
its torsion without relying on that parametrization.

We compute the exact extension class in
Lemma~\ref{lem:B-RD29-isogeny-Tate}, so that the Tate parameter
in (B-RD43.1) is specified rather than merely asserted to exist.
The tangent at $P$ is $y=0$. The tangent at $2P$ is
\[
 \ell_2=y-\frac{-9+\lambda}{2}x-(-3+3\lambda)=0.
\]
Put $x_2=2\lambda$ and $x_4=(-3+\lambda)/2$.
The Miller doubling identity gives a function with divisor
$8[P]-8[0]$:
\[
 f=\frac{y^4\ell_2^2}{(x-x_2)^4(x-x_4)}.
\]
Here the last tangent is the vertical line at the two-torsion point
$4P$. Reducing by the equation of $E$ cancels the denominator and gives
\[
\begin{split}
 f(x,y)={}&x^4-4\lambda x^3-8x^2
            +(2\lambda-6)x^2y\\
          &+(12\lambda+4)xy+16y.
\end{split}                                             \tag{B-RD43.5}
\]
This polynomial identity and all the displayed point additions
are verified exactly in $\mathbf Q(\lambda)$.
The divisor description of the dual-isogeny connecting homomorphism
is \citet[Appendix~A.1, Lemmas~A.1--A.2, p.~22]
{AntieauAuel2021Descent}. Its unspecified normalizing constant
cancels in the ratio at $3P$ and $2P$, since both avoid the zero
and pole of $f$. Thus
\[
 \delta(P)=\frac{f(3P)}{f(2P)}
           =\frac{576\lambda}{-432}
           =-\frac{4\lambda}{3}
                   \quad\hbox{in }K^\times/(K^\times)^8.
                                                        \tag{B-RD43.6}
\]
The quotient $B=E/\langle P\rangle$ therefore fits into the
Kummer extension
$0\to\mu_8\to B[8]\to\mathbf Z/8\to0$ with the indicated
class. Since $v_K(q)=2$, Lemma~\ref{lem:B-RD29-isogeny-Tate}
and Tate's explicit finite-order construction identify it with
$T_q[8]$. As an additional check, the Figure~8 representative
$\lambda^7(\lambda-1)^6(\lambda+1)^4$ divided by (B-RD43.6)
is exactly $(-3-\lambda)^8$.

The invariant calculation gives
\[
 c_4=-200-312\lambda,\qquad
 c_6=15488+576\lambda,\qquad
 \Delta=20736(\lambda-2),
\]
so
\[
                     v_K(c_4)=4,\quad v_K(\Delta)=8,
                     \quad v_K(j)=4.
                                                        \tag{B-RD43.7}
\]
For example, $c_4=-16(32+39\omega)$ and
$\lambda-2=-1+2\omega$ are respectively sixteen times a unit
and a unit. The displayed discriminant is minimal, since eight
is less than twelve and integral discriminant valuations are
nonnegative. Integral $j$ gives potential good reduction by
\citet[Section~2, p.~496]{SerreTate1968Good}; its positive valuation
gives supersingular special reduction in characteristic two,
as in the already verified proof of Proposition~\ref{prop:B-RD22}.

Over the completed maximal unramified field, let $m$ be a minimal
good-reduction degree. The change to a unit-discriminant equation
implies $12\mid8m$, hence $3\mid m$. Faithful special inertia
therefore contains an element of order three. Its trace on the
rank-two rational Dieudonne module is $-1$, as in
Proposition~\ref{prop:B-RD22}. The isogeny $E\to B$ preserves this
rational module and its finite inertia action. The Tate group's
two good rank-one graded pieces have trace $2$ on that element.
The order-three assertion concerns the finite special-inertia image;
no finite-order lift in the absolute inertia group is asserted.
Both groups satisfy the manuscript monodromy condition, and the
characters differ despite (B-RD43.2).

Finally, take an Eisenstein extension of $K$ of degree $e$.
The torsion isomorphism persists and the raw discriminant valuation
becomes $8e$. If $3\nmid e$, the relation $12\mid8em$ still
forces $3\mid m$, so the same character discrepancy survives.
Adding $h-2$ identical constant etale factors proves every height
$h\ge2$. Since eight-torsion is $M(2)$, its failure proves
(B-RD43.3).
\end{proof}

\par\smallskip\noindent\textbf{Exact certificate and scope.}
The program \path{code/verify_p2_character_level3.py} checks all eight
subgroup points and all $64$ additions, the exact invariant
polynomials, both tangent lines, the affine Miller-function identity,
its two evaluations and the eighth-power parameter comparison.
Its retained author output records $0.5159330129972659$ seconds.
The finite-scheme extension dictionary and the special-inertia
argument are proved above using the stated primary sources.
This counterexample raises the lower bound; it does not compute
the least dyadic index.

Independent reruns took $0.5299817019986222$ and
$0.6142895730008604$ seconds, preserving the author output. Both reviewers
checked the complete proof and the original descent and Tate sources.

